\chapter{Abelian-variety rigidity: morphisms from a genus-\(0\) curve into an abelian variety are constant}
\label{chap:AbelianVarietyRigidity}
\label{chap:avr_for_rr}
% archon:covers AlgebraicJacobian/AbelianVarietyRigidity.lean AlgebraicJacobian/Genus0BaseObjects.lean AlgebraicJacobian/Genus0BaseObjects/BareScheme.lean AlgebraicJacobian/Genus0BaseObjects/ChartIso.lean AlgebraicJacobian/Genus0BaseObjects/Points.lean AlgebraicJacobian/Genus0BaseObjects/GmScaling.lean AlgebraicJacobian/RigidityLemma.lean
% Local macro: diagrammatic (Lean `≫`) composition. Defined locally because this
% chapter is the writer's only editable file; flagged in the writer report for promotion
% to macros/common.tex if desired. \providecommand is a no-op if it is later added globally.
\providecommand{\fatsemi}{\mathbin{\,;\,}}
% STRATEGY NOTE (iter-157). This chapter is the dedicated home of the project's COMMITTED
% genus-0 route (route (c)): the characteristic-free abelian-variety rigidity stack. It is
% 1:1 with the new upstream Lean file AlgebraicJacobian/AbelianVarietyRigidity.lean, which
% imports AlgebraicJacobian.Genus and is imported by AlgebraicJacobian.Jacobian (breaking the
% RigidityKbar -> Rigidity -> Jacobian import cycle so that genusZeroWitness can consume the
% keystone directly). The headline prop:rigidity_genus0_curve_to_AV is the char-free
% replacement (no [CharZero]) for the over-k-bar declaration rigidity_over_kbar of
% \cref{chap:RigidityKbar} (which remains in tree as the fallback route (a) artifact and
% still carries [CharZero]). The df=0 / H^0(P^1,Omega)=0 / Serre-duality framing is the
% fallback (a), NOT the route developed here.
The genus-\(0\) arm of the Jacobian-witness construction has a \emph{trivial} witness object
(\(J = \Spec k\)); its only remaining mathematical content is a rigidity \emph{statement}:
every pointed morphism from a smooth proper geometrically irreducible curve of genus \(0\)
into an abelian variety is the constant morphism at the identity. We prove this
\textbf{characteristic-free} --- the protected goal is over an arbitrary field \(k\), so the
cheaper differential argument (via \(H^0(\mathbb P^1, \Omega) = H^0(\mathbb P^1, \mathcal
O(-2)) = 0\), hence \(df = 0\), hence \(f\) constant) is unavailable: it breaks on a
characteristic-\(p\) Frobenius-descent wall, where a nonconstant \(f\) can have \(df = 0\). The
rigidity route below sidesteps that wall.
The chain has two parts. The \emph{Rigidity Lemma} (\cref{thm:rigidity_lemma}, Mumford,
Form~I; equivalently Milne's Rigidity Theorem~1.1) is an elementary properness/closed-map
argument --- it uses neither the theorem of the cube nor any cohomology, is the entry point
for formalisation, and is now proven axiom-clean. From the Rigidity Lemma, Milne's \S I.3
derives the genus-\(0\) base case \emph{without} the theorem of the cube --- and, in this
project's reading, \emph{without} Milne's Theorem~3.2 (rational-map extension), Lemma~3.3,
Auslander--Buchsbaum, or differentials (the ``4th route''; see
\cref{rmk:base_case_fourth_route}). The decisive observation is that \(\mathbb P^1\) is
\emph{already complete} and \(f \colon \mathbb P^1 \to A\) is \emph{already a total morphism}, so
there is no rational defect map to extend over a complete surface and Theorem~3.2 never enters.
Instead the base case is closed by the \textbf{\(\mathbb G_m\)-scaling shortcut} (the project's
primary route, an Archon-original realisation of Milne Proposition~3.9): one uses the
\emph{multiplicative scaling} action \(\sigma_\times \colon \mathbb P^1 \times \mathbb G_m \to
\mathbb P^1\), \((x, \lambda) \mapsto \lambda x\) (the M\"obius scaling fixing both \(0\) and
\(\infty\)), which is a \emph{genuine total} scheme morphism on all of \(\mathbb P^1 \times \mathbb
G_m\) (\cref{def:gaTranslationP1}). Normalising \(f(0) = \eta_A\) and composing \(h := \sigma_\times
\fatsemi f\), one applies additivity of \(\mathrm{Hom}(-, A)\) over a product
(\cref{lem:hom_additivity_over_product}, Milne Corollary~1.5) --- whose Lean signature requires
only the \emph{first} factor complete --- with \(V = \mathbb P^1\) proper, \(W = \mathbb G_m\), and
base points \(0 \in \mathbb P^1\), \(1 \in \mathbb G_m\). Because \(0\) is a \emph{fixed point} of
scaling, the \(W\)-axis restriction \(\lambda \mapsto h(0, \lambda) = f(\lambda \cdot 0) = f(0)\) is
\emph{constant} at \(\eta_A\), so the decomposition collapses to \(h = p \fatsemi f\), i.e.\
\(f(\lambda x) = f(x)\). Specialising \(x = 1\) gives that \(f|_{\mathbb G_m}\) is constant, and since
\(\mathbb G_m\) is dense in \(\mathbb P^1\) with \(A\) separated, \(f\) is constant
(\cref{prop:morphism_P1_to_AV_constant}, via the separated-target/dominant-source handle
\cref{thm:GrpObj_eq_of_eqOnOpen}). This route uses \emph{only} the proven Corollary~1.5, a
scaling fixed point, and density --- it never forms the additive defect map, never invokes
\(\mathrm{Hom}(\mathbb G_a, A) = 0\), and so sidesteps the ``image of a group homomorphism is a
closed subgroup'' input that has no Mathlib support. The classical \(\mathbb G_a\)-additive
argument (\cref{lem:hom_from_Ga_trivial} via the \(\mathbb G_a/\mathbb G_m\) incompatibility
\cref{lem:hom_Ga_to_av_trivial}, the heart of Milne Proposition~3.9) is retained only as the
Milne-faithful alternative and is \emph{off} the genus-\(0\) critical path. The corollary that a
pointed regular map of abelian
varieties is a homomorphism (\cref{lem:av_regular_map_is_hom}, Milne Corollary~1.2) is recorded
for completeness and for Route~A but is \emph{not} needed on the genus-\(0\) path; Milne's
Theorem~3.2 (\cref{thm:rational_map_to_av_extends}) is likewise demoted to a Route-A-only item
(the Albanese universal property), no longer on the genus-\(0\) critical path
(\cref{rmk:thm32_codim1_mathlib_gap}). Combining these
gives that every morphism \(\mathbb P^1 \to A\) is constant
(\cref{prop:morphism_P1_to_AV_constant}); composing with the classification of genus-\(0\)
curves (\cref{thm:genus_zero_curve_iso_p1}, \(C \cong \mathbb P^1\) over \(\bar k\)) yields the
headline \cref{prop:rigidity_genus0_curve_to_AV}. The theorem of the cube (Milne \S I.5) is
\emph{not} used anywhere on this path (see \cref{rmk:cube_not_needed}); the chain invokes
neither Serre duality nor the FGA \(\Pic\) representability engine, so the genus-\(0\) arm is
fully decoupled from Route~A.
\medskip
Throughout, \(\bar k\) is an algebraically closed field of arbitrary characteristic. An
\emph{abelian variety} over \(\bar k\) is a complete (proper) irreducible group variety; in the
project's Lean encoding \(A\) is an object of \(\mathrm{Over}\,(\Spec \bar k)\) carrying
\(\mathtt{GrpObj}\), \(\mathtt{IsProper}\), \(\mathtt{Smooth}\), and \(\mathtt{GeometricallyIrreducible}\).
\section{The Rigidity Lemma}
\begin{lemma}
\leanok
[Rigidity Lemma (Mumford, Form I)]
  \label{thm:rigidity_lemma}
  \uses{lem:rigidity_eqOn_dense_open, lem:rigidity_core, lem:rigidity_snd_lift}
  
  \lean{AlgebraicGeometry.rigidity_lemma}
  % SOURCE: [Mumford, Abelian Varieties], Rigidity Lemma (Form I.), Ch.~II \S4, book p.~43 (read from references/mumford-abelian-varieties.pdf, PDF page 54)
  % SOURCE QUOTE: "RIGIDITY LEMMA. (Form I.) Let \(X\) be a complete variety, \(Y\) and \(Z\)
  % any varieties, and \(f \colon X \times Y \to Z\) a morphism such that for some
  % \(y_0 \in Y\), \(f(X \times \{y_0\})\) is a single point \(z_0\) of \(Z\). Then there is a
  % morphism \(g \colon Y \to Z\) such that if \(p_2 \colon X \times Y \to Y\) is the
  % projection, \(f = g \circ p_2\)."
  \textit{Source: Mumford, Abelian Varieties, Ch.~II \S4, Rigidity Lemma (Form~I), p.~43.}
  Let \(X\) be a complete variety, and let \(Y\), \(Z\) be any varieties over \(\bar k\). Let
  \(f \colon X \times Y \to Z\) be a morphism such that for some \(y_0 \in Y\) the image
  \(f(X \times \{y_0\})\) is a single point \(z_0 \in Z\). Then there is a morphism
  \(g \colon Y \to Z\) with \(f = g \circ p_2\), where \(p_2 \colon X \times Y \to Y\) is the
  projection; that is, \(f\) depends only on the \(Y\)-coordinate, and \(g(y) = f(x_0, y)\) for any
  fixed \(x_0 \in X\).
\end{lemma}
% SOURCE QUOTE PROOF: "PROOF. Choose any point \(x_0 \in X\), and define \(g \colon Y \to Z\) by
% \(g(y) = f(x_0, y)\). Since \(X \times Y\) is a variety, to show that \(f = g \circ p_2\), it is
% sufficient to show that these morphisms coincide on some open subset of \(X \times Y\). Let
% \(U\) be an affine open neighbourhood of \(z_0\) in \(Z\), \(F = Z - U\), and \(G = p_2(f^{-1}(F))\);
% then \(G\) is closed in \(Y\) since \(X\) is complete and hence \(p_2\) is a closed map. Further
% \(y_0 \notin G\) since \(f(X \times \{y_0\}) = \{z_0\}\). Therefore \(Y - G = V\) is a non-empty
% open subset of \(Y\). For each \(y \in V\), the complete variety \(X \times \{y\}\) gets mapped by
% \(f\) into the affine variety \(U\), and hence to a single point of \(U\). But this means that for
% any \(x \in X\), \(y \in V\), \(f(x, y) = f(x_0, y) = g \circ p_2(x, y)\), and this proves our
% assertion."
\begin{proof}
\leanok
% NOTE (iter-159 review): the iter-157 unsoundness REMAINS REPAIRED. The collapse hypothesis
% (_hf : f(X x {y0}) = {z0}) is threaded through all three helpers (rigidity_lemma ->
% rigidity_core -> rigidity_eqOn_dense_open) and GENUINELY CONSUMED in the body of
% rigidity_eqOn_dense_open (proves y0 \notin G, the load-bearing non-emptiness of V := Y - G).
% The counterexample f = fst is excluded (under f = fst, _hf forces X to collapse to a point).
% rigidity_eqOn_dense_open is sorry-free in its OWN body (hfib, the pullback-fibre fact over the
% k-bar-point y0, is closed). rigidity_snd_lift and snd_left_isClosedMap are axiom-clean (no sorryAx).
% NOTE (iter-160): the \uses edges run FORWARD along the true Lean dependency order
% (rigidity_lemma -> rigidity_eqOn_dense_open -> rigidity_eqOn_saturated_open_to_affine), so the
% not-proven status of the residual sorry propagates up and de-launders BOTH thm:rigidity_lemma and
% lem:rigidity_eqOn_dense_open. No backward edge from the leaf lemmas back up the chain.
% NOTE (iter-161): bridge 2 (rigidity_eqOn_saturated_open_to_affine) was DECOMPOSED into two named
% top-level sub-lemmas along the route-B (cohomology-free) split:
%   - Step 2, morphism_eq_of_eqAt_closedPoints: "two morphisms agreeing at every closed point of a
%     reduced Jacobson source into a separated target are equal" -- PROVEN axiom-clean (the
%     dense-closed-points coproduct probe + ext_of_isDominant; the connective Mathlib lacks).
%   - Step 1, rigidity_eqAt_closedPoint_of_proper_into_affine: Mumford's per-closed-slice constancy
%     -- the genuinely-deep residual, still `sorry`. THIS is now the chain's lone deep residual
%     (bridge 2's body is real assembly: Step 2 proven, the JacobsonSpace U instance now derivable).
% SIGNATURE CHANGE (iter-161, plan-authorized): Step 2's globalisation needs the closed points of
% the saturated open U to be DENSE, i.e. U Jacobson, which over Spec k-bar follows from (X tensor Y)
% being LOCALLY OF FINITE TYPE. So [LocallyOfFiniteType (X tensor Y).hom] is now carried on FIVE
% chain lemmas (rigidity_lemma, rigidity_core, rigidity_eqOn_dense_open,
% rigidity_eqOn_saturated_open_to_affine, rigidity_eqAt_closedPoint_of_proper_into_affine), in
% addition to [IsAlgClosed kbar]. Free downstream: curves and abelian varieties are of finite type
% over k-bar. The earlier "[IsAlgClosed] is the only added instance" claim is hereby retired.
  Fix any point \(x_0 \in X\) and define \(g \colon Y \to Z\) by \(g(y) = f(x_0, y)\) (the
  composite of \(y \mapsto (x_0, y)\) with \(f\)). Since \(X \times Y\) is a variety (in
  particular reduced and irreducible, so separated and with no embedded structure), to prove
  the equality of morphisms \(f = g \circ p_2\) it suffices to prove they agree on a non-empty
  open subset of \(X \times Y\).
  Let \(U\) be an affine open neighbourhood of \(z_0\) in \(Z\), put \(F = Z - U\) (a closed set),
  and set \(G = p_2(f^{-1}(F))\). Because \(X\) is complete, the projection
  \(p_2 \colon X \times Y \to Y\) is a \emph{closed} map (this is the defining property of
  completeness/properness); hence \(G\), the image of the closed set \(f^{-1}(F)\), is closed in
  \(Y\). The hypothesis \(f(X \times \{y_0\}) = \{z_0\} \subseteq U\) gives
  \(f^{-1}(F) \cap (X \times \{y_0\}) = \varnothing\), so \(y_0 \notin G\). Therefore
  \(V := Y - G\) is a non-empty open subset of \(Y\).
  For each \(y \in V\) we have \(X \times \{y\} \cap f^{-1}(F) = \varnothing\), i.e.
  \(f(X \times \{y\}) \subseteq U\). Thus the complete (proper) variety \(X \times \{y\}\) is
  mapped by \(f\) into the \emph{affine} variety \(U\). A proper connected variety mapping to an
  affine variety has image a single point (a global regular function on a complete connected
  variety is constant). Hence \(f(X \times \{y\})\) is one point, necessarily
  \(f(x_0, y) = g(y)\). So for all \(x \in X\) and \(y \in V\),
  \[ f(x, y) = f(x_0, y) = g(y) = (g \circ p_2)(x, y). \]
  The two morphisms \(f\) and \(g \circ p_2\) thus agree on the non-empty open
  \(X \times V\), which is dense in the irreducible \(X \times Y\); as \(Z\) is separated they
  agree everywhere. \qedhere
\end{proof}
\begin{remark}
  \label{rmk:rigidity_lemma_cube_free}
  The proof of \cref{thm:rigidity_lemma} uses \emph{only} (i) that completeness of \(X\) makes
  \(p_2\) a closed map, and (ii) that a proper connected variety has no nonconstant map to an
  affine variety. It uses neither the theorem of the cube nor any cohomology, and is valid in
  arbitrary characteristic. This is the designated formalisation entry point: it is the
  lowest-prerequisite link of the chain. (Mumford gives this as the geometric heart of the
  ``second proof'' of commutativity of an abelian variety; Milne packages the same statement
  as the Rigidity Theorem~1.1, requiring \(V\) complete and \(V \times W\) geometrically
  irreducible.)
\end{remark}
\begin{remark}
  \label{rmk:rigidity_lemma_decomposition}
  The Lean formalisation of \cref{thm:rigidity_lemma} splits into three helper steps, mirroring
  the proof above: \texttt{rigidity\_snd\_lift} (the cartesian-monoidal algebra identifying the
  collapsed map \((x, y) \mapsto (x_0, y)\) followed by \(f\) with \(g \circ p_2\));
  \texttt{rigidity\_core} (the scheme-level gluing: two maps out of the geometrically irreducible
  reduced \(X \times Y\) into the separated \(Z\) that agree on a non-empty open subset agree
  everywhere, by separatedness); and \texttt{rigidity\_eqOn\_dense\_open}
  (\cref{lem:rigidity_eqOn_dense_open}, the genuine geometry: the construction of the non-empty
  open \(X \times V\) and the slice-constancy of \(f\) on it). The collapse hypothesis
  \(f(X \times \{y_0\}) = \{z_0\}\) must be threaded through \emph{all three} helpers, since it is
  what makes \(V\) non-empty; a decomposition that drops it from the lower helpers is unsound
  (see \cref{lem:rigidity_eqOn_dense_open}).
  \emph{Status (iter-161).} Bridge~1 (the closed-map step) is \textbf{built} as the axiom-clean
  helper \texttt{snd\_left\_isClosedMap}. The fibre-over-\(y_0\) fact is \textbf{closed} (via the
  coarse \texttt{PullbackCarrier} layer), so \texttt{rigidity\_eqOn\_dense\_open} is
  \texttt{sorry}-free in its own body. Bridge~2 (slice-constancy on the saturated open \(U\),
  \cref{lem:rigidity_eqOn_saturated_open_to_affine}) has now been \textbf{decomposed} along its
  cohomology-free route-(c) split into two named sub-lemmas: the globalisation
  \cref{lem:morphism_eq_of_eqAt_closedPoints} (Step~2, ``agree at every closed point of a reduced
  Jacobson source \(\Rightarrow\) equal'') is \textbf{proven axiom-clean}, and the per-closed-slice
  constancy \cref{lem:rigidity_eqAt_closedPoint_of_proper_into_affine} (Step~1, Mumford's slice
  argument) is the chain's \emph{single genuinely-deep residual} \texttt{sorry}. The body of
  \cref{lem:rigidity_eqOn_saturated_open_to_affine} is therefore real assembly (Step~2 over Step~1).
  % NOTE: (iter-162 review) STALE as of iter-162 — Step~1 is now PROVEN and the WHOLE
  % Rigidity-Lemma chain (thm:rigidity_lemma … lem:rigidity_eqAt_closedPoint_of_proper_into_affine,
  % via the new helper lem:isIntegral_of_retract_of_integral) is sorry-free + axiom-clean
  % ({propext, Classical.choice, Quot.sound}, re-verified). The "single genuinely-deep residual
  % sorry" wording here and in the proof blocks of lem:rigidity_eqOn_dense_open / …saturated… /
  % …eqAt_closedPoint… should be refreshed by the plan/blueprint-writer to "chain closed iter-162".
  This route requires the base field to be algebraically closed \emph{and} the source to be of finite
  type: Step~2 needs the closed points of \(U\) to be dense (the Jacobson-space property), which over
  the algebraically closed \(\bar k\) holds because \(X \times Y\) is locally of finite type.
  Accordingly the formalization carries \textbf{both} \([\mathtt{IsAlgClosed}\ \bar k]\) and
  \([\mathtt{LocallyOfFiniteType}\ (X \otimes Y).\mathrm{hom}]\) across the chain lemmas
  (\texttt{rigidity\_lemma}, \texttt{rigidity\_core}, \texttt{rigidity\_eqOn\_dense\_open},
  \texttt{rigidity\_eqOn\_saturated\_open\_to\_affine},
  \texttt{rigidity\_eqAt\_closedPoint\_of\_proper\_into\_affine}); the earlier claim that
  \([\mathtt{IsAlgClosed}\ \bar k]\) is the \emph{only} added instance is retired. Both cost nothing
  downstream: curves and abelian varieties are of finite type over the algebraically closed \(\bar k\),
  and the consumers already assume it. The relative Stein-factorisation / proper-pushforward
  \(f_*\mathcal O = \mathcal O\) framing is a genuine Mathlib cohomology gap and is \emph{deliberately
  avoided}.
\end{remark}
\begin{lemma}\leanok
[Cartesian-monoidal collapse identity]
  \label{lem:rigidity_snd_lift}
  \lean{AlgebraicGeometry.rigidity_snd_lift}
  Let \(X\), \(Y\) be objects over \(\bar k\) and fix a \(\bar k\)-point \(x_0\) of \(X\).
  Post-composing the second projection \(p_2 \colon X \times Y \to Y\) with the slice section
  \(y \mapsto (x_0, y)\) equals the ``collapse the \(X\)-coordinate onto \(x_0\)'' endomorphism
  \((x, y) \mapsto (x_0, y)\) of \(X \times Y\). This is pure cartesian-monoidal algebra (no
  geometry): the product universal property distributes \(p_2\), the \(Y\)-component simplifies by
  the identity law, and the \(X\)-component collapses by uniqueness of maps to the terminal object.
\end{lemma}
\begin{lemma}\leanok
[Bridge~1: completeness makes the projection a closed map]
  \label{lem:snd_left_isClosedMap}
  \lean{AlgebraicGeometry.snd_left_isClosedMap}
  \uses{lem:projectiveLineBar_isProper}
  Let \(X\) be complete (proper) over \(\bar k\) and \(Y\) any object over \(\bar k\). Then the
  underlying topological map of the second projection \(p_2 \colon X \times Y \to Y\) is a
  \emph{closed} map. This is Mumford's ``completeness of \(X\) makes \(p_2\) a closed map'': the
  projection is the base change of \(X \to \Spec \bar k\) along \(Y \to \Spec \bar k\), properness
  of \(X\) gives universal closedness, and universal closedness is stable under base change, so the
  base-changed projection is universally closed and hence a closed map. Characteristic-free; no
  theorem of the cube, no cohomology.
\end{lemma}
\begin{lemma}\leanok
[Scheme-level gluing core of the Rigidity Lemma]
  \label{lem:rigidity_core}
  \lean{AlgebraicGeometry.rigidity_core}
  \uses{lem:rigidity_eqOn_dense_open}
  Let \(\bar k\) be algebraically closed, \(X\) proper, \(X \times Y\) geometrically irreducible,
  reduced and locally of finite type over \(\bar k\), and \(Z\) separated. Let
  \(f \colon X \times Y \to Z\) collapse the slice \(X \times \{y_0\}\) to a single point \(z_0\).
  Then \(f\) equals its composite with the collapse endomorphism \((x, y) \mapsto (x_0, y)\), i.e.\
  \(f = \mathtt{retract} \fatsemi f\). The proof takes the non-empty open \(U\) on which \(f\) and
  \(\mathtt{retract} \fatsemi f\) agree (\cref{lem:rigidity_eqOn_dense_open}); since \(X \times Y\)
  is geometrically irreducible over the one-point base \(\Spec \bar k\), it is irreducible, so the
  non-empty open \(U\) is dense and its inclusion is dominant; the dominant-source/separated-target
  rigidity handle then promotes agreement on \(U\) to equality of the two morphisms everywhere.
\end{lemma}
\begin{lemma}
\leanok
[Dense-open agreement: the geometric heart of the Rigidity Lemma]
  \label{lem:rigidity_eqOn_dense_open}
  \lean{AlgebraicGeometry.rigidity_eqOn_dense_open}
  \uses{lem:rigidity_eqOn_saturated_open_to_affine, lem:snd_left_isClosedMap}
  % SOURCE: [Mumford, Abelian Varieties], Rigidity Lemma (Form I.), Ch.~II \S4, book p.~43 (read from references/mumford-abelian-varieties.pdf, PDF page 54)
  % SOURCE QUOTE: "RIGIDITY LEMMA. (Form I.) Let \(X\) be a complete variety, \(Y\) and \(Z\)
  % any varieties, and \(f \colon X \times Y \to Z\) a morphism such that for some
  % \(y_0 \in Y\), \(f(X \times \{y_0\})\) is a single point \(z_0\) of \(Z\). Then there is a
  % morphism \(g \colon Y \to Z\) such that if \(p_2 \colon X \times Y \to Y\) is the
  % projection, \(f = g \circ p_2\)."
  \textit{Source: Mumford, Abelian Varieties, Ch.~II \S4, Rigidity Lemma (Form~I), p.~43 (the
  dense-open construction is the body of that proof).}
  Let \(X\) be a complete variety and \(Y\), \(Z\) any varieties over \(\bar k\). Let
  \(f \colon X \times Y \to Z\) be a morphism, fix a point \(x_0 \in X\), and suppose given
  \(y_0 \in Y\), \(z_0 \in Z\) with the \emph{collapse hypothesis}
  \[ f(X \times \{y_0\}) = \{z_0\} \qquad
     \text{(in Lean: } \mathtt{lift}\,(\mathbf 1_X)\,(\mathtt{toUnit}\,X \fatsemi y_0) \fatsemi f
     = \mathtt{toUnit}\,X \fatsemi z_0\text{).} \]
  Then there is a non-empty open subset \(U \subseteq X \times Y\) on which \(f\) agrees, as a
  scheme morphism, with the collapsed map \(\mathtt{retract} \fatsemi f\), where
  \(\mathtt{retract} := \mathtt{lift}\,(\mathtt{toUnit}\,(X \times Y) \fatsemi x_0)\,
  (p_2 \colon X \times Y \to Y)\) is the ``\((x, y) \mapsto (x_0, y)\)'' endomorphism. Concretely
  \(U = X \times V\) with \(V := Y - G\), \(G = p_2(f^{-1}(Z - U_0))\) for an affine open
  neighbourhood \(U_0\) of \(z_0\), and on \(U\) one has \(f(x, y) = f(x_0, y)\).
  \medskip
  \noindent\textbf{The collapse hypothesis is load-bearing.} The non-emptiness of \(V\) is exactly
  the assertion \(y_0 \notin G\), and this holds \emph{because} \(f(X \times \{y_0\}) = \{z_0\}
  \subseteq U_0\) keeps the slice over \(y_0\) inside the affine \(U_0\), away from \(F = Z - U_0\). A
  version of this lemma \emph{without} the collapse hypothesis is false: for \(X = Y = Z\) and
  \(f = p_1\) the first projection, \(f\) and \(\mathtt{retract} \fatsemi f \colon (x, y) \mapsto x_0\)
  agree on no non-empty open (the agreement locus is contained in \(\{x = x_0\}\), which has empty
  interior in the irreducible \(X \times Y\)). Therefore the formal target \emph{must} carry the
  collapse hypothesis as an explicit antecedent; dropping it (as an earlier decomposition
  silently did) makes the deferred goal unsatisfiable.
  \medskip
  \noindent\textbf{Formalization note: algebraically closed base, and locally-of-finite-type
  source.} The prose already takes \(\bar k\) algebraically closed, but the formal statement must
  \emph{encode} two properties. First, the cohomology-free proof of the slice-constancy bridge
  (below) needs each per-closed-point residue field \(\kappa(y)\) to be \(\bar k\) itself, so that a
  proper integral slice has global sections \(\Gamma = \bar k\) and maps to a single \(\bar k\)-point of
  the affine \(U_0\); this requires \([\mathtt{IsAlgClosed}\ \bar k]\). Second, the globalisation of the
  per-slice agreement (Step~2 of the bridge, \cref{lem:morphism_eq_of_eqAt_closedPoints}) needs the
  closed points of the saturated open \(U\) to be \emph{dense} --- i.e.\ \(U\) to be a Jacobson space.
  Over the algebraically closed \(\bar k\) this holds precisely because \(X \times Y\) is \emph{locally of
  finite type} over \(\bar k\) (\texttt{LocallyOfFiniteType.jacobsonSpace} transports the Jacobson
  property of \(\Spec \bar k\), and \texttt{closure\_closedPoints} then gives density); this requires
  \([\mathtt{LocallyOfFiniteType}\ (X \otimes Y).\mathrm{hom}]\). Accordingly the formalization carries
  \emph{both} \([\mathtt{IsAlgClosed}\ \bar k]\) and
  \([\mathtt{LocallyOfFiniteType}\ (X \otimes Y).\mathrm{hom}]\) on this lemma (in addition to
  \([\mathtt{Field}\ \bar k]\)), and propagates the same two instances to \cref{thm:rigidity_lemma},
  its core helper \texttt{rigidity\_core}, and the bridge sub-lemmas
  \cref{lem:rigidity_eqOn_saturated_open_to_affine} and
  \cref{lem:rigidity_eqAt_closedPoint_of_proper_into_affine}. Both cost nothing downstream: the
  consumers \cref{prop:morphism_P1_to_AV_constant} and \cref{prop:rigidity_genus0_curve_to_AV} (Lean
  \texttt{rigidity\_genus0\_curve\_to\_grpScheme}) \emph{already} assume
  \([\mathtt{IsAlgClosed}\ \bar k]\) --- the base-field variable is literally named \(\mathtt{kbar}\) ---
  and curves and abelian varieties are of finite type over \(\bar k\), so the
  locally-of-finite-type instance is automatic there.
\end{lemma}
% SOURCE QUOTE PROOF: "PROOF. Choose any point \(x_0 \in X\), and define \(g \colon Y \to Z\) by
% \(g(y) = f(x_0, y)\). Since \(X \times Y\) is a variety, to show that \(f = g \circ p_2\), it is
% sufficient to show that these morphisms coincide on some open subset of \(X \times Y\). Let
% \(U\) be an affine open neighbourhood of \(z_0\) in \(Z\), \(F = Z - U\), and \(G = p_2(f^{-1}(F))\);
% then \(G\) is closed in \(Y\) since \(X\) is complete and hence \(p_2\) is a closed map. Further
% \(y_0 \notin G\) since \(f(X \times \{y_0\}) = \{z_0\}\). Therefore \(Y - G = V\) is a non-empty
% open subset of \(Y\). For each \(y \in V\), the complete variety \(X \times \{y\}\) gets mapped by
% \(f\) into the affine variety \(U\), and hence to a single point of \(U\). But this means that for
% any \(x \in X\), \(y \in V\), \(f(x, y) = f(x_0, y) = g \circ p_2(x, y)\), and this proves our
% assertion."
\begin{proof}
  \leanok
  Let \(U_0\) be an affine open neighbourhood of \(z_0\) in \(Z\), put \(F = Z - U_0\), and set
  \(G = p_2(f^{-1}(F))\). Because \(X\) is complete, the projection \(p_2 \colon X \times Y \to Y\) is a
  \emph{closed} map, so \(G\) (the image of the closed set \(f^{-1}(F)\)) is closed in \(Y\). The
  collapse hypothesis \(f(X \times \{y_0\}) = \{z_0\} \subseteq U_0\) gives
  \(f^{-1}(F) \cap (X \times \{y_0\}) = \varnothing\), hence \(y_0 \notin G\); this is the only place
  the collapse hypothesis is used, and it is essential, for \(V := Y - G\) is non-empty precisely
  because \(y_0 \in V\). For each \(y \in V\) the proper variety \(X \times \{y\}\) is mapped by \(f\)
  into the affine \(U_0\), hence to a single point, necessarily \(f(x_0, y)\). So for all \(x \in X\),
  \(y \in V\), we have \(f(x, y) = f(x_0, y) = (\mathtt{retract} \fatsemi f)(x, y)\), i.e.\ \(f\) and
  \(\mathtt{retract} \fatsemi f\) agree on the non-empty open \(U := X \times V\). (In the Lean
  encoding the two slice-constancy bridges are the closed-map step from properness and the
  ``proper-into-affine has single-point image'' step; see the docstring of
  \texttt{rigidity\_core}.)
  \medskip
  \noindent\textbf{Formalization notes (the concrete char-free Mathlib route).} The textbook
  proof above is realised in Lean by three pieces. They cite Mathlib lemma names as the chosen
  route; none uses cohomology, and the relative Stein-factorisation framing is deliberately
  \emph{avoided} (see below).
  \emph{Bridge~1 (the closed-map step) --- BUILT.} ``\(X\) complete \(\Rightarrow p_2\) closed'' is
  the axiom-clean helper \texttt{snd\_left\_isClosedMap} (iter-158). It is obtained from
  \texttt{IsProper.toUniversallyClosed} on \(X\), base-changed across the canonical pullback square
  to make \(p_2\) universally closed, using that the monoidal projection \emph{is} the chosen
  pullback projection: \texttt{Over.snd\_left} rewrites \((\mathtt{snd}\,X\,Y).\mathtt{left}\) to
  \(\mathtt{Limits.pullback.snd}\ X.\mathtt{hom}\ Y.\mathtt{hom}\) exactly. This step is already
  proven.
  \emph{The fibre fact (non-emptiness of \(V\)) --- RESOLVED route, char-free.} That the slice
  \(X \times \{y_0\}\) over the \(\bar k\)-rational point \(y_0\) is exactly the image of the slice
  section \(s \colon X \to X \times Y\) is formalised \emph{without} the residue-field /
  \texttt{Triplet} / tensor-product machinery (that fine \texttt{carrierEquiv} layer is the wrong
  granularity and is to be avoided). One pastes the identity square for the section against the
  canonical pullback square (\texttt{IsPullback.of\_right},
  \texttt{IsPullback.of\_horiz\_isIso}; the section identities \(s \fatsemi p_1 = \mathbf 1\) and
  \(s \fatsemi p_2 = X.\mathtt{hom} \fatsemi y_0\) come from \texttt{lift\_fst}/\texttt{lift\_snd}
  with \texttt{Over.w} on the unit) and reads off the fibre with the \emph{coarse} topological
  layer of \texttt{PullbackCarrier}, namely
  \texttt{AlgebraicGeometry.Scheme.Pullback.image\_preimage\_eq\_of\_isPullback}. This needs no
  \([\mathtt{IsAlgClosed}]\).
  \emph{Bridge~2 (slice-constancy / the agreement equation) --- RESOLVED route, cohomology-FREE.}
  ``Proper connected slice into affine \(\Rightarrow\) single point, glued to a morphism equality on
  \(U\)'' is realised by a \textbf{global-sections + per-closed-point} route, \emph{not} by any
  relative Stein-factorisation or proper-pushforward \(f_*\mathcal O = \mathcal O\) statement: the
  latter is a genuine Mathlib gap (no Stein factorisation, no \(H^0\) base change), and route~(c)
  exists precisely to avoid coherent cohomology, so it must \emph{not} be attempted. The concrete
  stack is: for each closed point \(y \in V\) one has \(\kappa(y) = \bar k\) (alg-closed, finite type);
  the slice \(X_y\) is proper integral, so its global sections form a field
  (\texttt{isField\_of\_universallyClosed}) module-finite over \(\bar k\)
  (\texttt{finite\_appTop\_of\_universallyClosed}), and algebraic closedness kills the finite
  extension to give \(\Gamma(X_y) = \bar k\). A morphism into the affine \(U_0\) is pinned by its ring
  map on global sections (\texttt{ext\_of\_isAffine}), so \(f\) and
  \(\mathtt{retract} \fatsemi f\) agree on each closed slice. Closed points are dense in the
  locally-of-finite-type \(U\) (\texttt{closure\_closedPoints} / the Jacobson-space property of
  locally-of-finite-type \(\bar k\)-schemes; in the formalization this density is supplied by the new
  hypothesis \([\mathtt{LocallyOfFiniteType}\ (X \otimes Y).\mathrm{hom}]\)), so the per-slice
  agreement globalises to all of \(U\) via the reduced-source / separated-target rigidity
  \texttt{ext\_of\_isDominant\_of\_isSeparated'} (the same handle \texttt{rigidity\_core} already
  uses). The packaged ``morphisms agreeing on a dense set of closed points are equal''
  hom-extensionality --- the one connective Mathlib does not supply directly (it gives only the
  single-dominant-morphism \texttt{ext\_of\_isDominant}) --- has now been built as the axiom-clean
  \cref{lem:morphism_eq_of_eqAt_closedPoints} (Step~2): it assembles the closed points into one
  dominant coproduct probe and applies \texttt{ext\_of\_isDominant}. The whole argument is extracted
  into the dedicated obligation \cref{lem:rigidity_eqOn_saturated_open_to_affine} below, whose body
  wires that proven Step~2 over the per-slice Step~1
  (\cref{lem:rigidity_eqAt_closedPoint_of_proper_into_affine}); Step~1 (Mumford's per-closed-slice
  constancy) is now the chain's \emph{single} genuinely-deep residual \texttt{sorry}.
\end{proof}
\begin{lemma}
\leanok
[Bridge~2: slice-constancy on a saturated open mapping into an affine]
  \label{lem:rigidity_eqOn_saturated_open_to_affine}

  \lean{AlgebraicGeometry.rigidity_eqOn_saturated_open_to_affine}
  \uses{lem:rigidity_eqAt_closedPoint_of_proper_into_affine, lem:morphism_eq_of_eqAt_closedPoints}
  % NOTE (iter-161): the iter-160 SIGNATURE GAP is RESOLVED. The Lean proof wires Step 2
  % (morphism_eq_of_eqAt_closedPoints, PROVEN axiom-clean) over Step 1
  % (rigidity_eqAt_closedPoint_of_proper_into_affine, the residual sorry). Step 2 needs the closed
  % points of U to be DENSE, i.e. U Jacobson; over k̄ this follows from (X⊗Y) being locally of finite
  % type, so [LocallyOfFiniteType (X⊗Y).hom] is now a stated hypothesis of this lemma (and across the
  % chain). With it, the `JacobsonSpace U` instance is a routine discharge (LocallyOfFiniteType.
  % jacobsonSpace + JacobsonSpace.of_isOpenEmbedding), no longer an as-typed gap. The only `sorry`
  % reachable from this lemma is now Step 1 (Mumford's per-closed-slice constancy). The hypothesis
  % set is recorded in the statement and proof prose below (both [IsAlgClosed k̄] AND
  % [LocallyOfFiniteType (X⊗Y).hom]); the earlier "[IsAlgClosed] is the only added instance" claim
  % is retired.
  % SOURCE: [Mumford, Abelian Varieties], Rigidity Lemma (Form I.), Ch.~II \S4, book p.~43 (read from references/mumford-abelian-varieties.pdf, PDF page 54)
  % SOURCE QUOTE PROOF: "For each \(y \in V\), the complete variety \(X \times \{y\}\) gets mapped by
  % \(f\) into the affine variety \(U\), and hence to a single point of \(U\). But this means that for
  % any \(x \in X\), \(y \in V\), \(f(x, y) = f(x_0, y) = g \circ p_2(x, y)\), and this proves our
  % assertion." [Mumford, Ch.~II \S4, Rigidity Lemma, p.~43.]
  \textit{Source: Mumford, Abelian Varieties, Ch.~II \S4, Rigidity Lemma (Form~I), p.~43 (the
  ``for each \(y \in V\), the complete slice maps into the affine, hence to a single point'' step).}
  Let \(\bar k\) be algebraically closed and let \(X\), \(Y\), \(Z\) be objects of
  \(\mathrm{Over}\,(\Spec \bar k)\) with \(X\) proper, \(X \times Y\) geometrically irreducible, reduced,
  and \emph{locally of finite type} over \(\bar k\) (the formalization carries
  \([\mathtt{LocallyOfFiniteType}\ (X \otimes Y).\mathrm{hom}]\), which makes the saturated open \(U\)
  below a Jacobson space, so its closed points are dense --- the hypothesis Step~2 of the proof
  consumes), and \(Z\) separated. Let \(f \colon X \times Y \to Z\) be a morphism, fix a \(\bar k\)-point
  \(x_0 \colon \mathbf 1 \to X\), and write \(\mathtt{retract} := \mathtt{lift}\,
  (\mathtt{toUnit}\,(X \times Y) \fatsemi x_0)\,(p_2)\) for the ``\((x, y) \mapsto (x_0, y)\)''
  endomorphism. Suppose:
  \begin{itemize}
    \item \(U \subseteq X \times Y\) is a \(p_2\)-\emph{saturated} open, i.e.\ \(U = p_2^{-1}(V_{\mathrm{set}})\)
      for a set \(V_{\mathrm{set}} \subseteq Y\) (so \(U\) contains the whole fibre \(X_y\) over each of
      its points);
    \item \(U_0 \subseteq Z\) is an \emph{affine} open;
    \item \(f(U) \subseteq U_0\), i.e.\ \(f\) maps \(U\) into the affine \(U_0\).
  \end{itemize}
  Then \(f\) and \(\mathtt{retract} \fatsemi f\) agree on \(U\) as scheme morphisms:
  \(U.\iota \fatsemi f = U.\iota \fatsemi (\mathtt{retract} \fatsemi f)\).
\end{lemma}
\begin{proof}
  \leanok
  This is the formalization realization of Mumford's step ``for each \(y \in V\), the complete
  variety \(X \times \{y\}\) gets mapped by \(f\) into the affine variety \(U\), and hence to a single
  point of \(U\)''. We realise it by a \textbf{global-sections + per-closed-point} route that is
  \emph{cohomology-free}; in particular the relative Stein-factorisation / proper-pushforward
  \(f_*\mathcal O = \mathcal O\) framing is a confirmed Mathlib gap (no Stein factorisation, no
  \(H^0\) base change) and is \emph{deliberately avoided}. The argument splits into two named
  sub-lemmas, one residual and one proven.
  \medskip
  \noindent\textbf{Step 1 (per-closed-slice constancy) --- \cref{lem:rigidity_eqAt_closedPoint_of_proper_into_affine}, the residual.}
  Fix a closed point \(x\) of \(U\), lying over a closed point \(y = p_2(x) \in V_{\mathrm{set}}\).
  Since \(\bar k\) is algebraically closed and \(Y\) is of finite type over \(\bar k\), the residue field
  is \(\kappa(y) = \bar k\). Because \(U = p_2^{-1}(V_{\mathrm{set}})\) is saturated, \(U\) contains the
  \emph{entire} fibre \(X_y \cong X\) over \(y\); hence \(f\) maps the proper integral slice \(X_y\) into the
  affine \(U_0\). By \(\mathtt{isField\_of\_universallyClosed}\) the global sections \(\Gamma(X_y)\) form a
  field, module-finite over \(\bar k\) by \(\mathtt{finite\_appTop\_of\_universallyClosed}\), and
  algebraic closedness collapses the finite extension to give \(\Gamma(X_y) = \bar k\). A morphism into
  the affine \(U_0\) is pinned by its ring map on global sections (\(\mathtt{ext\_of\_isAffine}\)), so the
  restriction of \(f\) to \(X_y\) factors through a single \(\bar k\)-point of \(U_0\) --- necessarily the
  point \(f(x_0, y) = (\mathtt{retract} \fatsemi f)(x)\). Thus \(f\) and \(\mathtt{retract} \fatsemi f\)
  agree (after the residue-field probe) at every closed point of \(U\). This per-slice statement is the
  chain's single genuinely-deep residual \texttt{sorry}.
  \medskip
  \noindent\textbf{Step 2 (globalisation over dense closed points) --- \cref{lem:morphism_eq_of_eqAt_closedPoints}, proven.}
  The closed points are dense in the locally-of-finite-type \(\bar k\)-scheme \(U\)
  (\(\mathtt{closure\_closedPoints}\) / the Jacobson-space property, supplied here by the new hypothesis
  \([\mathtt{LocallyOfFiniteType}\ (X \otimes Y).\mathrm{hom}]\)). To turn ``\(f\) and
  \(\mathtt{retract} \fatsemi f\) agree at each closed point'' into a single morphism equality on \(U\),
  the proven globalisation \cref{lem:morphism_eq_of_eqAt_closedPoints} assembles the closed points
  into one dominant probe --- the coproduct
  \(\coprod_{x \in \mathrm{closedPoints}\,U} \Spec \kappa(x) \to U\), whose range is dense --- and feeds
  it to the reduced-source / separated-target rigidity handle \(\mathtt{ext\_of\_isDominant}\) (the same
  family \(\mathtt{rigidity\_core}\) uses). This yields the morphism equality
  \(U.\iota \fatsemi f = U.\iota \fatsemi (\mathtt{retract} \fatsemi f)\) on all of \(U\). The
  ``agrees on a dense set of closed points \(\Rightarrow\) hom-equality'' connective is the one piece
  Mathlib does not package directly (it supplies only the single-dominant-morphism
  \(\mathtt{ext\_of\_isDominant}\)); that bespoke build is now done in
  \cref{lem:morphism_eq_of_eqAt_closedPoints}, and it is \emph{not} cohomology. \qedhere
\end{proof}
\begin{lemma}\leanok
[Step~2: morphisms agreeing at every closed point of a Jacobson source are equal]
  \label{lem:morphism_eq_of_eqAt_closedPoints}
  
  \lean{AlgebraicGeometry.morphism_eq_of_eqAt_closedPoints}
  % Project-bespoke assembly: NO external source. Mathlib supplies only the single-dominant-morphism
  % ext_of_isDominant; the "dense closed points => hom-extensionality" connective is built here.
  Let \(W\), \(Z\) be schemes with \(W\) \emph{reduced}, \(W\) a \emph{Jacobson space}, and \(Z\)
  \emph{separated}. Let \(g_1, g_2 \colon W \to Z\) be morphisms. If, for every \emph{closed} point
  \(x \in W\), the two morphisms agree after restriction along the residue-field point
  \(\Spec \kappa(x) \to W\) (in Lean, \(W.\mathtt{fromSpecResidueField}\,x \fatsemi g_1 =
  W.\mathtt{fromSpecResidueField}\,x \fatsemi g_2\)), then \(g_1 = g_2\).
\end{lemma}
\begin{proof}
\leanok
  This is the ``dense closed points \(\Rightarrow\) hom-extensionality'' connective that Mathlib does
  not package directly: it supplies only the single-dominant-morphism extensionality
  \(\mathtt{ext\_of\_isDominant}\) (a dominant morphism into a separated target is an epimorphism for
  the purpose of right-cancellation). We assemble the per-closed-point hypotheses into one such
  dominant probe.
  Form the coproduct \(P := \coprod_{x \in \mathrm{closedPoints}\,W} \Spec \kappa(x)\) and the morphism
  \(\mathrm{probe} \colon P \to W\) defined by \(\mathtt{Sigma.desc}\) of the residue-field points
  \(W.\mathtt{fromSpecResidueField}\,x\). Its topological range contains every closed point of \(W\):
  the inclusion \(\iota_x \colon \Spec \kappa(x) \to P\) composes with \(\mathrm{probe}\) to
  \(W.\mathtt{fromSpecResidueField}\,x\) (the defining property of \(\mathtt{Sigma.desc}\)), and the
  range of \(W.\mathtt{fromSpecResidueField}\,x\) is the singleton \(\{x\}\)
  (\(\mathtt{Scheme.range\_fromSpecResidueField}\)). Because \(W\) is a Jacobson space, the closed points
  are dense (\(\mathtt{closure\_closedPoints}\), i.e.\ \(\overline{\mathrm{closedPoints}\,W} = W\)), so
  the range of \(\mathrm{probe}\) is dense and \(\mathrm{probe}\) is \emph{dominant}.
  Componentwise, \(g_1\) and \(g_2\) agree after \(\mathrm{probe}\): precomposing with each coproduct
  inclusion \(\iota_x\) reduces (by \(\mathtt{Sigma.\iota\_desc}\)) to
  \(W.\mathtt{fromSpecResidueField}\,x \fatsemi g_1 = W.\mathtt{fromSpecResidueField}\,x \fatsemi g_2\),
  which is the hypothesis at the closed point \(x\). By the universal property of the coproduct
  (\(\mathtt{Sigma.hom\_ext}\)) this gives \(\mathrm{probe} \fatsemi g_1 = \mathrm{probe} \fatsemi g_2\).
  Since \(\mathrm{probe}\) is dominant and \(Z\) is separated, \(\mathtt{ext\_of\_isDominant}\) cancels it
  on the left to yield \(g_1 = g_2\).
\end{proof}
\begin{lemma}
\leanok
[Global-sections constancy: two \(\bar k\)-points of a proper integral source into an affine agree]
  \label{lem:eq_comp_of_isAffine_of_properIntegral}
  \lean{AlgebraicGeometry.eq_comp_of_isAffine_of_properIntegral}
  % Project-bespoke assembly of Mathlib lemmas: NO external source. Realises, cohomology-free, the
  % classical "a global regular function on a proper integral variety over an algebraically closed
  % field is constant" --- precisely the affine-image step of Mumford's Rigidity Lemma proof quoted
  % above (the "complete slice into an affine maps to a single point" sentence, Ch.~II \S4, p.~43).
  Let \(\bar k\) be algebraically closed and let \(W\) be an integral scheme equipped with a structure
  morphism \(w \colon W \to \Spec \bar k\) that is universally closed (proper) and locally of finite
  type. Let \(g \colon W \to V\) be a morphism into an \emph{affine} scheme \(V\). Then any two
  \(\bar k\)-points of \(W\) --- i.e.\ any two sections \(a, b \colon \Spec \bar k \to W\) of \(w\) (so
  \(a \fatsemi w = \mathbf 1\) and \(b \fatsemi w = \mathbf 1\)) --- satisfy \(a \fatsemi g = b \fatsemi g\).
  \medskip
  \noindent\textbf{Every hypothesis is load-bearing} (confirmed by the iter-161 audit).
  \([\mathtt{IsAlgClosed}\ \bar k]\) collapses the finite extension \(\Gamma(W)/\bar k\); integrality of
  \(W\) makes \(\Gamma(W)\) a domain (without it a two-point disjoint union is a counterexample);
  \(\mathtt{UniversallyClosed}\) is what upgrades \(\Gamma(W)\) to a \emph{field} (without it
  \(\mathbb A^1\), with \(\Gamma = \bar k[t]\) and two distinct \(\bar k\)-points, is a counterexample);
  \(\mathtt{LocallyOfFiniteType}\) gives finiteness of \(\Gamma(W)\) over \(\bar k\); and
  \([\mathtt{IsAffine}\ V]\) is what lets \(\mathtt{ext\_of\_isAffine}\) pin the morphism by its ring map
  on global sections.
\end{lemma}
\begin{proof}
\leanok
  Because \(W\) is integral and \(w\) is universally closed, the global sections
  \(\Gamma(W, \mathcal O_W)\) form a field (\(\mathtt{isField\_of\_universallyClosed}\)), and the structure
  ring map \(w^\sharp \colon \bar k \to \Gamma(W)\) (the map \(\mathtt{wk.appTop}\) after the canonical
  identification \(\Gamma(\Spec \bar k) \cong \bar k\)) is integral / module-finite
  (\(\mathtt{finite\_appTop\_of\_universallyClosed}\)). Since \(\bar k\) is algebraically closed and
  \(\Gamma(W)\) is a domain finite over it, that map is bijective
  (\(\mathtt{IsAlgClosed.ringHom\_bijective\_of\_isIntegral}\)); thus \(\Gamma(W) = \bar k\) and
  \(w^\sharp\) is an isomorphism on global sections --- in particular an epimorphism. Both sections
  \(a\), \(b\) left-invert \(w^\sharp\) on global sections (they are sections of \(w\)), so they induce the
  \emph{same} ring map \(a^\sharp = b^\sharp\) on \(\Gamma(W)\) (cancel the epimorphism \(w^\sharp\)). A
  morphism into the \emph{affine} \(V\) is pinned by its global-sections ring map
  (\(\mathtt{ext\_of\_isAffine}\)), whence \(a \fatsemi g = b \fatsemi g\). This realises, without
  cohomology, the classical ``proper connected variety into an affine is a single point'' step
  quoted from Mumford above.
\end{proof}
\begin{lemma}
\leanok
[Integrality descends to a retract of an integral scheme]
  \label{lem:isIntegral_of_retract_of_integral}
  \lean{AlgebraicGeometry.isIntegral_of_retract}
  % NOTE: (iter-162 review) project-bespoke (elementary topology + algebra); NO external source.
  % LANDED axiom-clean as `AlgebraicGeometry.isIntegral_of_retract`, stated generically for any
  % retract `r : S ⟶ T`, `pr : T ⟶ S` with `r ≫ pr = 𝟙 S` of an integral `T` (more general than
  % the `X`/`X×Y`/`p_1` instance below). PROOF DIVERGENCE (cosmetic, sound either way): the Lean
  % proves reducedness PER-STALK (each `S.stalk x` embeds via `pr.stalkMap` into the reduced
  % `T.stalk`, then `isReduced_of_isReduced_stalk`), not via the global-sections split-injection
  % the prose below describes. Plan/blueprint-writer may align the prose to the stalk route.
  Let \(X \times Y\) be an integral scheme (here it is geometrically irreducible and reduced over
  \(\bar k\)). Suppose \(X\) is a \emph{retract} of \(X \times Y\): there is a section
  \(s \colon X \to X \times Y\) of the first projection \(p_1\), i.e.\ \(s \fatsemi p_1 = \mathbf 1_X\).
  Then \(X\) is integral.
\end{lemma}
\begin{proof}
\leanok
  Integrality is reducedness together with irreducibility, and the retract supplies each.
  \emph{Reduced.} The section identity \(s \fatsemi p_1 = \mathbf 1_X\) makes the pullback on global
  sections \(p_1^\sharp \colon \mathcal O_X(X) \to \mathcal O_{X \times Y}(X \times Y)\) split
  injective (it is left-inverted by \(s^\sharp\)). A split-injective ring map into the reduced ring
  \(\mathcal O_{X \times Y}(X \times Y)\) embeds \(\mathcal O_X(X)\) into a reduced ring, so it too is
  reduced (\(\mathtt{isReduced\_of\_injective}\) on the split injection); hence \(X\) is reduced.
  \emph{Irreducible.} The first projection \(p_1 \colon X \times Y \to X\) is surjective because it
  admits the section \(s\) (every point of \(X\) is \(p_1(s(\cdot))\)). The continuous image of an
  irreducible space under a surjection is irreducible; as \(X \times Y\) is irreducible and \(p_1\) is a
  continuous surjection onto \(X\), the space of \(X\) is irreducible.
  Reduced and irreducible together give that \(X\) is integral. This is elementary --- no cohomology,
  no external citation --- and feeds the Step-1 geometric assembly below, where the proper slice
  \(X_y \cong X\) must be presented as proper \emph{integral}.
\end{proof}
\begin{lemma}\leanok
[Step~1: per-closed-slice constancy of a proper slice into an affine]
  \label{lem:rigidity_eqAt_closedPoint_of_proper_into_affine}

  \lean{AlgebraicGeometry.rigidity_eqAt_closedPoint_of_proper_into_affine}
  \uses{lem:eq_comp_of_isAffine_of_properIntegral, lem:isIntegral_of_retract_of_integral}
  % SOURCE: [Mumford, Abelian Varieties], Rigidity Lemma (Form I.), Ch.~II \S4, book p.~43 (read from references/mumford-abelian-varieties.pdf, PDF page 54)
  % SOURCE QUOTE PROOF: "For each \(y \in V\), the complete variety \(X \times \{y\}\) gets mapped by
  % \(f\) into the affine variety \(U\), and hence to a single point of \(U\). But this means that for
  % any \(x \in X\), \(y \in V\), \(f(x, y) = f(x_0, y) = g \circ p_2(x, y)\), and this proves our
  % assertion." [Mumford, Ch.~II \S4, Rigidity Lemma, p.~43.]
  \textit{Source: Mumford, Abelian Varieties, Ch.~II \S4, Rigidity Lemma (Form~I), p.~43 (the
  ``for each \(y \in V\), the complete slice maps into the affine, hence to a single point'' step).}
  With the data of \cref{lem:rigidity_eqOn_saturated_open_to_affine} --- \(\bar k\) algebraically
  closed, \(X\) proper over \(\bar k\), \(X \times Y\) geometrically irreducible, reduced, and locally of
  finite type, \(Z\) separated, a \(\bar k\)-point \(x_0 \colon \mathbf 1 \to X\), a morphism
  \(f \colon X \times Y \to Z\), a \(p_2\)-saturated open \(U = p_2^{-1}(V_{\mathrm{set}})\) mapping into an
  affine open \(U_0 \subseteq Z\) --- fix a \emph{closed} point \(x\) of \(U\). Then \(f\) and the collapsed
  map \(\mathtt{retract} \fatsemi f\) (with \(\mathtt{retract} := \mathtt{lift}\,
  (\mathtt{toUnit}\,(X \times Y) \fatsemi x_0)\,(p_2)\), i.e.\ \((x, y) \mapsto (x_0, y)\)) agree at \(x\)
  after the residue-field probe \(\Spec \kappa(x) \to U\):
  \[ U.\mathtt{fromSpecResidueField}\,x \fatsemi (U.\iota \fatsemi f)
     = U.\mathtt{fromSpecResidueField}\,x \fatsemi (U.\iota \fatsemi (\mathtt{retract} \fatsemi f)). \]
\end{lemma}
\begin{proof}
  \leanok
  This is Mumford's per-slice step, realised \emph{cohomology-free} (route~B): the relative Stein /
  proper-pushforward \(f_*\mathcal O = \mathcal O\) framing is a confirmed Mathlib gap and is
  deliberately avoided. The closed point \(x\) of \(U\) lies over a closed point
  \(y = p_2(x) \in V_{\mathrm{set}}\); since \(\bar k\) is algebraically closed and \(Y\) is of finite type
  over \(\bar k\), the residue field is \(\kappa(y) = \bar k\). Saturation \(U = p_2^{-1}(V_{\mathrm{set}})\)
  puts the \emph{entire} proper integral fibre \(X_y \cong X\) over \(y\) inside \(U\), so \(f\) maps \(X_y\)
  into the affine \(U_0\). Because \(X_y\) is proper integral over \(\bar k\), its global sections form a
  field (\(\mathtt{isField\_of\_universallyClosed}\)), module-finite over \(\bar k\)
  (\(\mathtt{finite\_appTop\_of\_universallyClosed}\)); algebraic closedness collapses the finite
  extension, giving \(\Gamma(X_y, \mathcal O) = \bar k\).
  % NOTE (iter-162): the iter-161 documentary gap is RESOLVED --- this algebraic core ("two k-bar
  % points of a proper integral l.f.t. k-bar-scheme into an affine agree") now has its own node
  % lem:eq_comp_of_isAffine_of_properIntegral (above), and this proof's  edge points at it.
  A morphism into the affine \(U_0\) is pinned by
  its ring map on global sections (\(\mathtt{ext\_of\_isAffine}\)), so \(f|_{X_y}\) factors through a
  single \(\bar k\)-point of \(U_0\). That point is \(f(x_0, y)\), since the \(\bar k\)-point \((x_0, y)\) lies
  on \(X_y\); and \(f(x_0, y) = (\mathtt{retract} \fatsemi f)(x)\) by the definition of \(\mathtt{retract}\).
  Evaluating at the residue-field probe of the closed point \(x\) gives the asserted equality. This
  per-slice statement is the chain's single genuinely-deep residual \texttt{sorry}; its globalisation
  over the dense closed points is \cref{lem:morphism_eq_of_eqAt_closedPoints}.
\end{proof}
\section{The Milne \S I.3 chain: additivity, homomorphisms, extension}
This section assembles, from the proven Rigidity Lemma, the genus-\(0\) base case. The two
load-bearing ingredients of the \textbf{primary route} are the additive decomposition over a
product (\cref{lem:hom_additivity_over_product}, Milne Corollary~1.5) and the concrete \(\mathbb
G_m\)-scaling action \(\sigma_\times\) on \(\mathbb P^1\) (\cref{def:gaTranslationP1}); from these
the base case ``every morphism \(\mathbb P^1 \to A\) is constant'' falls out \emph{in one shot}
(\cref{prop:morphism_P1_to_AV_constant}): the scaling fixed point \(0\) collapses one axis of the
decomposition, giving \(f(\lambda x) = f(x)\) and hence constancy on the dense \(\mathbb G_m\). The
classical \(\mathbb G_a\)-additive homomorphism step
(\cref{lem:hom_from_Ga_trivial}) and the triviality lemma \(\mathrm{Hom}(\mathbb G_a, A) = 0\)
(\cref{lem:hom_Ga_to_av_trivial}) are retained as the Milne-faithful alternative but are
\emph{not} on the genus-\(0\) critical path. None of this uses the theorem of the cube, and --- crucially
--- none of it uses Milne's Theorem~3.2 / Lemma~3.3 (the rational-map extension): because \(f
\colon \mathbb P^1 \to A\) is \emph{already} a total morphism out of the complete \(\mathbb
P^1\), there is no defect map to extend over a complete surface (see
\cref{rmk:base_case_fourth_route}). Corollary~1.2 (\cref{lem:av_regular_map_is_hom}) and
Theorem~3.2 (\cref{thm:rational_map_to_av_extends}) are retained in this chapter only as
Route-A inputs (the Albanese universal property) and are explicitly \emph{off} the genus-\(0\)
path. The genus-\(0\) output is exactly the content of Milne's Proposition~3.9, ``every rational
(here: regular) map \(\mathbb P^1 \to A\) is constant''.
\begin{remark}
  \label{rmk:cube_not_needed}
  \textbf{The theorem of the cube is not needed on the genus-\(0\) path.} An earlier draft of
  this chapter recorded the theorem of the cube as a deferred prerequisite of the single-curve
  base case ``\(\mathbb P^1 \to A\) constant''. That was a misreading. The base case
  (Proposition~3.9, the \(d = 1\) case of Proposition~3.10) follows directly from the Rigidity
  Theorem~1.1 --- our \cref{thm:rigidity_lemma} --- via Corollary~1.5 (additivity,
  \cref{lem:hom_additivity_over_product}) composed with the total \(\mathbb G_m\)-scaling
  action \(\sigma_\times\) (\cref{def:gaTranslationP1}): the scaling fixed point \(0\) collapses one
  axis of the decomposition, yielding \(f(\lambda x) = f(x)\) and hence constancy on the dense
  \(\mathbb G_m\) (the primary \(\mathbb G_m\)-scaling shortcut; see
  \cref{rmk:base_case_fourth_route}). This route does not even use the triviality of
  \(\mathrm{Hom}(\mathbb G_a, A)\) (\cref{lem:hom_Ga_to_av_trivial}), which is retained only for the
  Milne-faithful \(\mathbb G_a\)-additive alternative. Neither the theorem of the cube nor Milne's Theorem~3.2 /
  Lemma~3.3 (rational-map extension) is used: since \(f \colon \mathbb P^1 \to A\) is already a
  total morphism out of the complete \(\mathbb P^1\), no defect map needs extending over a complete
  surface (this corrects the earlier claim that the surface extension was on the genus-\(0\) path;
  see \cref{rmk:thm32_codim1_mathlib_gap}). The theorem of the cube (Milne \S I.5) enters only for
  projectivity, the seesaw/square principle, dual abelian varieties, polarizations, and the
  Poincar\'e sheaf --- none on the genus-\(0\) path. For Route~A: the Albanese universal property
  that the positive-genus arm needs (Milne Proposition~6.1/6.4, \S III.6) is \emph{also} cube-free,
  resting on Theorem~3.2 (\cref{thm:rational_map_to_av_extends}) + Corollary~1.2/1.5 --- which is
  precisely why Theorem~3.2 and Corollary~1.2 are retained in this chapter as Route-A inputs.
\end{remark}
\begin{lemma}
\leanok
[Additive decomposition over a product (Milne Corollary 1.5)]
  \label{lem:hom_additivity_over_product}
  
  \uses{thm:rigidity_lemma}
  \lean{AlgebraicGeometry.hom_additive_decomp_of_rigidity}
  % NOTE: (iter-163 review) The Lean signature carries three product-level instance
  % hypotheses on V⊗W not spelled out in the prose: [GeometricallyIrreducible (V⊗W).hom],
  % [LocallyOfFiniteType (V⊗W).hom], [IsReduced (V⊗W).left] — i.e. "V×W is a variety".
  % These are mathematically free for genuine complete varieties V, W, but Mathlib does not
  % auto-derive the product instances from the factors, so they are carried explicitly to keep
  % the lemma axiom-clean. Also: the Lean states the EXISTENCE half only (canonical
  % axis-restriction witnesses inlined; uniqueness not formalized) and requires only V (not W)
  % complete — strictly more general than the prose. (lean-vs-blueprint-checker avr-iter163.)
  % NOTE: (iter-164 review) iter-164 hygiene dropped the A-side instances [Smooth A.hom] and
  % [GeometricallyIrreducible A.hom] from the Lean signature; the abelian-variety target now
  % requires only [GrpObj A] [IsProper A.hom]. The Lean signature is therefore STRICTLY MORE
  % GENERAL than the prose ("Let A be an abelian variety"): every situation the prose covers is
  % covered, and additionally any proper group object suffices. Axiom-clean preserved.
  % (lean-vs-blueprint-checker avr-iter164.)
  % SOURCE: [Milne, Abelian Varieties], Corollary 1.5, \S I.1, doc p.~9 (read from references/abelian-varieties.pdf, PDF page 15)
  % SOURCE QUOTE: "COROLLARY 1.5. Let V and W be complete varieties over k with k-rational
  % points v0 and w0, and let p and q be the projection maps V×W→V and V×W→W. Let A be an
  % abelian variety. Then a morphism h:V×W→A such that h(v0,w0)=0 can be written uniquely as
  % h = f∘p + g∘q with f:V→A and g:W→A morphisms such that f(v0)=0 and g(w0)=0."
  \textit{Source: Milne, Abelian Varieties, Corollary~1.5.}
  Let \(V\), \(W\) be complete varieties over \(\bar k\) with \(\bar k\)-points \(v_0 \in V\),
  \(w_0 \in W\), and let \(p \colon V \times W \to V\), \(q \colon V \times W \to W\) be the
  projections. Let \(A\) be an abelian variety over \(\bar k\). Then a morphism
  \(h \colon V \times W \to A\) with \(h(v_0, w_0) = 0\) can be written uniquely as
  \(h = f \circ p + g \circ q\), with morphisms \(f \colon V \to A\), \(g \colon W \to A\) satisfying
  \(f(v_0) = 0\) and \(g(w_0) = 0\).
  \emph{Lean encoding.} In the project \(A\) is a group object \texttt{[GrpObj A]} (not assumed
  commutative). The ``\(+\)'' here is the binary operation that \texttt{GrpObj} induces on each
  hom-set \(\mathrm{Hom}(X, A)\): for \(u, v \colon X \to A\), \(u + v\) is the composite
  \(X \xrightarrow{\langle u, v\rangle} A \times A \xrightarrow{\;\mathrm{mul}\;} A\) (and
  \(f \circ p + g \circ q\) is built this way from the projections). The decomposition is therefore
  read \emph{multiplicatively} as \(h = (f \circ p) \cdot (g \circ q)\) in the group
  \(\mathrm{Hom}(V \times W, A)\), and does \emph{not} require \(A\) commutative; commutativity of
  abelian varieties (Milne Cor 1.4, a separate consequence of \cref{lem:av_regular_map_is_hom})
  is needed only to rewrite the product symmetrically/additively, not for the statement or its
  proof. The prover should infer the hom-set group operation from the \texttt{GrpObj} instance,
  not from a \texttt{CommGrpObj} instance.
\end{lemma}
% SOURCE QUOTE PROOF: "PROOF. Set f=h|V×{w0}, g=h|{v0}×W, and identify V×{w0} and {v0}×W with
% V and W. On points, f(v) = h(v,w0) and g(w) = h(v0,w), and so φ =df h − (f∘p + g∘q) is the
% map that sends (v,w) ↦ h(v,w) − h(v,w0) − h(v0,w). Thus φ(V×{w0}) = 0 = φ({v0}×W) and so the
% theorem shows that φ = 0." [Milne, Cor 1.5, \S I.1, doc p.~9--10.]
\begin{proof}
  \leanok
  \textit{Source: Milne, Abelian Varieties, Corollary~1.5.}
  Define \(f := h|_{V \times \{w_0\}}\) and \(g := h|_{\{v_0\} \times W}\), identifying
  \(V \times \{w_0\} \cong V\) and \(\{v_0\} \times W \cong W\); on points \(f(v) = h(v, w_0)\) and
  \(g(w) = h(v_0, w)\). Form the difference
  \[ \varphi := h - (f \circ p + g \circ q) \colon V \times W \to A, \qquad
     \varphi(v, w) = h(v, w) - h(v, w_0) - h(v_0, w) \]
  (the subtraction is in the group \(A\), i.e.\ \(\varphi\) is the composite of
  \((h, -(f \circ p + g \circ q))\) with the group law of \(A\)). Then \(\varphi\) vanishes on both
  axes: \(\varphi(V \times \{w_0\}) = 0 = \varphi(\{v_0\} \times W)\) (direct computation, using
  \(h(v_0, w_0) = 0\)). In particular \(\varphi\) collapses the complete axis \(V \times \{w_0\}\) to
  the single point \(0 \in A\), so by the Rigidity Lemma (\cref{thm:rigidity_lemma}, with \(V\)
  complete and \(V \times W\) geometrically irreducible) \(\varphi\) factors through the projection
  \(q\); but it also vanishes on \(\{v_0\} \times W\), which is a section of \(q\), so the factoring
  morphism is identically \(0\), i.e.\ \(\varphi \equiv 0\). Hence \(h = f \circ p + g \circ q\).
  Uniqueness: restricting any such decomposition to the axes recovers \(f\) and \(g\) as the
  axis-restrictions of \(h\).
\end{proof}
\begin{lemma}
\leanok
[A pointed regular map of abelian varieties is a homomorphism (Milne Corollary 1.2)]
  \label{lem:av_regular_map_is_hom}

  \uses{lem:hom_additivity_over_product}
  \lean{AlgebraicGeometry.av_regularMap_isHom_of_zero}
  % NOTE: (iter-163 review) The Lean formalizes the POINTED direction
  % (η[A] ≫ α = η[B] ⟹ IsMonHom α), not the full "composite of a homomorphism with a
  % translation" form — which is what the genus-0 consumer needs. The conclusion IsMonHom α
  % (one_hom + mul_hom) is a faithful, slightly stronger rendering of "α is a homomorphism".
  % As with lem:hom_additivity_over_product, the signature carries the product-level instances
  % [GeometricallyIrreducible (A⊗A).hom], [LocallyOfFiniteType (A⊗A).hom], [IsReduced (A⊗A).left]
  % ("A×A is a variety"), free for a genuine abelian variety A but carried explicitly because
  % Mathlib lacks the product-instance inference. (lean-vs-blueprint-checker avr-iter163.)
  % NOTE: (iter-164 review) iter-164 hygiene dropped the B-side instances [Smooth B.hom] and
  % [GeometricallyIrreducible B.hom] from the Lean signature; the target B now requires only
  % [GrpObj B] [IsProper B.hom]. (The source-side A retains the full abelian-variety instance
  % set.) The Lean signature is strictly more general on B than the prose ("α : A → B of
  % abelian varieties") demands. The lean-auditor flagged the symmetric A-side
  % [Smooth A.hom]/[GeometricallyIrreducible A.hom] as ALSO unused — kept this iter as a
  % minimal-change follow-up. (lean-vs-blueprint-checker avr-iter164.)
  % SOURCE: [Milne, Abelian Varieties], Corollary 1.2, \S I.1, doc p.~9 (read from references/abelian-varieties.pdf, PDF page 15)
  % SOURCE QUOTE: "COROLLARY 1.2. Every regular map α:A→B of abelian varieties is the composite
  % of a homomorphism with a translation."
  \textit{Source: Milne, Abelian Varieties, Corollary~1.2.}
  Every regular map \(\alpha \colon A \to B\) of abelian varieties over \(\bar k\) is the composite
  of a homomorphism with a translation; equivalently, if \(\alpha(0_A) = 0_B\) then \(\alpha\) is a
  homomorphism, i.e.\ \(\alpha(a + a') = \alpha(a) + \alpha(a')\) for all \(a, a' \in A\).
\end{lemma}
% SOURCE QUOTE PROOF: "PROOF. The regular map α will send the k-rational point 0 of A to a
% k-rational point b of B. After composing α with translation by −b, we may assume that
% α(0)=0. Consider the map φ:A×A→B, φ(a,a′)=α(a+a′) − α(a) − α(a′). By this we mean that φ is
% the difference of the two regular maps A×A →(m)→ A →(α)→ B and A×A →(α×α)→ B×B →(m)→ B, which
% is a regular map. Then φ(A×0)=0=φ(0×A) and so φ=0. This means that α is a homomorphism."
% [Milne, Cor 1.2, \S I.1, doc p.~9.]
\begin{proof}
  \leanok
  \textit{Source: Milne, Abelian Varieties, Corollary~1.2.}
  The image \(\alpha(0_A)\) is a \(\bar k\)-point \(b\) of \(B\); composing \(\alpha\) with translation by
  \(-b\) we may assume \(\alpha(0_A) = 0_B\), and it suffices to show such an \(\alpha\) is a
  homomorphism. Form
  \[ \varphi \colon A \times A \to B, \qquad \varphi(a, a') = \alpha(a + a') - \alpha(a) - \alpha(a'), \]
  the difference of the two regular maps \(A \times A \xrightarrow{m} A \xrightarrow{\alpha} B\)
  and \(A \times A \xrightarrow{\alpha \times \alpha} B \times B \xrightarrow{m} B\). Then
  \(\varphi(0_A, 0_A) = 0\) and \(\varphi\) vanishes on both axes,
  \(\varphi(A \times 0) = 0 = \varphi(0 \times A)\) (since \(\alpha(0) = 0\)). By the additive
  decomposition \cref{lem:hom_additivity_over_product} (with \(V = W = A\)), the axis-restrictions
  of \(\varphi\) being \(0\) force \(\varphi \equiv 0\); equivalently, \(\varphi\) collapses the complete
  axis \(A \times 0\) and vanishes on \(0 \times A\), so Rigidity gives \(\varphi = 0\). Hence
  \(\alpha(a + a') = \alpha(a) + \alpha(a')\), i.e.\ \(\alpha\) is a homomorphism.
\end{proof}
\begin{remark}
  \label{rmk:thm32_codim1_mathlib_gap}
  \textbf{Formalisation risk: the pure-codimension-\(1\) step (Lemma 3.3) has no Mathlib
  Weil-divisor support.} Theorem~3.1's everywhere-extension-outside-codimension-\(2\) is within
  reach of Mathlib's valuative criterion of properness
  (\texttt{ValuativeCriterion}/\texttt{IsProper.of\_valuativeCriterion}), applied at the
  discrete valuation ring \(\mathcal O_{V, Z}\) of each codimension-\(1\) point \(Z\). But Lemma~3.3,
  which rules out codimension-\(1\) indeterminacy for maps into a group variety, is proved in
  Milne via Weil-divisor theory (\(\mathrm{div}(f)_0 - \mathrm{div}(f)_1\), prime divisors,
  tangent cones) for which Mathlib currently has \emph{no} usable API at this generality. This
  is the route's riskiest sub-build. An open question for the prover: whether a \emph{pointwise}
  valuative extension --- showing, at each codimension-\(1\) point \(Z\) separately, that the
  morphism \(\Spec k(V) \to A\) extends over \(\Spec \mathcal O_{V, Z}\) using only that \(A\) is
  proper and a group (so the difference-map argument can be run \emph{locally} at \(Z\) without
  globally building \(\mathrm{div}(f)\)) --- side-steps building divisor theory entirely.
  \textbf{This surface extension is \emph{not} on the genus-\(0\) critical path (corrected, iter-164).}
  An earlier draft routed the genus-\(0\) base case through this extension, by collapsing the
  additive-defect map \(\psi(x, y) = f(x + y) - f(x) - f(y)\) on the \emph{affine} surface
  \(\mathbb G_a \times \mathbb G_a\) and then extending \(\psi\) over the complete surface
  \(\mathbb P^1 \times \mathbb P^1\) --- which genuinely \emph{would} need Theorem~3.2's surface case
  (Milne Theorem~3.4 / Lemma~3.3). That detour is unnecessary. The 4th route
  (\cref{rmk:base_case_fourth_route}) replaces the two-variable affine defect map by the
  \emph{total} translation action \(\sigma \colon \mathbb P^1 \times \mathbb G_a \to \mathbb P^1\)
  (\cref{def:gaTranslationP1}): \(h := \sigma \fatsemi f\) is a total morphism on \(\mathbb P^1
  \times \mathbb G_a\) (no extension needed), and the complete-source factor is \(\mathbb P^1\)
  itself, so the additive decomposition \cref{lem:hom_additivity_over_product} applies directly.
  Consequently \cref{thm:rational_map_to_av_extends} (Theorem~3.2), the codim-\(1\) indeterminacy
  sub-build (Lemma~3.3), and the Weil-divisor gap discussed above are all confined to
  \textbf{Route~A} (the positive-genus Albanese universal property, Milne \S III.6); they are
  retained in this chapter solely as Route-A inputs and may be deferred until Route~A consumes
  them. The genus-\(0\) close no longer depends on any of them.
\end{remark}
\section{The genus-\(0\) base objects and the \(\mathbb G_m\)-scaling action (the primary route)}
\begin{definition}
[Concrete genus-\(0\) base objects: \(\mathbb P^1\), \(\mathbb G_a\), \(\mathbb G_m\) over \(\bar k\)]
  \label{def:genus0_base_objects}
  \lean{AlgebraicGeometry.ProjectiveLineBar}
  % Archon-original Lean encoding of standard objects; the mathematics is textbook (Hartshorne
  % I.2--3 for the projective line; the group laws on \(\mathbb A^1\) are elementary). No external
  % % SOURCE QUOTE is needed for the Lean encoding itself.
  Work over \(\Spec \bar k\) with \(\bar k\) algebraically closed. The primary \(\mathbb G_m\)-scaling
  route (and its \(\mathbb G_a\)-additive alternative) needs only three concrete objects of
  \(\mathrm{Over}\,(\Spec \bar k)\):
  \begin{itemize}
    \item \(\mathbb P^1 = \mathbb P^1_{\bar k}\), the projective line --- a smooth proper
      geometrically irreducible curve of genus \(0\) (Mathlib: projective space, or \(\mathrm{Proj}\)
      of the graded polynomial ring; \lean name \texttt{ProjectiveLineBar}). It carries
      the distinguished \(\bar k\)-points \(0\), \(1\), \(\infty\) and the two standard affine charts
      \(\mathbb A^1 = \mathbb P^1 \setminus \{\infty\}\) (coordinate \(x\)) and
      \(\mathbb P^1 \setminus \{0\}\) (coordinate \(u = 1/x\)).
    \item \(\mathbb G_a = \mathbb A^1_{\bar k}\), the affine line as the \emph{additive} group object
      (Mathlib: \texttt{AffineSpace}\,\(\mathbb A(1; \Spec \bar k)\) with the \(\mathtt{GrpObj}\)
      additive structure \((x, y) \mapsto x + y\), identity \(0\); \lean name \texttt{Ga}).
      As a scheme \(\mathbb G_a = \mathbb A^1 = \mathbb P^1 \setminus \{\infty\}\), and the open
      immersion \(\mathbb G_a \hookrightarrow \mathbb P^1\) has dense (cofinite) image.
    \item \(\mathbb G_m = \mathbb A^1 \setminus \{0\} = \mathbb P^1 \setminus \{0, \infty\}\), the
      \emph{multiplicative} group object (group law \((x, y) \mapsto xy\), identity \(1\); \lean name
      \texttt{Gm}). Its open immersion into \(\mathbb P^1\) also has dense image.
  \end{itemize}
  All three are of finite type over \(\bar k\); \(\mathbb P^1\) is proper, while \(\mathbb G_a\) and
  \(\mathbb G_m\) are affine (hence not complete). These are the only ``genus-\(0\)'' geometric inputs
  the base case needs.
\end{definition}
% Per-decl Lean coverage hooks for the components of \cref{def:genus0_base_objects}. The bundled
% definition above pins \texttt{ProjectiveLineBar} in its own \lean{...} block; the sibling
% definition blocks below pin the remaining Lean targets one-by-one so the blueprint-doctor can
% statically link each declaration.
\begin{definition}
[\(\mathbb G_a\) as an object of \(\mathrm{Over}\,(\Spec \bar k)\)]
  \label{def:ga}
  \lean{AlgebraicGeometry.Ga}
  \uses{def:genus0_base_objects}
  % Archon-original Lean encoding (per-decl pin for the additive-group bullet of
  % \cref{def:genus0_base_objects}); no external % SOURCE QUOTE needed.
  The additive group object \(\mathbb G_a = \mathbb A^1_{\bar k}\) described in
  \cref{def:genus0_base_objects}, packaged as an object of \(\mathrm{Over}\,(\Spec \bar k)\).
\end{definition}
% NOTE (iter-172 encoding clarification): the Lean encoding of \(\mathbb G_m\) realises
% \(\mathbb A^1 \setminus \{0\}\) as the AFFINE scheme \(\Spec(\mathrm{Localization.Away}\,t)\)
% rather than as the basic open subscheme \(D(t) \subset \mathbb A^1\). The two are
% canonically isomorphic; the affine-Spec choice is preferred in this project because
% downstream consumers (the scaling action \(\sigma_\times\) of
% \cref{def:gaTranslationP1}, the additivity decomposition input
% \(W = \mathbb G_m\) of \cref{lem:hom_additivity_over_product}) need an affine scheme,
% not just an open subscheme of one.
\begin{definition}
[\(\mathbb G_m\) as an object of \(\mathrm{Over}\,(\Spec \bar k)\)]
  \label{def:gm}
  \lean{AlgebraicGeometry.Gm}
  \uses{def:genus0_base_objects}
  % Archon-original Lean encoding (per-decl pin for the multiplicative-group bullet of
  % \cref{def:genus0_base_objects}); no external % SOURCE QUOTE needed.
  The multiplicative group object \(\mathbb G_m = \mathbb A^1 \setminus \{0\}\) described in
  \cref{def:genus0_base_objects}, packaged as an object of \(\mathrm{Over}\,(\Spec \bar k)\).
  In Lean this is encoded as the \emph{affine} scheme
  \(\Spec(\mathrm{Localization.Away}\,t)\) (the spectrum of the localisation
  \(\bar k[t, t^{-1}]\)), \emph{not} as the basic open subscheme \(D(t) \subset \mathbb A^1\);
  the two are canonically isomorphic but downstream consumers (the scaling action of
  \cref{def:gaTranslationP1} and the additivity-over-product input
  \(W = \mathbb G_m\) in \cref{lem:hom_additivity_over_product}) require an affine carrier.
\end{definition}
\begin{definition}
[The distinguished \(\bar k\)-point \(0 \in \mathbb P^1\)]
  \label{def:p1bar_zero}
  \lean{AlgebraicGeometry.ProjectiveLineBar.zeroPt}
  \uses{def:genus0_base_objects}
  % Archon-original Lean encoding (per-decl pin for the \(\bar k\)-point \(0 \in \mathbb P^1\) of
  % \cref{def:genus0_base_objects}); no external % SOURCE QUOTE needed.
  The \(\bar k\)-point \(0 \in \mathbb P^1(\bar k)\) realised as the morphism
  \(\Spec \bar k \to \mathbb P^1\) in \(\mathrm{Over}\,(\Spec \bar k)\), defined via
  \(\mathrm{Proj.fromOfGlobalSections}\) on the projective line of \cref{def:genus0_base_objects}.
  It is one of the two fixed points of the \(\mathbb G_m\)-scaling action used in the
  primary genus-\(0\) proof.
\end{definition}
\begin{definition}
[The distinguished \(\bar k\)-point \(1 \in \mathbb P^1\)]
  \label{def:p1bar_one}
  \lean{AlgebraicGeometry.ProjectiveLineBar.onePt}
  \uses{def:genus0_base_objects}
  % Archon-original Lean encoding (per-decl pin for the \(\bar k\)-point \(1 \in \mathbb P^1\) of
  % \cref{def:genus0_base_objects}); no external % SOURCE QUOTE needed.
  The \(\bar k\)-point \(1 \in \mathbb P^1(\bar k)\) realised as the morphism
  \(\Spec \bar k \to \mathbb P^1\) in \(\mathrm{Over}\,(\Spec \bar k)\), defined via
  \(\mathrm{Proj.fromOfGlobalSections}\) on the projective line of \cref{def:genus0_base_objects}.
\end{definition}
\begin{definition}
[The distinguished \(\bar k\)-point \(\infty \in \mathbb P^1\)]
  \label{def:p1bar_infty}
  \lean{AlgebraicGeometry.ProjectiveLineBar.inftyPt}
  \uses{def:genus0_base_objects}
  % Archon-original Lean encoding (per-decl pin for the \(\bar k\)-point \(\infty \in \mathbb P^1\) of
  % \cref{def:genus0_base_objects}); no external % SOURCE QUOTE needed.
  The \(\bar k\)-point \(\infty \in \mathbb P^1(\bar k)\) realised as the morphism
  \(\Spec \bar k \to \mathbb P^1\) in \(\mathrm{Over}\,(\Spec \bar k)\), defined via
  \(\mathrm{Proj.fromOfGlobalSections}\) on the projective line of \cref{def:genus0_base_objects}.
  It is the second fixed point of the \(\mathbb G_m\)-scaling action.
\end{definition}
\begin{definition}
[The multiplicative identity \(1 \in \mathbb G_m\)]
  \label{def:gm_one}
  \lean{AlgebraicGeometry.Gm.onePt}
  \uses{def:genus0_base_objects}
  % Archon-original Lean encoding (per-decl pin for the unit point of the multiplicative group
  % of \cref{def:genus0_base_objects}); no external % SOURCE QUOTE needed.
  The multiplicative identity \(1 \in \mathbb G_m(\bar k)\), supplied in Lean as the unit map
  \(\eta[\mathbb G_m] \colon \Spec \bar k \to \mathbb G_m\) in \(\mathrm{Over}\,(\Spec \bar k)\).
  It serves as the \(W\)-base point for the additive decomposition applied in the primary
  \(\mathbb G_m\)-scaling proof.
\end{definition}
\begin{lemma}
[\(\mathbb P^1_{\bar k}\) is proper over \(\Spec \bar k\)]
  \label{lem:projectiveLineBar_isProper}
  \lean{AlgebraicGeometry.projectiveLineBar_isProper}
  \uses{def:genus0_base_objects}
  % Archon-original instance; the mathematics is textbook (Hartshorne
  % II.4 + the standard properness of Proj of a finitely generated graded
  % algebra over a Noetherian base; here the base is a field, hence trivially
  % Noetherian). No external SOURCE QUOTE needed.
  The structure morphism \((\mathtt{ProjectiveLineBar}\,\bar k).\mathrm{hom}
  \colon \mathbb P^1_{\bar k} \to \Spec \bar k\) is proper.
  \medskip
  \noindent\textbf{Proof sketch.} Unfold \((\mathtt{ProjectiveLineBar}\,\bar k)
  .\mathrm{hom}\) as the composition
  \(\Proj.\mathrm{toSpecZero}\,(\mathrm{projectiveLineBarGrading}\,\bar k)
  \fatsemi \Spec.\mathrm{map}\,(\mathrm{algebraMap}\,\bar k\,
  \uparrow{\mathcal A_0})\). The first factor is proper by Mathlib's
  \(\mathtt{Proj.instIsProperToSpecZero}\): the polynomial algebra
  \(\bar k[X_0, X_1]\) is finite type over its degree-\(0\) piece
  \(\mathcal A_0\), and the \(\mathtt{IsScalarTower}\) +
  \(\mathtt{Algebra.FiniteType}\) chain (via
  \(\mathtt{Algebra.FiniteType.of\_restrictScalars\_finiteType}\)) supplies the
  finite-type witness. The second factor \(\Spec.\mathrm{map}\,(\mathrm{algebraMap}
  \,\bar k\,\uparrow{\mathcal A_0})\) is an isomorphism: the algebra map
  \(\bar k \to \uparrow{\mathcal A_0}\) is bijective because \(\mathcal A_0\)
  is the constants subalgebra of \(\bar k[X_0, X_1]\) (injectivity via
  \(\mathtt{MvPolynomial.C\_injective}\); surjectivity via
  \(\mathtt{MvPolynomial.homogeneousComponent\_of\_mem}\) +
  \(\mathtt{homogeneousComponent\_zero}\)). Composition of proper with iso is
  proper. LOC: \(\sim 40\) (already landed axiom-clean).
\end{lemma}
\begin{lemma}


\leanok
[\(\mathbb P^1_{\bar k}\) is smooth of relative dimension \(1\)]
  \label{lem:projectiveLineBar_smoothOfRelDim}
  \lean{AlgebraicGeometry.projectiveLineBar_smoothOfRelDim}
  \uses{def:genus0_base_objects, lem:mvPolynomialFin_isStandardSmoothOfRelativeDimension}
  % Archon-original instance; the mathematics is textbook (Hartshorne
  % II.8.20 for projective space smoothness over a field; the relative
  % dimension is the projective dimension n=1 for \(\mathbb P^1\)). No
  % external SOURCE QUOTE needed.
  The structure morphism \((\mathtt{ProjectiveLineBar}\,\bar k)
  .\mathrm{hom} \colon \mathbb P^1_{\bar k} \to \Spec \bar k\) is smooth
  of relative dimension \(1\). The proof uses the standard 2-chart
  affine cover by \(D_+(X_0)\) and \(D_+(X_1)\) (each chart is
  \(\Spec(\bar k[t])\) where \(t\) is the affine coordinate
  \(X_i/X_{1-i}\)); each chart is standard smooth via the polynomial
  ring presentation, hence smooth over \(\Spec \bar k\); the relative
  dimension equals \(1\) via the single generator \(t\) of the
  polynomial ring. The smoothness on each chart lifts to smoothness on
  \(\mathbb P^1_{\bar k}\) via the open-cover Zariski-locality criterion
  for \(\mathtt{SmoothOfRelativeDimension}\), namely Mathlib's
  \(\mathtt{IsZariskiLocalAtSource.of\_openCover}\) (specialised to
  \(\mathtt{SmoothOfRelativeDimension}\,1\) via the
  \(\mathtt{HasRingHomProperty} \Rightarrow \mathtt{IsZariskiLocalAtSource}\)
  chain).
  \medskip
  \noindent\textbf{Per-chart substrate dependency.} The per-chart closure
  decomposes as: (i) reduce smoothness to a
  \(\mathtt{RingHom.Locally\,IsStandardSmoothOfRelativeDimension\,1}\)
  statement via Mathlib's \(\mathtt{HasRingHomProperty.iff\_of\_isAffine}\)
  (both source and target are affine); (ii) reduce \(\mathtt{Locally}\) to the
  direct property via
  \(\mathtt{RingHom.locally\_of\,isStandardSmoothOfRelativeDimension\_respectsIso}\);
  (iii) close on the composite \(\bar k \to \Gamma(\mathrm{chart}_i, \top)\)
  ring map by combining Mathlib's
  \(\mathtt{Algebra.IsStandardSmoothOfRelativeDimension.of\_algEquiv}\) with
  the project's \cref{lem:mvPolynomialFin_isStandardSmoothOfRelativeDimension}
  (the \(\mathrm{MvPolynomial\,(Fin\,n)\,R}\) instance) transported along the
  chart-ring iso \cref{def:proj_chart_ring_iso}
  (\(\mathtt{homogeneousLocalizationAwayIso}\)) reparameterised as an
  \(\bar k\)-algebra equivalence \(\mathrm{MvPolynomial}\,(\mathrm{Fin}\,1)\,
  \bar k \cong \Gamma(\mathrm{chart}_i, \top)\).
  \medskip
  \noindent\textbf{Status (iter-196).} The cover-reduction structural step is
  landed axiom-clean. The per-chart aux (Lean:
  \(\mathtt{projectiveLineBar\_smooth\_chart\_aux}\), private) carries the
  residual sorry, because the chart-ring iso
  \(\mathtt{homogeneousLocalizationAwayIso}\) lives in
  \(\mathtt{AlgebraicJacobian/Genus0BaseObjects/ChartIso.lean}\), which is
  \emph{downstream} of \(\mathtt{BareScheme.lean}\) — so the per-chart closure
  cannot consume it without an import cycle. \textbf{Recommended refactor:}
  relocate \(\mathtt{projectiveLineBar\_smoothOfRelDim}\) (and its private
  per-chart aux) to \(\mathtt{ChartIso.lean}\) or a new
  \(\mathtt{Genus0BaseObjects/Smooth.lean}\) where the chart-ring iso is in
  scope; then close the per-chart gap via
  \(\mathtt{Algebra.IsStandardSmoothOfRelativeDimension.of\_algEquiv}\) in
  \(\sim 10\) LOC. This relocation is being dispatched separately to the
  refactor subagent (slug \(\mathtt{barescheme\text{-}smoothness\text{-}relocation}\))
  this iter; once landed, the residual sorry closes.
\end{lemma}
\begin{lemma}
[Project-side substrate: \(\mathrm{MvPolynomial}\,(\mathrm{Fin}\,n)\,R\) is
\(R\)-standard smooth of relative dimension \(n\)]
  \label{lem:mvPolynomialFin_isStandardSmoothOfRelativeDimension}
  \lean{AlgebraicGeometry.mvPolynomialFin_isStandardSmoothOfRelativeDimension}
  \uses{def:mvpoly_submersive}
  % Project-local Mathlib supplement (axiom-clean, iter-196 landed). Builds the
  % canonical SubmersivePresentation of MvPolynomial as an R-algebra (no
  % relations; generators indexed by Fin n). No external textbook source: the
  % construction is pure category-of-presentations infrastructure on the
  % Mathlib `Algebra.SubmersivePresentation` machinery.
  For any commutative ring \(R\) and natural number \(n\), the polynomial ring
  \(\mathrm{MvPolynomial}\,(\mathrm{Fin}\,n)\,R\) is \(R\)-standard smooth of
  relative dimension \(n\). At the pinned Mathlib commit \(\mathtt{b80f227}\)
  Mathlib does \emph{not} ship a direct instance
  \(\mathtt{Algebra.IsStandardSmoothOfRelativeDimension}\,n\,R\,
  (\mathrm{MvPolynomial}\,(\mathrm{Fin}\,n)\,R)\) (per iter-182 prover scout);
  the project supplies it via a 5-declaration
  \(\mathtt{Algebra.SubmersivePresentation}\) chain (Lean names with
  \(\mathtt{Algebra}\)-namespace target):
  \begin{enumerate}
    \item \(\mathtt{mvPolyGenerators}\) — canonical generators
      (\(\mathtt{val} := \mathrm{MvPolynomial.X}\); section \(\sigma' := \mathrm{id}\)).
    \item \(\mathtt{mvPolyPresentation}\) — presentation with empty relation set
      (\(\sigma = \mathtt{PEmpty}\); kernel of \(\mathrm{aeval}\,\mathrm{X}\) is \(\bot\)).
    \item \(\mathtt{mvPolyPreSubmersivePresentation}\) — pre-submersive with empty
      \(\mathtt{map} \colon \mathtt{PEmpty} \to \alpha\).
    \item \(\mathtt{mvPolySubmersivePresentation}\) — submersive presentation;
      Jacobian determinant of the zero-sized matrix is \(1\) (a unit).
    \item \(\mathtt{mvPolynomialFin\_isStandardSmoothOfRelativeDimension}\) —
      the final instance, obtained from
      \(\mathtt{Algebra.SubmersivePresentation.isStandardSmoothOfRelativeDimension}\)
      with dimension count \(\mathrm{Nat.card}\,(\mathrm{Fin}\,n)
      - \mathrm{Nat.card}\,\mathtt{PEmpty} = n\).
  \end{enumerate}
  Total LOC across the chain: \(\sim 50\) (all axiom-clean, iter-196 landed).
  This substrate is the load-bearing input to the per-chart smoothness step of
  \cref{lem:projectiveLineBar_smoothOfRelDim}.
\end{lemma}
\begin{lemma}
[\(\mathbb P^1_{\bar k}\) is geometrically irreducible]
  \label{lem:projectiveLineBar_geomIrred}
  \lean{AlgebraicGeometry.projectiveLineBar_geomIrred}
  \uses{def:genus0_base_objects}
  % Archon-original instance; the mathematics is textbook (Hartshorne
  % II.4 + I.1; \(\Proj\) of an integral domain over a field is
  % integral and geometrically integral when the field is algebraically
  % closed and the polynomial ring is well-defined). No external SOURCE
  % QUOTE needed.
  The structure morphism \((\mathtt{ProjectiveLineBar}\,\bar k)
  .\mathrm{hom} \colon \mathbb P^1_{\bar k} \to \Spec \bar k\) is
  geometrically irreducible: for every field extension \(K/\bar k\),
  the base change \(\mathbb P^1_{\bar k} \times_{\Spec \bar k} \Spec K
  = \mathbb P^1_K\) is irreducible. Since \(\bar k\) is algebraically
  closed and \(\bar k[X_0, X_1]\) is the polynomial ring (integral
  domain), \(\mathbb P^1_{\bar k} = \Proj(\bar k[X_0, X_1])\) is
  integral, and the base change to any field extension \(K\) preserves
  integrality (the homogeneous coordinate ring \(K[X_0, X_1]\) is
  itself an integral domain, hence \(\mathbb P^1_K\) is integral and
  irreducible).
\end{lemma}
\begin{proof}
  \uses{def:genus0_base_objects}
  % NOTE (iter-197 blueprint-writer): the iter-196 lean-vs-blueprint-checker
  % flagged this proof as under-specified for Lean closure. The expanded sketch
  % below decomposes the argument into five named sub-helpers (A--E), pinning
  % the load-bearing Proj-base-change iso (Helper A) to a future Lean target;
  % until Helper A lands, the public `projectiveLineBar_geomIrred` instance
  % remains a Tier-3 named scaffold sorry.
  We decompose the argument into five sub-helpers. The load-bearing piece is
  Helper~A (the Proj base-change isomorphism), which is currently a Mathlib gap
  at the pinned commit \(\mathtt{b80f227}\) and must be supplied as a project-side
  substrate (estimated 100--200 LOC). The remaining helpers are then either
  Mathlib-direct or short combinatorial bridges. Total estimated closure cost
  across the chain: 200--350 LOC over multiple iters.
  \medskip
  \noindent\textbf{Sub-helper~A (load-bearing): base-change of \(\Proj\).}
  \(\lean{AlgebraicGeometry.Proj.baseChangeIso}\) (future target;
  Mathlib upstreaming candidate or project supplement). For a graded ring
  \(\mathcal A \colon \mathbb N \to \sigma\) over a base ring \(R\) and an
  \(R\)-algebra \(K\), there is a natural isomorphism of schemes over \(\Spec R\):
  \[ \Proj\,(K \otimes_R \mathcal A)
     \ \cong\ \Proj\,\mathcal A \,\times_{\Spec R}\, \Spec K. \]
  This is the standard Stacks Project base-change formula for relative
  \(\Proj\) (Stacks tag 0BLW, ``base change of relative Proj''). At the pinned
  Mathlib commit \(\mathtt{b80f227}\) no version of this isomorphism is shipped
  for the abstract graded-ring \(\Proj\) construction; the project-side build
  needs to construct the iso functorially, prove naturality in \(K\), and
  expose the \(\Spec K\)-projection compatibility for downstream consumers.
  \medskip
  \noindent\textbf{Sub-helper~B: chart-iso reparameterisation over an extension
  field.} \(\lean{AlgebraicGeometry.homogeneousLocalizationAwayIso.baseChange}\)
  (future target). The chart-ring iso of \cref{def:proj_chart_ring_iso}
  identifies \(\mathrm{Away}\,\mathcal A\,X_i \cong \bar k[u]\) over the base
  field \(\bar k\). Helper~B re-runs the same construction over an arbitrary
  field extension \(K / \bar k\) to obtain
  \(K \otimes_{\bar k} \mathrm{Away}\,\mathcal A\,X_i \cong K[u]\) (the affine
  chart of \(\mathbb P^1_K\)). This is a routine base-change of the iso of
  \cref{def:proj_chart_ring_iso} and carries no new substrate gap once
  Helper~A is in scope.
  \medskip
  \noindent\textbf{Sub-helper~C: the base-changed graded ring is a graded
  integral domain.} For \(K\) any field extension of \(\bar k\), the tensor
  product \(K \otimes_{\bar k} \bar k[X_0, X_1]\) is a graded \(K\)-algebra
  canonically identified with the polynomial ring \(K[X_0, X_1]\), which is
  itself an integral domain. Consequently the underlying scheme of
  \(\Proj\,(K \otimes_{\bar k} \bar k[X_0, X_1])\) is integral on every basic
  open \(D_+(X_i)\) (each chart ring is a localisation of the integral domain
  \(K[X_0, X_1]\), hence a domain), and a pointwise check at generic points
  gives integrality of the whole scheme. (Currently un-pinned; iter-197+ prover
  names the Lean target as it materialises the helper.)
  \medskip
  \noindent\textbf{Sub-helper~D: \(\Proj\) of an integral graded domain over a
  field is integral.} Mathlib's
  \(\mathtt{Proj.instIsIntegral}\)-style instance (when the irrelevant-ideal
  cofinality assumption holds, as it does for the standard polynomial grading)
  packages the global integrality of \(\Proj\,(K[X_0, X_1])\). Integrality
  implies irreducibility of the underlying scheme. (Currently un-pinned;
  iter-197+ prover names the Lean target as it materialises the helper.)
  \medskip
  \noindent\textbf{Sub-helper~E: assemble via Helper~A.} For every field
  extension \(K / \bar k\), Helper~A gives
  \((\mathtt{ProjectiveLineBar}\,\bar k) \times_{\Spec \bar k} \Spec K
  \cong \Proj\,(K \otimes_{\bar k} \bar k[X_0, X_1]) = \mathbb P^1_K\); Helper~C
  + Helper~D give that \(\mathbb P^1_K\) is integral, hence irreducible.
  Combined with the algebraic-closure hypothesis on \(\bar k\) (which makes the
  generic-point check uniform across all extensions), this gives
  \(\mathtt{GeometricallyIrreducible}\,(\mathtt{ProjectiveLineBar}\,\bar k).\mathrm{hom}\).
  (Currently un-pinned; iter-197+ prover names the Lean target as it
  materialises the helper.)
  \medskip
  \noindent\textbf{Mathlib references and gap budget.} Stacks tag 0BLW (Proj
  base-change) is \emph{not} in Mathlib \(\mathtt{b80f227}\); project-side
  substrate is required and is the dominant cost (Helper~A,
  \(\sim 100\)--\(200\) LOC). Mathlib does ship
  \(\mathtt{Proj.instIsIntegral}\) (Helper~D) and the
  \(\mathrm{Algebra.TensorProduct}\) machinery (Helper~C) but \emph{not} the
  graded-tensor-product-of-polynomials integral-domain instance directly, so
  Helper~C requires a brief project-side bridge. Total estimated closure cost:
  200--350 LOC across the five helpers; Helper~A dominates. This sorry is
  acceptable as a Tier-3 named scaffold until the Proj base-change
  infrastructure lands.
\end{proof}
\begin{definition}
[The multiplicative \(\mathtt{GrpObj}\) structure on \(\mathbb G_m\)]
  \label{def:gm_grpObj}
  \lean{AlgebraicGeometry.gm_grpObj}
  \uses{def:gm}
  % Archon-original Lean encoding (per-decl pin for the multiplicative-group structure on
  % \cref{def:gm}); no external % SOURCE QUOTE needed.
  The multiplicative group-object structure \((\mathbb G_m, \cdot, 1)\) in
  \(\mathrm{Over}\,(\Spec \bar k)\), recording the multiplication \((x, y) \mapsto xy\) and unit \(1\)
  of \cref{def:gm} as a Lean \(\mathtt{GrpObj}\) instance. This structure is consumed only as a
  carrier for the scaling action \(\sigma_\times\) of \cref{def:gaTranslationP1}; the primary
  \(\mathbb G_m\)-scaling proof of \cref{prop:morphism_P1_to_AV_constant} uses \(\mathbb G_m\) as a
  parameter scheme (via \(W\) in the additive decomposition), not its group-object axioms.
\end{definition}
\subsection*{Chart cover and chart-ring iso (formalisation infrastructure)}
The next three declarations are iter-168 formalisation infrastructure that the chartwise
construction of the \(\mathbb G_m\)-scaling action \(\sigma_\times\) (\cref{def:gaTranslationP1})
and its downstream consumers (\cref{prop:morphism_P1_to_AV_constant}) ultimately rest on.
They realise, at the Lean level, the standard ``affine chart of \(\mathbb P^1\) is \(\mathbb
A^1\)'' picture: a two-chart affine open cover of \(\overline{\mathbb P^1}\) by \(D_+(X_0)\) and
\(D_+(X_1)\); the identification of each chart ring with the one-variable polynomial ring
\(\bar k[u]\); and the resulting reducedness of \(\overline{\mathbb P^1}\) as a scheme. None
of them is on the genus-\(0\) critical path conceptually, but each is consumed by the Lean
formalisation downstream.
\begin{definition}
[Two-chart affine open cover of \(\overline{\mathbb P^1}\)]
  \label{def:projlinebar_affine_cover}
  \lean{AlgebraicGeometry.projectiveLineBarAffineCover}
  \uses{def:genus0_base_objects}
  % Archon-original Lean encoding of the textbook fact that \(\mathbb P^1\) is covered by its
  % two standard affine charts \(D_+(X_0)\) and \(D_+(X_1)\); the construction is built on top of
  % Mathlib's \(\mathrm{Proj.affineOpenCoverOfIrrelevantLESpan}\). No external % SOURCE QUOTE
  % needed for the Lean encoding itself.
  The two-chart affine open cover of \(\overline{\mathbb P^1}\) by the standard \(D_+(X_0)\) and
  \(D_+(X_1)\), built by specialising Mathlib's
  \(\mathrm{Proj.affineOpenCoverOfIrrelevantLESpan}\) to the index set
  \(\iota = \{0, 1\} = \mathrm{Fin}\,2\) with family \(f = (X_0, X_1)\) and weights
  \(m = (1, 1)\). The non-trivial input is that the irrelevant ideal of the bi-graded
  \(\bar k[X_0, X_1]\) is contained in \(\mathrm{Ideal.span}\{X_0, X_1\}\), which holds by
  writing any irrelevant element as a sum of monomials of strictly positive total degree
  (so each monomial is divisible by \(X_0\) or \(X_1\)). This cover packages the two affine
  charts into the data that Mathlib's \(\mathrm{IsReduced.of\_openCover}\) and
  \(\mathrm{Scheme.Cover.glueMorphisms}\) APIs consume, and is the chartwise glue for the
  iter-169 attack on the body of \(\sigma_\times = \mathtt{gmScalingP1}\)
  (\cref{def:gaTranslationP1}).
\end{definition}
\begin{definition}

\leanok
[Chart-ring iso: \(\mathrm{Away}\,\mathcal A\,X_i \cong \bar k[u]\)]
  \label{def:proj_chart_ring_iso}
  \lean{AlgebraicGeometry.homogeneousLocalizationAwayIso}
  % Archon-original Lean encoding of the textbook identification ``the affine chart of
  % \(\mathbb P^1\) at \(X_i\) is \(\mathbb A^1\) with coordinate \(X_{1-i}/X_i\)'', expressed at
  % the homogeneous-localization level. No external % SOURCE QUOTE needed.
  For each \(i \in \{0, 1\}\), the chart-\(i\) ring of \(\overline{\mathbb P^1}\) --- the
  degree-zero part \(\mathrm{Away}\,\mathcal A\,X_i\) of the localization of the standard
  \(\mathbb N\)-graded polynomial ring \(\bar k[X_0, X_1]\) at \(X_i\) --- is isomorphic as a
  \(\bar k\)-algebra to the one-variable polynomial ring \(\bar k[u]\), where the affine
  chart coordinate is \(u = X_{1-i}/X_i\). This is the formalisation-level realisation of
  the standard ``affine chart of \(\mathbb P^1\) is \(\mathbb A^1\)'' identification expressed
  at the homogeneous-localization level. The forward direction \(\mathrm{Away}\,\mathcal
  A\,X_i \to \bar k[u]\) is built via \(\mathrm{Localization.awayLift}\) from the chart
  evaluation \(X_i \mapsto 1\), \(X_{1-i} \mapsto u\); the inverse direction is built via the
  universal property of \(\mathrm{MvPolynomial}\,\{*\}\,\bar k\), sending \(u\) to the
  homogeneous-localization element \(X_{1-i}/X_i \in \mathrm{Away}\,\mathcal A\,X_i\); the
  two directions are packaged as a \(\mathrm{RingEquiv}\) via \(\mathrm{RingEquiv.ofRingHom}\)
  applied to the two round-trip identities.
  % NOTE (iter-172 refresh of iter-171): the reverse round-trip body
  % \cref{lem:proj_chart_ring_iso_aux_left} is a REAL ``cancel surjective'' proof
  % (axiom-clean once the surjectivity helper is closed), and the residual
  % \texttt{sorryAx} taint of the iso now flows entirely through the named
  % helper \cref{lem:mvPoly_to_homogeneousLocalization_away_surjective}
  % (Lean: \texttt{AlgebraicGeometry.mvPolyToHomogeneousLocalizationAway\_surjective},
  % \texttt{Genus0BaseObjects.lean:372}). The forward round-trip
  % (the identity on \(\bar k[u]\)) remains axiom-clean. Closing the surjectivity
  % helper (via \texttt{Away.adjoin\_mk\_prod\_pow\_eq\_top} specialised to \(d = 1\),
  % \(v = (X_0, X_1)\), \(dv = (1, 1)\)) discharges the chart-ring iso's last sorry.
\end{definition}
\begin{lemma}

\leanok
[Reverse round-trip of the chart-ring iso]
  \label{lem:proj_chart_ring_iso_aux_left}
  \lean{AlgebraicGeometry.homogeneousLocalizationAwayIso_aux_left}
  \uses{def:proj_chart_ring_iso, lem:mvPoly_to_homogeneousLocalization_away_surjective}
  % Archon-original Lean encoding; this is the residual identity needed for the chart-ring
  % map of \cref{def:proj_chart_ring_iso} to upgrade from a pair of mutually inverse-looking
  % ring maps to a genuine \(\mathrm{RingEquiv}\). No external % SOURCE QUOTE.
  The composite \(\mathrm{inverse} \circ \mathrm{forward} \colon \mathrm{Away}\,\mathcal
  A\,X_i \to \mathrm{Away}\,\mathcal A\,X_i\) is the identity ring map. This is the reverse
  half of the round-trip witnessing that the chart-ring map of
  \cref{def:proj_chart_ring_iso} is a ring isomorphism; the forward half \(\mathrm{forward}
  \circ \mathrm{inverse} = \mathrm{id}_{\bar k[u]}\) is already proven axiom-clean.
  As of iter-171, this reverse round-trip body has \emph{landed} as a real
  ``cancel surjective'' proof: from surjectivity of the inverse
  \(\bar k[u] \to \mathrm{Away}\,\mathcal A\,X_i\)
  (\cref{lem:mvPoly_to_homogeneousLocalization_away_surjective}) and the already-proven
  forward round-trip \(\mathrm{forward} \circ \mathrm{inverse} = \mathrm{id}_{\bar k[u]}\)
  one cancels to conclude \(\mathrm{inverse} \circ \mathrm{forward} = \mathrm{id}\) on the
  \(\mathrm{Away}\) side. The cancellation step itself is mechanical; the only remaining
  substantive content is the surjectivity sub-lemma below.
\end{lemma}
\begin{lemma}

\leanok
[Surjectivity of the inverse chart-ring map \(\bar k[u] \to \mathrm{Away}\,\mathcal A\,X_i\)]
  \label{lem:mvPoly_to_homogeneousLocalization_away_surjective}
  \lean{AlgebraicGeometry.mvPolyToHomogeneousLocalizationAway_surjective}
  % Archon-original Lean encoding of an elementary corollary of a Mathlib lemma. The
  % external authority is the Mathlib statement
  % \texttt{HomogeneousLocalization.Away.adjoin\_mk\_prod\_pow\_eq\_top}
  % (\texttt{Mathlib/RingTheory/GradedAlgebra/HomogeneousLocalization.lean:1064}); we do
  % not include a verbatim \% SOURCE QUOTE because no Mathlib source quote is currently
  % retained under \texttt{references/} for this lemma (Mathlib source files are not
  % mirrored in the project's reference cache). The mathematics is standard.
  For each \(i \in \{0, 1\}\), the ring map
  \[ \mathrm{mvPolyToHomogeneousLocalizationAway}_i \colon
     \mathrm{MvPolynomial}\,\{*\}\,\bar k \longrightarrow
     \mathrm{Away}\,\mathcal A\,X_i \]
  --- the inverse direction of the chart-ring iso of \cref{def:proj_chart_ring_iso},
  built via \(\mathrm{MvPolynomial.eval}_2\mathrm{Hom}\) as the ring map sending the
  generator \(u \in \mathrm{MvPolynomial}\,\{*\}\,\bar k\) to the homogeneous-localisation
  element \(X_{1-i}/X_i \in \mathrm{Away}\,\mathcal A\,X_i\) and the constants
  \(r \in \bar k\) along the composite \(\bar k \to \mathcal A_0 \to \mathrm{Away}\,\mathcal
  A\,X_i\) --- is surjective onto \(\mathrm{Away}\,\mathcal A\,X_i\).
  \medskip
  \textit{Proof.} The image of \(\mathrm{mvPolyToHomogeneousLocalizationAway}_i\)
  is the \(\mathcal A_0\)-subalgebra of \(\mathrm{Away}\,\mathcal A\,X_i\) generated by the
  single localised element \(X_{1-i}/X_i\), since \(\mathrm{MvPolynomial}\,\{*\}\,\bar k\)
  is generated as a \(\bar k\)-algebra by its single generator \(u\) and the base-ring map
  \(\bar k \to \mathcal A_0\) is bijective (via the canonical identification of degree-zero
  scalars in the \(\bar k\)-algebra \(\bar k[X_0, X_1]\) with \(\bar k\)). Mathlib's lemma
  \(\mathrm{HomogeneousLocalization.Away.adjoin\_mk\_prod\_pow\_eq\_top}\) says that for any
  multi-index \(v \colon \iota' \to A\) of homogeneous elements of degree \(dv\) generating
  \(\mathcal A\) up to degree \(d\), the \(\mathcal A_0\)-subalgebra of
  \(\mathrm{Away}\,\mathcal A\,(\prod v_i^{dv_i})\) generated by the family of ratios
  \(\{v_i^{dv_i}/\prod v_j^{dv_j}\}\) is the whole ring \(\top\). Specialising to
  \(d = 1\), \(\iota' = \mathrm{Fin}\,2\), \(v = (X_0, X_1)\), and exponent vector
  \(dv = (1, 1)\), the relevant generator \(\prod v_i^{dv_i} = X_0 \cdot X_1\) collapses
  along the localisation at \(X_i\) (the \(X_i\)-component is already a unit, so only the
  \(X_{1-i}/X_i\) ratio remains as a non-trivial generator), and the conclusion is that
  \(\mathrm{Algebra.adjoin}\,\mathcal A_0\,\{X_{1-i}/X_i\} = \top\) inside
  \(\mathrm{Away}\,\mathcal A\,X_i\). This is exactly the image of
  \(\mathrm{mvPolyToHomogeneousLocalizationAway}_i\), so the map is surjective. \qed
\end{lemma}
\begin{lemma}

\leanok
[Chart-ring iso preserves \(\mathrm{algebraMap}\,\bar k\)]
  \label{lem:chart_ring_iso_preserves_algebraMap}
  \lean{AlgebraicGeometry.homogeneousLocalizationAwayIso_algebraMap}
  % Archon-original Lean encoding; the \(\bar k\)-algebra-preservation half of the chart-ring iso
  % (\cref{def:proj_chart_ring_iso}). The cited Mathlib fact
  % \texttt{HomogeneousLocalization.algebraMap\_eq} is a structural unfolding of how the
  % \(\bar k\)-algebra structure on \(\mathrm{Away}\,\mathcal A\,X_i\) is wired through
  % \texttt{Algebra.compHom} from \(\mathtt{HomogeneousLocalization.fromZeroRingHom}\), not a
  % statement copied from a textbook; no external % SOURCE QUOTE is retained because the
  % project does not mirror Mathlib source files under \texttt{references/}.
  For each \(i \in \{0, 1\} = \mathrm{Fin}\,2\), the forward direction of the chart-ring iso
  of \cref{def:proj_chart_ring_iso} preserves the canonical \(\bar k\)-algebra structure:
  composing
  \[ \mathrm{forward}_i \colon
     \mathrm{Away}\,\mathcal A\,X_i
     \longrightarrow
     \mathrm{MvPolynomial}\,\{*\}\,\bar k \]
  with the \(\bar k\)-algebra map \(\mathrm{algebraMap}\,\bar k\,(\mathrm{Away}\,\mathcal A\,X_i)\)
  on the source side yields the canonical
  \(\mathrm{algebraMap}\,\bar k\,(\mathrm{MvPolynomial}\,\{*\}\,\bar k)\) on the target side:
  \[ \mathrm{algebraMap}\,\bar k\,(\mathrm{Away}\,\mathcal A\,X_i)
     \,\fatsemi\,
     \mathrm{forward}_i
     \ =\
     \mathrm{algebraMap}\,\bar k\,(\mathrm{MvPolynomial}\,\{*\}\,\bar k). \]
  Equivalently, the chart-ring iso sends each constant element \(r \in \bar k\) --- viewed as
  a degree-\(0\) element of \(\mathrm{Away}\,\mathcal A\,X_i\) via
  \(\mathrm{algebraMap}\,\bar k\,(\mathrm{Away}\,\mathcal A\,X_i)\) --- to the corresponding
  constant \(\mathrm{MvPolynomial.C}\,r \in \mathrm{MvPolynomial}\,\{*\}\,\bar k\).
  \medskip
  \noindent\textbf{Proof sketch.} The \(\bar k\)-algebra map
  \(\mathrm{algebraMap}\,\bar k\,(\mathrm{Away}\,\mathcal A\,X_i)\) unfolds (via Mathlib's
  \(\mathtt{HomogeneousLocalization.algebraMap\_eq}\)) as the
  \(\mathrm{Algebra.compHom}\) of \(\mathtt{HomogeneousLocalization.fromZeroRingHom}\) along the
  canonical \(\bar k \to \mathcal A_0\) inclusion of constants into the degree-\(0\) piece of
  the standard \(\mathbb N\)-graded ring \(\bar k[X_0, X_1]\). The forward direction
  \(\mathrm{forward}_i\) of \cref{def:proj_chart_ring_iso} sends a degree-\(0\) element of
  \(\mathrm{Away}\,\mathcal A\,X_i\) to its image under the chart evaluation \(X_i \mapsto 1\),
  \(X_{1-i} \mapsto u\); on a constant \(r \in \bar k\) this evaluation is exactly
  \(\mathrm{MvPolynomial.C}\,r\). The identification of the two ring maps on every constant
  \(r\) is therefore the round-trip identity
  \(\mathrm{forward}_i \circ \mathrm{inverse}_i = \mathrm{id}_{\bar k[u]}\) (the axiom-clean
  half of \cref{lem:proj_chart_ring_iso_aux_left}) applied to the constant
  \(\mathrm{MvPolynomial.C}\,r\), together with the surjectivity of
  \(\mathrm{inverse}_i = \mathrm{mvPolyToHomogeneousLocalizationAway}_i\)
  (\cref{lem:mvPoly_to_homogeneousLocalization_away_surjective}) that lifts the
  pointwise check to a ring-map equality. This is the load-bearing \(\bar k\)-algebra
  preservation step consumed by Step (B) of \cref{lem:gmscaling_chart_PLB_eq}, where the
  chart-\(i\) ring map is post-composed with \(\mathrm{algebraMap}\,\bar k\) of the relative
  tensor product to identify it with
  \(\mathrm{algebraMap}\,\bar k\,(\mathrm{Away}\,\mathcal A\,X_i \otimes_{\bar k}
  \mathtt{GmRing}\,\bar k)\).
\end{lemma}
\begin{lemma}

\leanok
[Chart-\(1\) \(\mathrm{Proj.appIso}\) evaluation: \(\mathtt{isLocElem} \mapsto [X_0/X_1] \mapsto 1\)]
  \label{lem:iotaGm_chart1_appIso_eval}
  \lean{AlgebraicGeometry.iotaGm_chart1_appIso_eval}
  \uses{def:proj_chart_ring_iso, lem:chart_ring_iso_preserves_algebraMap, lem:iotaGm_r_1_eq_specMap}
  % Archon-original Lean encoding of the chart-\(1\) \(\mathrm{Proj.appIso}\) evaluation step.
  % Extracted in iter-192 (per progress-critic STUCK verdict and blueprint-reviewer MF-1
  % finding: 4 consecutive iters - iter-188/189/190/191 - hit the LOC budget wall on this
  % evaluation inside the body close of \texttt{iotaGm\_chart1\_composition\_isOpenImmersion}).
  % Pulled out as a named standalone hook so the prover can target it directly. The
  % mathematical content is the standard fact that \(\mathrm{Proj.appIso}\) of the chart-\(i\)
  % homogeneous chart map identifies the canonical local section \(\mathtt{isLocElem}\) of
  % \(\Spec(\mathrm{Away}\,\mathcal A\,X_i)\) with the homogeneous-localisation ratio
  % \([X_{1-i}/X_i] \in \mathrm{HomogeneousLocalization.Away}\,\mathcal A\,X_i\). No
  % external % SOURCE QUOTE: this is a Mathlib-side identification + project-side composition.
  Let \(\mathcal A = \bar k[X_0, X_1]\) be the standard \(\mathbb N\)-graded polynomial ring,
  and consider the homogeneous chart morphism at the degree-\(1\) generator \(X_1\),
  \[ \mathrm{Proj.away}\iota\,X_1 \colon \Spec\,(\mathrm{Away}\,\mathcal A\,X_1)
     \longrightarrow \Proj\,\mathcal A, \]
  realising the basic open \(D_+(X_1) \subseteq \overline{\mathbb P^1} = \Proj\,\mathcal A\) as
  an affine open. Mathlib's \(\mathtt{Proj.appIso}\) (applied to the open immersion
  \(\mathrm{Proj.away}\iota\,X_1\)) packages, at the level of global sections over \(\top\),
  the canonical scheme-side identification
  \[ \Gamma(\Spec(\mathrm{Away}\,\mathcal A\,X_1),\, \top)
     \ =\ \mathrm{Away}\,\mathcal A\,X_1
     \ \xleftarrow{\ \sim\ }\
     \Gamma(\Proj\,\mathcal A,\, D_+(X_1))
     \ =\ \mathrm{HomogeneousLocalization.Away}\,\mathcal A\,X_1. \]
  Under this identification, the canonical local-section element
  \(\mathtt{isLocElem} \in \Gamma(\Spec(\mathrm{Away}\,\mathcal A\,X_1), \top)\)
  (the global section of \(\Spec\) corresponding under the affine Spec / global-sections
  adjunction to the identity ring map) is the image, via
  \((\mathtt{Proj.appIso}\,(\mathrm{Proj.away}\iota\,X_1)).\mathrm{inv}\), of the homogeneous
  fraction
  \[ [X_0/X_1] \ \in\ \mathrm{HomogeneousLocalization.Away}\,\mathcal A\,X_1 \]
  --- the ratio of the two degree-\(1\) generators, i.e.\ the chart-\(1\) affine coordinate
  \(u = X_0/X_1\). Composing further with the chart-\(1\) ring identification
  \(\mathrm{forward}_1 \colon \mathrm{Away}\,\mathcal A\,X_1 \xrightarrow{\sim}
  \mathrm{MvPolynomial}\,\{*\}\,\bar k\) of \cref{def:proj_chart_ring_iso} sends
  \([X_0/X_1]\) to the polynomial generator \(u\); then pre-composing with the source-side
  \(\mathtt{pullbackSpecIso}\) bridge of \cref{def:gmscaling_chart} (with the trivial
  \(\mathbb G_m\) factor at the identity section \(\lambda = 1\) collapsed by
  \(\mathtt{Algebra.TensorProduct.lid}\)) lands at \(1\) in the codomain ring. The combined
  composite ring map therefore sends \(\mathtt{isLocElem}\) to \(1 \in \bar k\):
  \[ (\mathtt{iotaGm\_chart1\_composition}).\mathrm{appTop}\,(\mathtt{isLocElem})
     \ =\ 1 \quad \text{in the codomain ring.} \]
\end{lemma}
\begin{proof}
  The Mathlib identification of \(\mathrm{Proj.away}\iota\,X_1\) as the open immersion
  presenting the basic open \(D_+(X_1) \subseteq \Proj\,\mathcal A\) packages
  \(\mathtt{Proj.appIso}\,(\mathrm{Proj.away}\iota\,X_1)\) as the canonical iso on global
  sections; its concrete value \(\mathtt{Proj.appIso\_apply}\) traces the canonical local
  section back to the homogeneous-localisation ratio of the two distinguished generators. The
  canonical local-section \(\mathtt{isLocElem}\) is, by the affine Spec / global-sections
  adjunction, the global section of \(\Spec(\mathrm{Away}\,\mathcal A\,X_1)\) corresponding to
  the identity ring map; under the inverse identification it becomes the element
  \([X_0/X_1] \in \mathrm{HomogeneousLocalization.Away}\,\mathcal A\,X_1\), the ratio of the
  two degree-\(1\) generators of \(\mathcal A\). Applying the chart-\(1\) forward map
  \(\mathrm{forward}_1\) of \cref{def:proj_chart_ring_iso} then sends \([X_0/X_1]\) to the
  polynomial-ring generator \(u \in \mathrm{MvPolynomial}\,\{*\}\,\bar k\) (by the defining
  equation \(X_1 \mapsto 1\), \(X_0 \mapsto u\) of the chart-\(1\) iso). Finally, the source-side
  \(\mathtt{pullbackSpecIso}\) bridge with the Gm-twist collapse is the chart-\(1\) instance
  of the same recipe \cref{def:gmscaling_chart} uses: at the identity section \(\lambda = 1\)
  of \(\mathbb G_m\), the Gm-twist of the chart-\(1\) generator is the chart-\(1\) generator
  itself, and \(\mathtt{Algebra.TensorProduct.lid}\) eliminates the trivial
  \(\mathtt{GmRing}\,\bar k\)-factor to send \(u\) to \(1 \in \bar k\). The combined Mathlib
  helper \(\mathtt{IsOpenImmersion.lift\_app}\) (which expresses the \(\mathtt{lift}\)
  identity at the global-sections level for an open immersion) packages this evaluation in
  the form that the consumer \(\mathtt{iotaGm\_chart1\_composition\_isOpenImmersion}\) (the
  open-immersion certificate that III.c step \(3\)/\(4\) discharges) and
  \(\mathtt{iotaGm\_r\_1\_fac}\) (the iter-190/191 lift identity) directly consume. The
  evaluation is purely Mathlib-side once the chart-\(1\) forward identification is in
  hand --- no project-bespoke ring-map invariant is introduced.
\end{proof}
\subsection*{Iter-195 analogist recipe --- \(\mathtt{Proj.awayι\_app\_basicOpen}\) port (Lane E close)}
% SOURCE: Mathlib `Mathlib/AlgebraicGeometry/AffineScheme.lean:560-564`
% (the `IsAffineOpen.fromSpec_app_self` formula) and
% `Mathlib/AlgebraicGeometry/ProjectiveSpectrum/Basic.lean:143-185`
% (the `basicOpenToSpec_app_top` + `basicOpenIsoSpec` + `basicOpenIsoAway`
% chain). Read by the iter-195 mathlib-analogist
% `lane-e-proj-appiso-pivot` (cross-domain-inspiration mode);
% verdict ANALOGUE_FOUND.
% NOTE (iter-196 blueprint-writer avr-lane-e-recipe): this subsection records
% the explicit Mathlib-named recipe the iter-196 prover dispatches to close the
% Lane E chart-1 `appTop` residual at line 273 of
% `AlgebraicJacobian/AbelianVarietyRigidity.lean` (the typed `sorry` carried by
% \texttt{kbarChart1Ring\_specMap\_fac}). The iter-196 progress-critic primary
% corrective is precisely \emph{explicit Lean API specification in the
% blueprint}: the analogist file
% \texttt{analogies/lane-e-proj-appiso-pivot.md} is too long and cross-domain
% to feed the prover directly. The two new \(\mathtt{\lean}\) pins below
% (\cref{lem:awayi_app_basicOpen} and
% \cref{lem:awayi_appIso_top_inv_apply_isLocElem}) name the new declarations
% the prover will land; the third block records the consumer dispatch at the
% existing typed sorry.
% Lane E history (iter-187 -- iter-195 STUCK). After the iter-187 pivot to
% the separated-locus alternative (\cref{lem:gmscaling_chart_agreement}, path
% (III.c)), the chart-\(1\) close devolved to a single \(\mathtt{Proj.appIso}\)
% evaluation residual: identify
% \(((\mathrm{Proj.awayι}\,X_1).\mathrm{appIso}\,\top).\mathrm{inv}\,(\mathtt{isLocElem})\)
% with \([X_0/X_1] \in \mathrm{HomogeneousLocalization.Away}\,\mathcal A\,X_1\). The
% iter-188/189/190/191 prover attempts (direct \(\mathtt{simp + rw}\) chains
% through \(\mathtt{Proj.basicOpen\_eq\_basicOpenIso} \blacktriangleright
% \mathtt{Proj.appIso\_inv}\)) hit Lean term-unification timeouts; iter-192
% factored through the named helper \cref{lem:iotaGm_chart1_appIso_eval} but
% that helper carries the same residual; iter-193 re-routed through
% \(\mathtt{IsOpenImmersion.lift\_uniq}\) (eliminating the outer
% morphism-level reasoning, leaving the inner \(\mathtt{appTop}\) residual
% intact); iter-194 structurally isolated the residual at the consumer
% \(\mathtt{kbarChart1Ring\_specMap\_fac}\) (line 273 of
% \(\mathtt{AlgebraicJacobian/AbelianVarietyRigidity.lean}\)) without
% closing it. The iter-195 \texttt{cross-domain-inspiration} analogist found
% the right Mathlib model: the substantive content of the residual is the
% \(\mathrm{Proj}\)-side mirror of
% \(\mathtt{IsAffineOpen.fromSpec\_app\_self}\)
% (\(\mathtt{Mathlib/AlgebraicGeometry/AffineScheme.lean:560\text{--}564}\)),
% whose proof shape transports verbatim. The three blocks below capture the
% port and its consumer dispatch.
% NOTE (iter-197 blueprint-writer): the iter-196 prover identified the
% intermediate \(\mathrm{Proj.basicOpenIsoSpec}.\mathrm{inv}.\mathrm{app}\,\top\)
% identity as the missing named helper whose absence caused the dependent-motive
% blocker on \cref{lem:awayi_app_basicOpen}. It is now pinned below as
% \cref{lem:basicOpenIsoSpec_inv_app_top}, immediately preceding the section-level
% port. The two iter-196 landed primitives (the preimage identity and the
% \(\mathrm{Spec.map} \fatsemi \mathrm{fromSpec}\) bridge) are also pinned just
% above for blueprint visibility.
\begin{lemma}

\leanok
[Preimage of \(\mathrm{Proj.basicOpen}\) under \(\mathrm{Proj.awayι}\) is \(\top\)]
  \label{lem:awayi_preimage_basicOpen_self}
  \lean{AlgebraicGeometry.Proj.awayι_preimage_basicOpen_self}
  % Project-local Mathlib supplement (axiom-clean, iter-196 landed). Combines
  % Mathlib's \(\mathrm{Proj.opensRange\_awayι}\) with
  % \(\mathrm{Scheme.Hom.preimage\_opensRange}\). No external textbook source.
  Let \(\mathcal A \colon \mathbb N \to \sigma\) be an \(\mathbb N\)-graded ring,
  \(f \in \mathcal A_m\) homogeneous of degree \(m > 0\). The preimage of the
  basic open \(\mathrm{Proj.basicOpen}\,\mathcal A\,f\) under the canonical
  Spec-morphism \(\mathrm{Proj.awayι}\,\mathcal A\,f\,(\mathtt{f\_deg})\,
  (\mathtt{hm})\) is the whole space:
  \[ \mathrm{Proj.awayι}\,\mathcal A\,f\,(\mathtt{f\_deg})\,(\mathtt{hm})
       \,\Vert^{-1}\,\mathrm{Proj.basicOpen}\,\mathcal A\,f
     \ =\ \top. \]
  This is a direct corollary of the range of \(\mathrm{Proj.awayι}\) being
  exactly \(D_+(f)\), via the canonical
  \(\mathrm{Scheme.Hom.preimage\_opensRange}\) rewrite. \(\sim 2\) LOC in Lean.
\end{lemma}
\begin{lemma}

\leanok
[Bridge \(\mathrm{Proj.awayι}\) to \(\mathtt{IsAffineOpen.fromSpec}\)]
  \label{lem:awayi_eq_specMap_fromSpec}
  \lean{AlgebraicGeometry.Proj.awayι_eq_specMap_fromSpec}
  % Project-local Mathlib supplement (axiom-clean, iter-196 landed). Step (i) of
  % the iter-195 analogist recipe `lane-e-proj-appiso-pivot`. No external
  % textbook source: the substantive content is the categorical
  % `Iso.ext`-via-`basicOpenIsoSpec_hom` cancellation; the proof shape is
  % project-bespoke.
  Let \(\mathcal A \colon \mathbb N \to \sigma\) be an \(\mathbb N\)-graded ring,
  \(f \in \mathcal A_m\) homogeneous of degree \(m > 0\). Then the canonical
  Spec-morphism \(\mathrm{Proj.awayι}\) factors through Mathlib's
  \(\mathtt{IsAffineOpen.fromSpec}\) via the chart-ring iso
  \(\mathrm{basicOpenIsoAway}\):
  \[ \mathrm{Proj.awayι}\,\mathcal A\,f\,(\mathtt{f\_deg})\,(\mathtt{hm})
     \ =\
     \Spec.\mathrm{map}\,(\mathrm{Proj.basicOpenIsoAway}\,\mathcal A\,f\,
       (\mathtt{f\_deg})\,(\mathtt{hm})).\mathrm{inv}
     \,\fatsemi\,
     (\mathrm{Proj.isAffineOpen\_basicOpen}\,\mathcal A\,f\,(\mathtt{f\_deg})\,
       (\mathtt{hm})).\mathrm{fromSpec}. \]
  The proof unfolds \(\mathrm{Proj.awayι}\) and \(\mathtt{IsAffineOpen.fromSpec}\)
  to expose the parallel \(\mathrm{basicOpenIsoSpec}.\mathrm{inv} \fatsemi \iota\)
  factorisation, then applies \(\mathtt{Iso.ext}\) on the forward direction via
  the chain \(\mathtt{Proj.basicOpenIsoSpec\_hom}\),
  \(\mathtt{Proj.basicOpenToSpec}\), \(\mathtt{IsAffineOpen.isoSpec\_hom}\),
  \(\mathtt{Functor.mapIso\_hom}\). \(\sim 10\) LOC in Lean.
\end{lemma}
\begin{lemma}


\leanok
[Inversion of \(\mathrm{Proj.basicOpenIsoSpec}\) at \(\top\)]
  \label{lem:basicOpenIsoSpec_inv_app_top}

  \lean{AlgebraicGeometry.Proj.basicOpenIsoSpec_inv_app_top}
  % Project-bespoke section-level intermediate isolated from the proof of
  % \cref{lem:awayi_app_basicOpen}. The iter-196 prover's task report explicitly
  % named this as the missing helper whose absence caused the dependent-motive
  % blocker on `Scheme.Hom.app` (`.archon/task_results/AlgebraicJacobian/
  % AbelianVarietyRigidity.md:66-71`). Closed form via the same Mathlib
  % `Scheme.invApp` + `Proj.basicOpenIsoSpec_hom` + `Proj.basicOpenToSpec_app_top`
  % chain; no external textbook source.
  Let \(\mathcal A \colon \mathbb N \to \sigma\) be an \(\mathbb N\)-graded ring,
  let \(f \in \mathcal A_m\) be a homogeneous element of strictly positive
  degree \(m > 0\), and let \(\mathrm{Proj.basicOpenIsoSpec}\,\mathcal A\,f\,
  (\mathtt{f\_deg})\,(\mathtt{hm})\) be the canonical iso of schemes
  \(\Proj.\mathrm{basicOpen}\,\mathcal A\,f \,\cong\,
  \Spec\,(\mathrm{HomogeneousLocalization.Away}\,\mathcal A\,f)\). Then the
  inverse evaluated at the whole space \(\top\) admits the closed-form
  factorisation
  \begin{align*}
    (\mathrm{Proj.basicOpenIsoSpec}\,\mathcal A\,f\,(\mathtt{f\_deg})\,
       (\mathtt{hm})).\mathrm{inv}.\mathrm{app}\,\top
    \ =\ \;
      &(\mathrm{Proj.basicOpen}\,\mathcal A\,f).\mathrm{topIso.hom} \\
      &\quad\fatsemi\,
      (\mathrm{Proj.basicOpenIsoAway}\,\mathcal A\,f\,(\mathtt{f\_deg})\,
        (\mathtt{hm})).\mathrm{inv} \\
      &\quad\fatsemi\,
      (\mathrm{Scheme}.\Gamma\mathrm{SpecIso}\,\_).\mathrm{inv}.
  \end{align*}
  Equivalently, the inverse of the scheme iso \(\mathrm{basicOpenIsoSpec}\),
  evaluated on the whole space \(\top\), factors as the chain
  \(\mathrm{topIso}\)-followed-by-inversions of the two component isos
  (chart-ring iso and \(\Gamma\mathrm{SpecIso}\)).
\end{lemma}
\begin{proof}
  Set \(\beta := \mathrm{Proj.basicOpenIsoSpec}\,\mathcal A\,f\,(\mathtt{f\_deg})
  \,(\mathtt{hm})\). By Mathlib's \(\mathtt{Scheme.invApp}\) (the
  inverse-of-iso lemma for scheme presheaf maps), the section
  \(\beta.\mathrm{inv}.\mathrm{app}\,\top\) is the categorical inverse of the
  ring map \(\beta.\mathrm{hom}.\mathrm{app}\,\top\) on the appropriate
  presheaf. Substitute the explicit value
  \(\beta.\mathrm{hom} = \mathrm{Proj.basicOpenToSpec}\,\mathcal A\,f\,
  (\mathtt{f\_deg})\,(\mathtt{hm})\) via Mathlib's
  \(\mathtt{Proj.basicOpenIsoSpec\_hom}\). Then evaluate
  \(\mathrm{Proj.basicOpenToSpec}.\mathrm{app}\,\top\) via Mathlib's
  \(\mathtt{Proj.basicOpenToSpec\_app\_top}\), which exposes the explicit
  factor chain
  \[ (\mathrm{Scheme}.\Gamma\mathrm{SpecIso}\,\_).\mathrm{hom}
     \,\fatsemi\,
     \mathrm{awayToSection}\,\mathcal A\,f
     \,\fatsemi\,
     (\mathrm{Proj.basicOpen}\,\mathcal A\,f).\mathrm{topIso.inv}, \]
  where \(\mathrm{awayToSection}\,\mathcal A\,f\) is the canonical (forward)
  ring map that \(\mathrm{Proj.basicOpenIsoAway}.\mathrm{hom}\) is implemented
  as. Inverting this composition (composition of isos reverses) yields the RHS
  of the lemma statement. The simplification is mechanical
  (\(\mathtt{Iso.inv\_hom\_id}\) cancels each factor's round trip in turn);
  LOC estimate: \(\sim 5\)--\(15\).
\end{proof}
\begin{lemma}


\leanok
[Section-level formula for \(\mathrm{Proj.awayι}\) on its image: the
\(\mathtt{IsAffineOpen.fromSpec\_app\_self}\) port]
  \label{lem:awayi_app_basicOpen}
  \uses{lem:basicOpenIsoSpec_inv_app_top, lem:awayi_eq_specMap_fromSpec, lem:awayi_preimage_basicOpen_self}

  \lean{AlgebraicGeometry.Proj.awayι_app_basicOpen}
  % Project-bespoke port of `IsAffineOpen.fromSpec_app_self` (Mathlib
  % `AffineScheme.lean:560-564`) onto the `Proj.awayι` open immersion.
  % The cited Mathlib lemma is the structural template; no external
  % textbook source is invoked. The port is documented in detail by the
  % iter-195 analogist at `analogies/lane-e-proj-appiso-pivot.md`
  % (cross-domain-inspiration mode, ranked analogue #1, ANALOGUE_FOUND).
  Let \(\mathcal A \colon \mathbb N \to \sigma\) be a graded \(\mathbb N\)-graded ring
  (the project's \(\mathcal A = \mathrm{projectiveLineBarGrading}\,\bar k\) on
  \(A = \mathrm{MvPolynomial}\,(\mathrm{Fin}\,2)\,\bar k\) is the load-bearing
  instance), let \(f \in \mathcal A_m\) be a homogeneous element of strictly
  positive degree \(m > 0\) (e.g.\ \(f = X_1\) of degree \(1\)), and let
  \[ \mathrm{Proj.awayι}\,\mathcal A\,f\,(\mathtt{f\_deg})\,(\mathtt{hm}) \colon
     \Spec\,(\mathrm{HomogeneousLocalization.Away}\,\mathcal A\,f)
     \longrightarrow \Proj\,\mathcal A \]
  be the canonical Spec-morphism realising the basic open
  \(D_+(f) \subseteq \Proj\,\mathcal A\) as an affine open. Then the section-level
  map of \(\mathrm{Proj.awayι}\) at the codomain open \(\mathrm{Proj.basicOpen}\,\mathcal A\,f\)
  admits the explicit closed-form formula
  \[ (\mathrm{Proj.awayι}\,\mathcal A\,f\,(\mathtt{f\_deg})\,(\mathtt{hm}))
       .\mathrm{app}\,(\mathrm{Proj.basicOpen}\,\mathcal A\,f)
     \ =\
     (\mathrm{Proj.basicOpenIsoAway}\,\mathcal A\,f\,(\mathtt{f\_deg})\,(\mathtt{hm})).\mathrm{inv}
       \,\fatsemi\,
     (\mathrm{Scheme}.\Gamma\mathrm{SpecIso}\,\_).\mathrm{inv}
       \,\fatsemi\,
     (\Spec\,\_).\mathrm{presheaf.map}\,
       (\mathrm{eqToHom}\,(\mathtt{opensRange\_awayι})).\mathrm{op}. \]
  Equivalently, the \emph{forward} (covariant) section-level pullback by
  \(\mathrm{Proj.awayι}\) on the canonical basic open \(\mathrm{Proj.basicOpen}\,\mathcal A\,f\)
  is exactly the composite of (i) the inverse of the project-recognised iso
  \(\mathrm{Proj.basicOpenIsoAway} \colon \mathrm{Away}\,\mathcal A\,f
  \xrightarrow{\ \sim\ } \Gamma(\Proj\,\mathcal A, \mathrm{basicOpen}\,\mathcal A\,f)\),
  (ii) the inverse of the canonical \(\Gamma\)-\(\Spec\) adjunction iso
  \(\mathrm{Scheme}.\Gamma\mathrm{SpecIso}\), and (iii) a pure
  \(\mathrm{eqToHom}\) presheaf-map adjustment along the equality
  \(\mathtt{opensRange\_awayι} \colon (\mathrm{Proj.awayι}\,\_\,f\,\_\,\_)
  ''\,\top = \mathrm{basicOpen}\,\mathcal A\,f\) (the equation that the range of
  \(\mathrm{Proj.awayι}\) is the basic open \(D_+(f)\)). No \(\mathtt{appIso}\)
  chase is needed: the formula is direct.
\end{lemma}
\begin{proof}

  The proof MIRRORS the proof of \(\mathtt{IsAffineOpen.fromSpec\_app\_self}\)
  in Mathlib (\(\mathtt{Mathlib/AlgebraicGeometry/AffineScheme.lean:560\text{--}564}\))
  and proceeds in four steps. The substrate primitives below are all named
  Mathlib lemmas at the pinned commit \(\mathtt{b80f227}\) (verified by the
  iter-195 analogist); no Mathlib upstream gap is incurred.
  \medskip
  \noindent\textbf{Step 1 (factor \(\mathrm{Proj.awayι}\) through
  \(\mathtt{IsAffineOpen.fromSpec}\)).} Apply the iter-196 landed bridge
  \(\cref{lem:awayi_eq_specMap_fromSpec}\) (Lean:
  \(\mathtt{Proj.awayι\_eq\_specMap\_fromSpec}\)) to rewrite
  \(\mathrm{Proj.awayι}\,\mathcal A\,f\,(\mathtt{f\_deg})\,(\mathtt{hm})\) on the
  LHS as
  \[ \Spec.\mathrm{map}\,(\mathrm{Proj.basicOpenIsoAway}\,\mathcal A\,f\,
       (\mathtt{f\_deg})\,(\mathtt{hm})).\mathrm{inv}
     \ \fatsemi\
     (\mathrm{Proj.isAffineOpen\_basicOpen}\,\mathcal A\,f\,(\mathtt{f\_deg})\,
       (\mathtt{hm})).\mathrm{fromSpec}. \]
  Then apply \(\mathtt{Scheme.Hom.comp\_app}\)
  \((g_1 \fatsemi g_2).\mathrm{app}\,U = g_2.\mathrm{app}\,U
   \,\fatsemi\, g_1.\mathrm{app}(g_2^{-1}\,U)\)
  to split the section-level value at
  \(\mathrm{Proj.basicOpen}\,\mathcal A\,f\) into the composition of the two
  scheme-map sections:
  \((\mathtt{IsAffineOpen.fromSpec}).\mathrm{app}(\mathrm{Proj.basicOpen}\,\mathcal A\,f)\)
  followed by \(\Spec.\mathrm{map}\,(\mathrm{basicOpenIsoAway}).\mathrm{inv}.\mathrm{app}\)
  at the preimage. \textbf{This route avoids the dependent-motive obstruction
  that blocked iter-188--iter-196:} both sides of the factorisation keep the
  same \(\mathrm{app}\)-codomain type, because the preimage of
  \(\mathrm{Proj.basicOpen}\,\mathcal A\,f\) under either \(\mathrm{Proj.awayι}\)
  or the composite \(\Spec.\mathrm{map}(\ldots) \fatsemi \mathrm{fromSpec}\)
  agrees with \(\top\) under
  \(\cref{lem:awayi_preimage_basicOpen_self}\) (Lean:
  \(\mathtt{Proj.awayι\_preimage\_basicOpen\_self}\), iter-196 landed).
  \medskip
  \noindent\textbf{Step 2 (evaluate \(\mathrm{fromSpec}.\mathrm{app}\,\top\) via
  Mathlib).} Apply Mathlib's
  \(\mathtt{IsAffineOpen.fromSpec\_app\_self}\) (at
  \(\mathtt{Mathlib/AlgebraicGeometry/AffineScheme.lean:560\text{--}564}\)) to
  the affine open \(\mathrm{Proj.basicOpen}\,\mathcal A\,f\) (witness:
  \(\mathrm{Proj.isAffineOpen\_basicOpen}\)). This yields the explicit
  closed-form formula
  \[ (\mathrm{Proj.isAffineOpen\_basicOpen}\,\mathcal A\,f\,(\mathtt{f\_deg})\,
       (\mathtt{hm})).\mathrm{fromSpec}.\mathrm{app}\,
       (\mathrm{Proj.basicOpen}\,\mathcal A\,f)
     \ =\
     (\mathrm{Scheme}.\Gamma\mathrm{SpecIso}\,\_).\mathrm{inv}
     \,\fatsemi\,
     (\Spec\,\_).\mathrm{presheaf.map}\,
       (\mathrm{eqToHom}\,(\ldots)).\mathrm{op}, \]
  i.e.\ the \(\mathrm{fromSpec}\) section at the canonical basic open factors
  as \(\Gamma\mathrm{SpecIso}\)-inverse followed by a pure \(\mathrm{eqToHom}\)
  presheaf-map adjustment. (This is the iter-195 analogist's identified
  Mathlib template; \cref{lem:basicOpenIsoSpec_inv_app_top} packages the
  morally-equivalent intermediate on the \(\mathrm{basicOpenIsoSpec}.\mathrm{inv}\)
  side of the iso and is consumed by the alternative
  \(\mathrm{basicOpenIsoSpec}.\mathrm{inv} \fatsemi \iota\) route.)
  \medskip
  \noindent\textbf{Step 3 (evaluate \(\Spec.\mathrm{map}(\ldots).\mathrm{app}\,\_\)
  via \(\mathtt{Spec.map\_appTop}\)).} The remaining
  \(\Spec.\mathrm{map}\,(\mathrm{basicOpenIsoAway}\,\mathcal A\,f\,(\mathtt{f\_deg})\,
  (\mathtt{hm})).\mathrm{inv}\) factor of Step~1 evaluates on sections by
  Mathlib's \(\mathtt{Spec.map\_appTop}\): for any ring map
  \(\varphi \colon R \to S\), \((\Spec.\mathrm{map}\,\varphi).\mathrm{app}\,\top
  = (\mathrm{Scheme}.\Gamma\mathrm{SpecIso}\,R).\mathrm{hom}\,\fatsemi\,\varphi
  \,\fatsemi\,(\mathrm{Scheme}.\Gamma\mathrm{SpecIso}\,S).\mathrm{inv}\) (modulo
  the canonical \(\mathrm{eqToHom}\) bookkeeping). Specialised to
  \(\varphi = (\mathrm{basicOpenIsoAway}).\mathrm{inv}\), this gives the
  \(\mathrm{basicOpenIsoAway}.\mathrm{inv}\) factor of the target formula
  (surrounded by the canonical \(\Gamma\mathrm{SpecIso}\) bookkeeping).
  \medskip
  \noindent\textbf{Step 4 (combine).} Composing the section-value formulas of
  Steps~2 (the \(\mathrm{fromSpec}.\mathrm{app}\) factor) and~3 (the
  \(\Spec.\mathrm{map}.\mathrm{app}\) factor) via the
  \(\mathtt{Scheme.Hom.comp\_app}\) split of Step~1, the canonical
  \(\Gamma\mathrm{SpecIso}.\mathrm{hom} \fatsemi \Gamma\mathrm{SpecIso}.\mathrm{inv}
  = \mathbf 1\) round-trip cancels (via \(\mathtt{Iso.hom\_inv\_id\_assoc}\)),
  and the surviving \(\mathrm{eqToHom}\) presheaf-map collapses against
  \(\mathtt{opensRange\_awayι}\) listed in the conclusion. The result is exactly
  the RHS of the lemma statement. This mirrors the three-line Mathlib proof of
  \(\mathtt{IsAffineOpen.fromSpec\_app\_self}\) step-for-step on the
  \(\mathrm{Proj}\)-side, replacing \(\mathtt{IsAffineOpen.fromSpec}\) with
  \(\mathrm{Proj.awayι}\) via the iter-196 bridge
  \cref{lem:awayi_eq_specMap_fromSpec} and Mathlib's affine
  \(\mathtt{ΓSpecIso\_hom\_fromSpec\_app}\) with the \(\mathrm{Proj}\)-side
  companion \(\mathtt{Proj.basicOpenToSpec\_app\_top}\) (consumed in
  \cref{lem:basicOpenIsoSpec_inv_app_top} as its substrate). LOC estimate:
  \(\sim 10\)--\(15\).
\end{proof}
\begin{lemma}


\leanok
[Point-value evaluation of \(((\mathrm{Proj.awayι}).\mathrm{appIso}\,\top).\mathrm{inv}\)
on the canonical generator \(\mathtt{isLocElem}\)]
  \label{lem:awayi_appIso_top_inv_apply_isLocElem}
  \uses{lem:awayi_app_basicOpen}

  \lean{AlgebraicGeometry.Proj.awayι_appIso_top_inv}
  % NOTE (iter-197): renamed pin from `_apply_isLocElem` (point-value form) to
  % `awayι_appIso_top_inv` (morphism-level form) per the iter-197 prover task
  % report. The morphism-level form is strictly stronger (it implies the
  % point-value form via `congr`), and is what the consumer `kbarChart1Ring_specMap_fac`
  % actually uses through the `IsOpenImmersion.lift_app` chain.
  % Project-bespoke point-value corollary of \cref{lem:awayi_app_basicOpen}.
  % Combines that section-level formula with Mathlib's
  % `Scheme.Hom.appIso_hom` (`Mathlib/AlgebraicGeometry/OpenImmersion.lean:199`)
  % to obtain the closed-form evaluation of the inverse section-level iso on
  % the canonical generator `isLocElem` --- exactly the residual the
  % consumer `kbarChart1Ring_specMap_fac` needs.
  Continue the setup of \cref{lem:awayi_app_basicOpen}: \(\mathcal A\) a graded
  ring, \(f \in \mathcal A_m\) homogeneous of degree \(m > 0\). Mathlib's
  \(\mathtt{Scheme.Hom.appIso}\,\top\) (at
  \(\mathtt{Mathlib/AlgebraicGeometry/OpenImmersion.lean:199}\)) supplies the
  section-level iso
  \[ ((\mathrm{Proj.awayι}\,\mathcal A\,f\,(\mathtt{f\_deg})\,(\mathtt{hm})).\mathrm{appIso}\,\top)
     \colon \Gamma(\Proj\,\mathcal A, \mathrm{basicOpen}\,\mathcal A\,f)
     \ \xrightarrow{\ \sim\ }\
     \Gamma(\Spec\,(\mathrm{Away}\,\mathcal A\,f),\,\top), \]
  packaging the canonical identification at the
  \(\mathrm{opensRange}\,\mathrm{Proj.awayι} = \mathrm{basicOpen}\,\mathcal A\,f\)
  pair. Evaluated in the inverse direction at the canonical generator
  \(\mathtt{isLocElem} \in \Gamma(\Spec\,(\mathrm{Away}\,\mathcal A\,f),\,\top)\)
  --- equivalently \(\mathrm{IsLocalization.mk}'\,(\mathrm{Localization.Away}\,f)
  \,1\,\langle f, 1, \mathrm{rfl}\rangle\) in the affine global-sections form,
  i.e.\ the homogeneous-localisation element corresponding to the affine
  \(\Spec\)-identity-ring-map under the
  \(\Gamma(\Spec(\cdot))\cong(\cdot)\) adjunction ---  it sends
  \(\mathtt{isLocElem}\) to the homogeneous-localisation ratio
  \([X_0/X_1] \in \mathrm{HomogeneousLocalization.Away}\,\mathcal A\,f\)
  \((\)concretely \(\mathrm{Proj.basicOpenIsoAway.hom}\,(\mathtt{mk\_isLocElem})\)
  in the equivalent statement form\()\); for the project's
  \(\mathcal A = \mathrm{projectiveLineBarGrading}\,\bar k\) and \(f = X_1\) it
  recovers the long-documented residue identification
  \(\mathtt{isLocElem} \mapsto [X_0/X_1]\) of the iter-188 -- iter-194
  Lane E narrative.
\end{lemma}
\begin{proof}

  The proof combines the section-level formula \cref{lem:awayi_app_basicOpen}
  with two named Mathlib substrate lemmas; no further substantive content is
  needed.
  \medskip
  \noindent\textbf{Step 1 (\(\mathrm{appIso.hom}\) in terms of \(\mathrm{app}\)).}
  Mathlib's \(\mathtt{Scheme.Hom.appIso\_hom}\) at
  \(\mathtt{Mathlib/AlgebraicGeometry/OpenImmersion.lean:199}\) exposes the
  forward direction \(((\mathrm{Proj.awayι}).\mathrm{appIso}\,\top).\mathrm{hom}\)
  as \((\mathrm{Proj.awayι}).\mathrm{app}\,(\mathrm{opensRange}\,\mathrm{Proj.awayι}
  = \mathrm{basicOpen}\,\mathcal A\,f)\) composed with a pure
  \(\mathrm{eqToHom}\) presheaf-map adjustment. Substituting the closed-form
  expression for \((\mathrm{Proj.awayι}).\mathrm{app}\,(\mathrm{basicOpen}\,\mathcal A\,f)\)
  from \cref{lem:awayi_app_basicOpen} gives an explicit closed form for the
  forward iso \(((\mathrm{Proj.awayι}).\mathrm{appIso}\,\top).\mathrm{hom}\) in
  terms of \(\mathrm{Proj.basicOpenIsoAway}.\mathrm{inv}\),
  \((\mathrm{Scheme}.\Gamma\mathrm{SpecIso}\,\_).\mathrm{inv}\), and named
  \(\mathrm{eqToHom}\) presheaf-map factors.
  \medskip
  \noindent\textbf{Step 2 (flip to the forward direction via
  \(\mathtt{Iso.eq\_inv\_apply}\)).} Apply
  \(\mathtt{Iso.eq\_inv\_apply}\) (Mathlib
  \(\mathtt{CategoryTheory/Iso.lean}\)) to reduce the inverse-evaluation
  \(((\mathrm{Proj.awayι}).\mathrm{appIso}\,\top).\mathrm{inv}\,(\mathtt{isLocElem})
  = [X_0/X_1]\) to the equivalent forward statement
  \[ ((\mathrm{Proj.awayι}).\mathrm{appIso}\,\top).\mathrm{hom}\,([X_0/X_1])
     \ =\ \mathtt{isLocElem}, \]
  i.e.\ \(\mathrm{Proj.basicOpenIsoAway.hom}\,(\mathtt{mk\_isLocElem})
  = (\mathrm{Proj.awayι}).\mathrm{app}\,(\mathrm{basicOpen}\,\mathcal A\,f)\,
  (\mathtt{isLocElem})\). The forward direction is structurally cleaner:
  \(\mathrm{Proj.basicOpenIsoAway}.\mathrm{hom}\) is \emph{implemented as}
  \(\mathrm{awayToSection}\,\mathcal A\,f\) (see
  \(\mathtt{Mathlib/.../ProjectiveSpectrum/Basic.lean:178}\)), so the
  RHS of the forward statement is the section-level evaluation of
  \(\mathrm{Proj.awayι}.\mathrm{app}\) at \(\mathrm{basicOpen}\,\mathcal A\,f\)
  on the canonical generator.
  \medskip
  \noindent\textbf{Step 3 (close by \(\mathrm{awayToSection\_apply}\)).} The
  forward equation
  \(\mathrm{awayToSection}\,\mathcal A\,f\,(\mathtt{mk\_isLocElem})
  = (\mathrm{Proj.awayι}).\mathrm{app}\,(\mathrm{basicOpen}\,\mathcal A\,f)\,
  (\mathtt{isLocElem})\) evaluates by Mathlib's
  \(\mathtt{awayToSection\_apply}\) in
  \(\mathtt{Mathlib/AlgebraicGeometry/ProjectiveSpectrum/Scheme.lean}\):
  closed form on \(\mathrm{IsLocalization.map}\,\_\,\mathrm{HomogeneousLocalization.val}\)
  identifies both sides as the canonical homogeneous-localisation ratio. LOC
  estimate: \(\sim 5\)--\(10\).
\end{proof}
\noindent\textbf{Consumer dispatch (iter-196 prover).} The existing
typed sorry at line \(273\) of
\(\mathtt{AlgebraicJacobian/AbelianVarietyRigidity.lean}\) in
\(\mathtt{kbarChart1Ring\_specMap\_fac}\) reduces, after
\cref{lem:awayi_appIso_top_inv_apply_isLocElem} is in scope, to a
forward equation between two ring-map expressions on
\(\mathtt{appTop}\). The iter-196 prover dispatch is:
\begin{enumerate}
  \item Rewrite the LHS using
    \(\mathtt{simp\ only}\,[\,\mathtt{Proj.awayι\_appIso\_top\_inv\_apply\_isLocElem}\,]\)
    to substitute the closed form supplied by
    \cref{lem:awayi_appIso_top_inv_apply_isLocElem}.
  \item Identify the RHS by \(\mathtt{ext\_of\_isAffine}\) plus
    \(\mathtt{awayToSection\_apply}\), exhibiting the resulting ring map as
    \(\mathrm{IsLocalization.map}\,\_\,\mathrm{HomogeneousLocalization.val}\)
    on the canonical generator.
  \item Close by definitional unfolding of \(\mathtt{kbarChart1Ring}\)
    and the canonical
    \(\mathrm{MvPolynomial.eval}\,(\lambda \_ \mapsto 1)\) collapse,
    matching the chart-\(1\) ring map.
\end{enumerate}
This dispatch is the substantive content the iter-188 -- iter-194
attempts converged on; it is now packaged in the two named helpers
\cref{lem:awayi_app_basicOpen} and
\cref{lem:awayi_appIso_top_inv_apply_isLocElem}, with the consumer-site
discharge estimated at \(\sim 5\)--\(10\) LOC. The same recipe also closes
the residual sorry inside \cref{lem:iotaGm_chart1_appIso_eval} (line
\(481\) of \(\mathtt{AlgebraicJacobian/AbelianVarietyRigidity.lean}\)),
which shares the same \(\mathtt{Proj.appIso}\) evaluation residual modulo
the iso-chain functoriality already discharged in stages 5--8 of that
helper's proof.
\begin{lemma}
\leanok
[\(\overline{\mathbb P^1}\) is reduced]
  \label{lem:projlinebar_isReduced}
  \lean{AlgebraicGeometry.projectiveLineBar_isReduced}
  \uses{def:projlinebar_affine_cover}
  % Archon-original Lean realisation of an elementary textbook fact (projective space over
  % a field is reduced) via the explicit chart cover of \cref{def:projlinebar_affine_cover}
  % and the canonical \(\mathrm{val}\)-injection of a homogeneous localization into the
  % corresponding plain localization. No external % SOURCE QUOTE needed.
  The scheme underlying \(\overline{\mathbb P^1}\) is reduced. By the two-chart affine cover
  (\cref{def:projlinebar_affine_cover}) it suffices to show that each chart ring is
  reduced. The chart-\(i\) ring \(\mathrm{Away}\,\mathcal A\,X_i\) embeds into the ordinary
  localization \(\bar k[X_0, X_1][X_i^{-1}]\) of the polynomial ring \(\bar k[X_0, X_1]\) at
  the powers of \(X_i\) via the canonical injection \(\mathrm{val} \colon
  \mathrm{Away}\,\mathcal A\,X_i \hookrightarrow \bar k[X_0, X_1][X_i^{-1}]\)
  (\(\mathtt{HomogeneousLocalization.val\_injective}\)); the codomain is a localization of a
  polynomial ring at a non-zero-divisor (the powers of \(X_i\)), hence a domain, hence
  reduced. A subring of a reduced ring is reduced
  (\(\mathtt{Function.Injective.isDomain}\) on the \(\mathrm{val}\)-embedding), so each chart
  ring is a domain --- in particular reduced --- and the global reducedness follows.
\end{lemma}
\begin{definition}
\leanok
[The \(\mathbb G_m\)-scaling action \(\sigma_\times \colon \mathbb P^1 \times \mathbb G_m \to \mathbb P^1\) (primary), and the \(\mathbb G_a\)-translation companion \(\sigma\)]
  \label{def:gaTranslationP1}

  \lean{AlgebraicGeometry.gmScalingP1}
  \uses{def:genus0_base_objects, def:gmscaling_cover, def:gmscaling_chart, lem:gmscaling_chart_agreement, lem:gmscaling_over_coherence}
  % NOTE (iter-171 refresh of iter-170 decision): the option (c) inline chart-glue
  % commitment LANDED its body skeleton iter-171. `gmScalingP1` now has a CONCRETE
  % body via `Over.homMk ((gmScalingP1_cover).glueMorphisms gmScalingP1_chart
  % gmScalingP1_chart_agreement) gmScalingP1_over_coherence` — no buried sorries.
  % Step A (chart-side ring maps `gmScalingP1_chart{0,1}_ringMap`) + the cover
  % `gmScalingP1_cover` are axiom-clean iter-171. The body's `sorryAx` taint flows
  % through three named top-level internal scaffold sorries:
  % `gmScalingP1_chart` (per-chart scheme morphism), `gmScalingP1_chart_agreement`
  % (cocycle), `gmScalingP1_over_coherence` (over-side coherence). iter-172+ closes
  % these via `pullbackSpecIso ≫ Spec.map _ ≫ Proj.awayι` + `Scheme.Cover.hom_ext`.
  % See `analogies/gmscaling-deep.md` for the full 6-step decomposition.
  % Archon-original Lean encoding; the chart computations below are elementary and characteristic-free.
  % No external % SOURCE QUOTE for the encoding. The mathematical fact -- \(\mathbb G_m\) acts on
  % \(\mathbb P^1\) by scalings fixing \(0\) and \(\infty\), \(\mathbb G_a\) by translations fixing \(\infty\)
  % -- is standard (the Borel subgroup of \(\mathrm{PGL}_2\) acting on \(\mathbb P^1\)).
  % NOTE (iter-164/167): the load-bearing action for the PRIMARY genus-0 route is now the
  % MULTIPLICATIVE scaling action \(\sigma_\times\) (\lean name gmScalingP1, landed iter-165/166;
  % pinned in this block's \lean{...} hook). The \(\mathbb G_a\)-translation \(\sigma\) (\lean name
  % gaTranslationP1 [expected]) is retained only for the demoted additive alternative and is
  % intentionally not yet a Lean target.
  \textbf{The scaling action (primary).} The \emph{multiplicative scaling action}
  \[ \sigma_\times \colon \mathbb P^1 \times \mathbb G_m \to \mathbb P^1, \qquad
     \text{``}\sigma_\times(x, \lambda) = \lambda x\text{''} \]
  (\lean name \texttt{gmScalingP1}) is a \emph{genuine total} scheme morphism on all of
  \(\mathbb P^1 \times \mathbb G_m\) --- no indeterminacy, characteristic-free (M\"obius scaling
  fixing both \(0\) and \(\infty\)) --- defined chartwise:
  \begin{itemize}
    \item On \(\mathbb A^1 \times \mathbb G_m\) (target chart \(x\)): \((x, \lambda) \mapsto \lambda x\),
      a polynomial map, regular everywhere.
    \item Near \(\infty\) (target chart \(u = 1/x\)): the target coordinate \(1/(\lambda x)\) rewrites as
      \(u/\lambda\), regular everywhere on this chart because \(\lambda \in \mathbb G_m\) is invertible.
  \end{itemize}
  The two target charts cover the image and overlap consistently, so \(\sigma_\times\) is defined on
  the whole surface. It is a left action of the group \(\mathbb G_m\) on \(\mathbb P^1\):
  \(\sigma_\times(x, 1) = x\) and \(\sigma_\times(\sigma_\times(x, \lambda), \mu) = \sigma_\times(x,
  \lambda\mu)\). Crucially it has \emph{two fixed points}, \(\sigma_\times(0, \lambda) = 0\) and
  \(\sigma_\times(\infty, \lambda) = \infty\) for all \(\lambda\); the fixed point \(0\) is what makes the
  \(W\)-axis of the additive decomposition collapse in the primary base-case proof
  (\cref{prop:morphism_P1_to_AV_constant}, \cref{rmk:base_case_fourth_route}).
  \medskip
  \noindent\textbf{Companion: the \(\mathbb G_a\)-translation action (demoted route only).} For the
  Milne-faithful additive alternative we also record the \emph{translation action}
  \[ \sigma \colon \mathbb P^1 \times \mathbb G_a \to \mathbb P^1, \qquad
     \text{``}\sigma(x, y) = x + y\text{''} \]
  (\lean name \texttt{gaTranslationP1} [expected]), likewise a \emph{total} scheme morphism on all
  of \(\mathbb P^1 \times \mathbb G_a\), defined chartwise: on \(\mathbb A^1 \times \mathbb G_a\)
  (chart \(x\)) it is the polynomial map \((x, y) \mapsto x + y\); near \(\infty\) (\(u = 1/x\)) the
  coordinate \(1/(x + y)\) rewrites as \(u/(1 + u y)\), regular at \(u = 0\) (value \(0\) there --- \(\sigma\)
  fixes \(\infty\) for every \(y\)), while the locus \(x + y = 0\) is covered by the chart \(x\) on which it
  is the finite value \(0\). It is a left action of \(\mathbb G_a\) on \(\mathbb P^1\) fixing \(\infty\):
  \(\sigma(x, 0) = x\) and \(\sigma(\sigma(x, y), y') = \sigma(x, y + y')\). This companion underlies
  only the alternative proof of \cref{lem:hom_Ga_to_av_trivial} and the demoted \(\mathbb
  G_a\)-additive close of the base case (\cref{lem:hom_from_Ga_trivial}); it is \emph{not} on the
  genus-\(0\) critical path.
\end{definition}
\subsection*{Scaffolds for the body of \(\sigma_\times = \mathtt{gmScalingP1}\) (iter-171 chart-glue split)}
% NOTE (iter-173, blueprint-writer g0bo-pin-scaffolds): the iter-171 option-(c) inline chart-glue
% commitment factored the body of \texttt{gmScalingP1} (\cref{def:gaTranslationP1}) into a
% \texttt{Scheme.Cover.glueMorphisms} over the pullback cover, with three top-level named
% scaffold sorries naming the per-chart scheme morphism, its cocycle on overlap, and the
% \texttt{Over}-side structure-morphism coherence. The four blocks below pin those scaffolds
% one-by-one so the dependency graph and \texttt{sync\_leanok} can track each declaration
% independently (the iter-172 lean-vs-blueprint-checker \texttt{g0bo172} MAJOR finding).
% Iter-173 prover work on the body of \(\sigma_\times\) targets exactly these three obligations.
% NOTE (iter-186 writer, gmscaling-chart-agreement-expansion): pinning the two project-side
% simp helpers `pullback_map_{fst,snd}_proj` (analogist Recipe 1, iter-184) here so that the
% expanded proof sketch of \cref{lem:gmscaling_chart_agreement} below can `\cref{...}` them
% without an unresolved cross-reference.
\begin{lemma}

\leanok
[Projection simp helper: \(\mathtt{pullback.map} \fatsemi \mathtt{pullback.fst}\)]
  \label{lem:pullback_map_fst_proj}
  \lean{AlgebraicGeometry.pullback_map_fst_proj}
  % Archon-original project-side simp helper (iter-184 Lane B, analogist Recipe 1 of
  % `analogies/gmscaling-projection-idiom.md`). No external % SOURCE QUOTE.
  Let \(C\) be a category and let
  \(f_1 \colon W \to S\), \(f_2 \colon X \to S\), \(g_1 \colon Y \to T\),
  \(g_2 \colon Z \to T\), \(i_1 \colon W \to Y\), \(i_2 \colon X \to Z\), \(i_3 \colon S \to T\)
  be morphisms with \(f_1 \fatsemi i_3 = i_1 \fatsemi g_1\) and
  \(f_2 \fatsemi i_3 = i_2 \fatsemi g_2\) (and both relevant pullbacks existing). Then
  \[ \mathtt{pullback.map}\,f_1\,f_2\,g_1\,g_2\,i_1\,i_2\,i_3\,(\ldots) \,\fatsemi\,
       \mathtt{pullback.fst}\,g_1\,g_2 \ =\
     \mathtt{pullback.fst}\,f_1\,f_2 \,\fatsemi\, i_1. \]
  This identity holds definitionally (its body is \(\mathtt{pullback.lift\_fst}\)), but
  Mathlib ships \(\mathtt{pullback.lift\_fst}\) as \(\mathtt{@[reassoc]}\) only --- not
  \(\mathtt{@[simp]}\) --- so the \(\mathtt{abbrev}\)-form
  \(\mathtt{pullback.map} \fatsemi \mathtt{pullback.fst}\) is not collapsed by \(\mathtt{simp}\).
  Promoting this projection identity to a \(\mathtt{@[reassoc (attr := simp)]}\) lemma is
  the iter-184 Lane B unblock for the
  \(\mathtt{gmScalingP1\_cover\_intersection\_X\_iso}\) chain (whose STEP~4 uses
  \(\mathtt{asIso (pullback.map \ldots)}\) for the inner
  \(\mathtt{Proj.pullbackAwayιIso}\) lift). Candidate for upstream Mathlib contribution.
\end{lemma}
\begin{lemma}

\leanok
[Projection simp helper: \(\mathtt{pullback.map} \fatsemi \mathtt{pullback.snd}\)]
  \label{lem:pullback_map_snd_proj}
  \lean{AlgebraicGeometry.pullback_map_snd_proj}
  % Archon-original project-side simp helper (iter-184 Lane B, analogist Recipe 1 of
  % `analogies/gmscaling-projection-idiom.md`). No external % SOURCE QUOTE.
  In the same data as \cref{lem:pullback_map_fst_proj}, the dual identity holds:
  \[ \mathtt{pullback.map}\,f_1\,f_2\,g_1\,g_2\,i_1\,i_2\,i_3\,(\ldots) \,\fatsemi\,
       \mathtt{pullback.snd}\,g_1\,g_2 \ =\
     \mathtt{pullback.snd}\,f_1\,f_2 \,\fatsemi\, i_2. \]
  Same provenance and same rationale as the \(\mathtt{fst}\) variant: definitional via
  \(\mathtt{pullback.lift\_snd}\), but promoted here to the simp set so the
  iso-chain projection collapse in
  \(\mathtt{gmScalingP1\_cover\_intersection\_X\_iso}\) fires through the
  \(\mathtt{asIso (pullback.map \ldots)}\) link. Used by the expanded proof sketch of
  \cref{lem:gmscaling_chart_agreement} (pickup path (III.a)).
\end{lemma}
\begin{definition}
\leanok
[Pullback open cover of \(\mathbb P^1 \times \mathbb G_m\) along the chart cover of \(\mathbb P^1\)]
  
  \label{def:gmscaling_cover}

  \lean{AlgebraicGeometry.gmScalingP1_cover}
  \uses{def:projlinebar_affine_cover}
  % Archon-original Lean encoding: the affine open cover of \(\mathbb P^1 \times \mathbb G_m\)
  % whose charts the chartwise body of \(\sigma_\times\) is built on. No external % SOURCE QUOTE.
  The two-chart affine open cover of \((\mathbb P^1 \otimes \mathbb G_m).\mathrm{left}\) obtained by
  pulling back the chart cover of \cref{def:projlinebar_affine_cover} along the first pullback
  projection \(\mathrm{pr}_1 \colon \mathbb P^1 \times \mathbb G_m \to \mathbb P^1\) (Mathlib's
  \(\mathtt{pullback.fst}\,\mathbb P^1.\mathrm{hom}\,\mathbb G_m.\mathrm{hom}\)). Concretely the \(i\)-th
  chart is
  \[ (\mathtt{gmScalingP1\_cover})_i =
     \mathrm{pullback}\bigl(\mathtt{pullback.fst}\,\mathbb P^1.\mathrm{hom}\,\mathbb G_m.\mathrm{hom},\
                            \mathrm{Proj.away}\,\iota\,\mathcal A\,X_i\bigr), \]
  i.e.\ the affine open \(D_+(X_i) \times \mathbb G_m \subseteq \mathbb P^1 \times \mathbb G_m\). This
  cover is the \(\mathtt{Scheme.Cover.glueMorphisms}\) source on which the body of \(\sigma_\times =
  \mathtt{gmScalingP1}\) (\cref{def:gaTranslationP1}) is constructed: each per-chart scheme
  morphism (\cref{def:gmscaling_chart}) is built on chart \((\mathtt{gmScalingP1\_cover})_i\), the
  cocycle (\cref{lem:gmscaling_chart_agreement}) checks agreement on each
  pairwise pullback of cover charts, and the over-side coherence
  (\cref{lem:gmscaling_over_coherence}) intertwines the glued total morphism with the structure
  maps to \(\Spec \bar k\).
\end{definition}
\begin{definition}
\leanok
  
[The chart-\(i\) scheme morphism for \(\sigma_\times\)]
  \label{def:gmscaling_chart}

  \lean{AlgebraicGeometry.gmScalingP1_chart}
  \uses{def:gmscaling_cover, def:proj_chart_ring_iso}
  % Archon-original Lean scaffold for the iter-171 chart-glue body of \(\sigma_\times\)
  % (\cref{def:gaTranslationP1}). No external % SOURCE QUOTE.
  For each chart index \(i \in \{0, 1\} = \mathrm{Fin}\,2\), the chart-\(i\) scheme morphism
  \[ \mathtt{gmScalingP1\_chart}\,\bar k\,i \colon
     (\mathtt{gmScalingP1\_cover}\,\bar k).X\,i
     \longrightarrow \mathtt{ProjectiveLineBarScheme}\,\bar k \]
  realising \(\sigma_\times = \mathtt{gmScalingP1}\) (\cref{def:gaTranslationP1}) on the \(i\)-th
  chart of the pullback cover \(\mathtt{gmScalingP1\_cover}\) (\cref{def:gmscaling_cover}). On
  chart \(1\) (target affine chart \(D_+(X_1)\), coordinate \(u = X_0/X_1\)) the morphism realises
  \((x, \lambda) \mapsto \lambda x\) as the ring-level map \(u \mapsto u \otimes \lambda\) in
  \(\mathrm{Localization.Away}\,t \otimes_{\bar k} \mathtt{GmRing}\,\bar k\) (here \(t\) is the
  \(\mathrm{MvPolynomial}\,\{*\}\,\bar k\) generator); on chart \(0\) (target affine chart
  \(D_+(X_0)\), coordinate \(t = X_1/X_0\)) it realises \(1/(\lambda x) = u/\lambda\) as the ring-level
  map \(t \mapsto t \otimes \lambda^{-1}\). The two chart-side ring maps are packaged in Lean as
  \(\mathtt{gmScalingP1\_chart0\_ringMap}\) and \(\mathtt{gmScalingP1\_chart1\_ringMap}\) (both
  axiom-clean iter-171).
  \medskip
  \noindent\textbf{Construction recipe.} The scheme morphism is assembled from the chart-side
  ring map by the standard ``ring map \(\Rightarrow\) scheme map between \(\Spec\)'s of relative tensor
  products'' bridge, applied to the affine carrier of each cover chart. Concretely the chain is
  \(\mathtt{pullbackSpecIso}\) (identifying \((\mathtt{gmScalingP1\_cover})_i\) with \(\Spec\) of an
  \(\bar k\)-tensor product of the chart-\(i\) chart ring with \(\mathtt{GmRing}\)), followed by
  \(\Spec.\mathrm{map}\) of the ring map above, followed by the chart-ring iso of
  \cref{def:proj_chart_ring_iso} (which identifies \(\mathrm{Away}\,\mathcal A\,X_i\) with the
  affine chart ring \(\bar k[u]\) at the \(\mathrm{MvPolynomial}\,\{*\}\,\bar k\)-level), and then
  \(\mathrm{Proj.away}\,\iota\) to land in \(\mathbb P^1.\mathrm{left}\). The substantive step that is
  \emph{not} algebraic (the chart-ring iso \cref{def:proj_chart_ring_iso} is already axiom-clean
  via the iter-172 \cref{lem:mvPoly_to_homogeneousLocalization_away_surjective}) is the
  source-side \(\mathtt{pullbackSpecIso}\) identification, which navigates the abstract
  \(\mathtt{pullback}\,(\mathtt{pullback.fst}\,\mathbb P^1.\mathrm{hom}\,\mathbb G_m.\mathrm{hom})\,
  (\mathrm{Proj.away}\,\iota\,\mathcal A\,X_i)\) into the \(\Spec\)-of-tensor-product form. The full
  recipe is recorded in the iter-173 analogist note \(\mathtt{analogies/chart{\text-}bridge.md}\)
  (in flight), and the prior 6-step decomposition in \(\mathtt{analogies/gmscaling{\text-}deep.md}\).
\end{definition}
\begin{lemma}
  
\leanok
[Cocycle agreement for \(\mathtt{gmScalingP1\_chart}\) on pairwise overlaps]
  \label{lem:gmscaling_chart_agreement}

  \lean{AlgebraicGeometry.gmScalingP1_chart_agreement}
  \uses{def:gmscaling_chart, def:gmscaling_cover, lem:pullback_map_fst_proj, lem:pullback_map_snd_proj, thm:IsClosedImmersion_lift_iff_range_subset}
  % Archon-original Lean scaffold for the iter-171 chart-glue body of \(\sigma_\times\)
  % (\cref{def:gaTranslationP1}). No external % SOURCE QUOTE.
  For every pair of chart indices \(x, y \in \mathrm{Fin}\,2\) the two restrictions of
  \(\mathtt{gmScalingP1\_chart}\,\bar k\) (\cref{def:gmscaling_chart}) to the pairwise pullback
  \(\mathrm{pullback}\,((\mathtt{gmScalingP1\_cover})_x.f)\,((\mathtt{gmScalingP1\_cover})_y.f)\) agree:
  \[ \mathrm{pullback.fst}\,((\mathtt{gmScalingP1\_cover})_x.f)\,((\mathtt{gmScalingP1\_cover})_y.f)
     \fatsemi \mathtt{gmScalingP1\_chart}\,\bar k\,x
     \ =\
     \mathrm{pullback.snd}\,((\mathtt{gmScalingP1\_cover})_x.f)\,((\mathtt{gmScalingP1\_cover})_y.f)
     \fatsemi \mathtt{gmScalingP1\_chart}\,\bar k\,y. \]
  This is the cocycle / overlap-coherence equation that \(\mathtt{Scheme.Cover.glueMorphisms}\)
  consumes in order to assemble the chart-wise pieces \(\mathtt{gmScalingP1\_chart}\,\bar k\,i\) into
  a single morphism on the total \((\mathbb P^1 \otimes \mathbb G_m).\mathrm{left}\). The diagonal
  cases \((x, y) = (0, 0)\) and \((x, y) = (1, 1)\) are immediate via
  \(\mathtt{pullback.condition}\); the substantive content is the cross case \((0, 1)\) (and its
  symmetric \((1, 0)\)), broken down below.
  \medskip
  % NOTE (iter-186 writer, gmscaling-chart-agreement-expansion): two recurring Mathlib /
  % ergonomic gaps have repeatedly blocked closure of this lemma and are recorded once
  % here for the prover.
  % (i) The product-stability instances `gm_geomIrred` and `projGm_isReduced` rest on
  % `Algebra.TensorProduct.isDomain_of_isAlgClosed_left` (or any bridge "factor algebraically
  % closed + other a domain ⟹ tensor a domain"), which is absent at commit b80f227. Both
  % stay as typed sorries until upstream Mathlib lands the bridge or a project shim via
  % `Function.Injective.isDomain` on the chain `algebraMap (kbar ⊗ A) (kbar ⊗ A)` is
  % contributed; see \cref{lem:gmscaling_chart_agreement}'s pickup path (III.c) and the
  % stability-instance mini-blocks at the end of this subsection.
  % (ii) The current `gmScalingP1_cover_intersection_X_iso` is constructed in tactic mode
  % (`refine ≪≫ ?_; refine; exact`), so its elaboration does not display its
  % `Iso.trans`-spine syntactically; the canonical Mathlib simp chain
  % `Iso.trans_inv` / `Iso.inv_comp_eq` / `pullback.congrHom_inv` / `asIso_inv`
  % therefore fails to fire ("unused" simp arguments). This is a project ergonomic gap,
  % not a Mathlib lemma gap, addressed by pickup path (III.a).
  \noindent\textbf{(I) Setup and structural lift.} Restricting the indices to the cross case
  \((x, y) = (0, 1)\), the two charts being compared are
  \(\mathtt{gmScalingP1\_chart}\,\bar k\,0\) over the cover patch
  \((\mathtt{gmScalingP1\_cover}\,\bar k).X\,0 = D_+(X_0) \times \mathbb G_m\) (containing
  \(0 \in \mathbb A^1\)) and \(\mathtt{gmScalingP1\_chart}\,\bar k\,1\) over
  \((\mathtt{gmScalingP1\_cover}\,\bar k).X\,1 = D_+(X_1) \times \mathbb G_m\) (containing
  \(\infty\)). Their pairwise pullback intersection sits over the basic-open intersection
  \(D_+(X_0 \cdot X_1) \times \mathbb G_m\), and the project records an explicit iso
  identifying the intersection scheme with an affine \(\Spec\):
  \[ \mathtt{gmScalingP1\_cover\_intersection\_X\_iso}\,\bar k \colon
     \mathrm{pullback}\,((\mathtt{gmScalingP1\_cover}\,\bar k)_0.f)\,
       ((\mathtt{gmScalingP1\_cover}\,\bar k)_1.f)
     \xrightarrow{\ \sim\ }
     \Spec\bigl(\mathrm{Away}\,\mathcal A\,(X_0 \cdot X_1) \otimes_{\bar k}
       \mathtt{GmRing}\,\bar k\bigr). \]
  The cleanest entry point is to post-compose both sides of the cocycle with the iso's
  inverse and cancel the resulting epimorphism --- the
  \(\mathtt{cancel\_epi}\,(\mathtt{gmScalingP1\_cover\_intersection\_X\_iso}\,\bar k).\mathrm{inv}\)
  move. This lifts the agreement to a morphism equality between two arrows
  \(\Spec\bigl(\mathrm{Away}\,\mathcal A\,(X_0 \cdot X_1) \otimes_{\bar k}
     \mathtt{GmRing}\,\bar k\bigr) \rightrightarrows \overline{\mathbb P^1}\), where both
  sides should factor through the merged-generator chart
  \(\mathrm{Proj.awayι}\,(X_0 \cdot X_1)\). The cocycle then reduces to an equality of two
  ring maps into \(\mathrm{Away}\,\mathcal A\,(X_0 \cdot X_1) \otimes_{\bar k}
  \mathtt{GmRing}\,\bar k\).
  \medskip
  \noindent\textbf{(II) Why the canonical simp chain is blocked.} After
  \(\mathtt{cancel\_epi}\) the goal contains the inverse of
  \(\mathtt{gmScalingP1\_cover\_intersection\_X\_iso}\) composed with each pullback
  projection. In principle the Mathlib-canonical simp set
  \(\{\mathtt{Iso.trans\_inv},\, \mathtt{Iso.inv\_comp\_eq},\,
  \mathtt{pullback.congrHom\_inv},\, \mathtt{asIso\_inv},\,
  \mathtt{pullbackRightPullbackFstIso\_inv\_*},\,
  \mathtt{pullbackLeftPullbackSndIso\_inv\_*},\,
  \mathtt{pullbackSymmetry\_inv\_comp\_*},\,
  \mathtt{Proj.pullbackAwayιIso\_inv\_*}\}\)
  together with the project-side simp helpers
  \(\mathtt{pullback\_map\_fst\_proj}\) / \(\mathtt{pullback\_map\_snd\_proj}\)
  (\cref{lem:pullback_map_fst_proj}, \cref{lem:pullback_map_snd_proj}; iter-184 Recipe~1)
  would collapse the iso projections to a single \(\Spec.\mathrm{map}\) of a known ring
  map. Empirically, however, \(\mathtt{gmScalingP1\_cover\_intersection\_X\_iso}\) is
  produced by a tactic-mode \(\mathtt{refine\ \ldots\ \fatsemi\ \ldots\ \fatsemi\ ?\_}\)
  construction whose elaboration does not present its
  \(\mathtt{Iso.trans}\)-spine syntactically; consequently \emph{none} of the canonical
  inv-projection simp lemmas matches the elaborated form, and the iter-181 through
  iter-185 attempts at closure via \(\mathtt{simp [Iso.trans\_inv, \ldots]}\) report every
  simp argument ``unused''. This is the precise diagnostic recorded in
  \(\mathtt{analogies/gmscaling{\text-}projection{\text-}idiom.md}\) (iter-184 analogist
  consult, which supersedes the now-stale \(\mathtt{analogies/chart{\text-}bridge.md}\)
  cited in earlier iterations of this proof sketch).
  \medskip
  \noindent\textbf{(III) Three concrete pickup paths, with iter-187 status.} Three
  orthogonal strategies have been considered for closing this lemma. As of iter-187 the
  first two are off the table for the reasons recorded below, and only the third
  (III.c) --- the \emph{separated-locus alternative} --- is the live route. The
  three-way enumeration is preserved so that future iterations have an audit trail of
  what has and has not been tried.
  \begin{description}
    \item[(III.a) Refactor the iso to term-mode \(\fatsemi\)-spine ---
      \textbf{BLOCKED at Mathlib commit \texttt{b80f227}.}]
      The original idea was to rewrite
      \(\mathtt{gmScalingP1\_cover\_intersection\_X\_iso}\) as a single
      \(\mathtt{Iso.trans}\) (\(\fatsemi\)-)chain of named Mathlib isos
      (\(\mathtt{pullbackRightPullbackFstIso}\),
      \(\mathtt{pullback.congrHom}\),
      \(\mathtt{pullbackSymmetry}\),
      \(\mathtt{Proj.pullbackAwayιIso}\), and the final
      \(\mathtt{pullbackSpecIso}\)), with \emph{no} intermediate
      \(\mathtt{refine \ldots ?\_}\) holes discharged in tactic mode, so that
      \(\mathtt{Iso.trans\_inv}\) and the canonical inv-projection simp lemmas would
      unfold the chain link-by-link.
      \emph{Why this is blocked.} The shim
      \(\mathtt{Algebra.TensorProduct.isDomain\_of\_isAlgClosed\_left}\) (the bridge
      ``factor algebraically closed \(+\) other a domain \(\Rightarrow\) tensor a
      domain'') on which this path's product-stability instances rest is \emph{not
      present at the pinned Mathlib commit} \texttt{b80f227}. Independently, the
      iter-186 Lane B prover empirically verified that even with the term-mode
      refactor in place, the cocycle's residual simp normal form cannot be reduced by
      Mathlib's existing \(\mathtt{pullbackRightPullbackFstIso\_inv\_*}\) family: those
      simp lemmas require the projection to be \emph{syntactically displayed} at the
      top level, whereas in the elaborated cocycle the projections are buried inside
      outer \(\mathtt{pullback.map}\) constructs (\(\mathtt{simp}\) reports ``made no
      progress''). Both obstacles are external to the project and require upstream
      Mathlib motion. The route is therefore permanently blocked at this commit
      pending Mathlib upstream.

    \item[(III.b) Package named projection lemmas via
      \(\mathtt{pullback.hom\_ext}\) --- \textbf{DESCOPED iter-187.}]
      The fallback was, without refactoring the iso, to prove two
      \(\mathtt{@[reassoc (attr := simp)]}\) lemmas
      \(\mathtt{gmScalingP1\_cover\_intersection\_X\_iso\_inv\_fst}\) /
      \(\mathtt{\_inv\_snd}\) whose body is a hand-written
      \(\mathtt{pullback.hom\_ext}\) against the relevant pullback projection of the
      intersection. The universal property of the pullback would force each
      projection's value diagram-wise, sidestepping the opaque elaborated iso form
      because \(\mathtt{pullback.hom\_ext}\) does not require syntactic unfolding.
      \emph{Why this is descoped this iteration.} The iter-186 Lane B prover
      empirically determined that the (III.a) refactor (term-mode iso) does remove
      the elaboration-shape obstacle on the outer iso, but the simp coverage gap
      \emph{persists in mid-chain}: the \(\mathtt{pullback.map} \fatsemi
      \mathtt{pullbackRightPullbackFstIso}.\mathtt{inv}\) adjacency that the cocycle
      surfaces is not Mathlib-canonical, so even with the named projection lemmas
      built, the residual goal does not collapse without further Mathlib motion. The
      iter-185 \texttt{STRATEGY.md} \texttt{Open Q} ``failure-mode trigger'' fired in
      iter-186, and the iter-187 progress-critic returned a \texttt{STUCK +
      OVER\_BUDGET} verdict against this lemma (sorry count flat 4/4/4/4 across
      iter-183 to iter-186, four consecutive iterations of pivot/recipe language with
      no sorry closure, mandatory-decrement gate missed iter-186). Accordingly this
      path is descoped for iter-187 to focus all Lane B effort on the alternative
      (III.c); it may be revisited only if (III.c) itself surfaces unexpected
      blockers.

    \item[(III.c) Separated-locus alternative ---
      \textbf{iter-187 MANDATORY PIVOT} \emph{(per the iter-185
      \texttt{STRATEGY.md} \texttt{Open Q} failure-mode trigger, the iter-187
      progress-critic \texttt{STUCK} verdict, and the iter-187 blueprint-reviewer
      MF-2 directive).}]
      The chart-agreement equation
      \(\mathrm{pullback.fst} \fatsemi \mathtt{gmScalingP1\_chart}\,\bar k\,0 =
      \mathrm{pullback.snd} \fatsemi \mathtt{gmScalingP1\_chart}\,\bar k\,1\)
      on the cover intersection
      \(\mathrm{pullback}\,((\mathtt{gmScalingP1\_cover})_0.f)\,
      ((\mathtt{gmScalingP1\_cover})_1.f)\) is the statement that two morphisms
      \(\mathrm{chart}_0, \mathrm{chart}_1 \colon \mathrm{(intersection)} \to
      \overline{\mathbb P^1}\) agree on that intersection scheme. \emph{The key
      observation is that \(\overline{\mathbb P^1}\) is separated over \(\bar k\)},
      so the diagonal
      \(\Delta_{\overline{\mathbb P^1}} \colon
        \overline{\mathbb P^1} \longrightarrow
        \overline{\mathbb P^1} \times_{\bar k} \overline{\mathbb P^1}\)
      is a closed immersion. Because \(\overline{\mathbb P^1} = \Proj(\bar k[X_0, X_1])\)
      is the \(\Proj\) of a graded \(\bar k\)-algebra, its separatedness over \(\Spec
      \bar k\) is the standard Stacks fact ``\(\Proj(S)\) is separated'' (Stacks tag
      \texttt{01KU} / lemma \texttt{lemma-proj-separated}); Mathlib supplies the
      categorical packaging via \(\mathtt{IsSeparated}\) and
      \(\mathtt{IsSeparated.diagonal\_isClosedImmersion}\).
      \emph{Bypass.} Reformulate the cocycle as: the pair-morphism
      \((\mathrm{chart}_0, \mathrm{chart}_1) \colon \mathrm{(intersection)}
      \longrightarrow \overline{\mathbb P^1} \times_{\bar k} \overline{\mathbb P^1}\)
      (built via \(\mathtt{CategoryTheory.Limits.prod.lift}\) on
      \(\mathrm{chart}_0 := \mathrm{pullback.fst} \fatsemi \mathtt{gmScalingP1\_chart}\,
      \bar k\,0\) and
      \(\mathrm{chart}_1 := \mathrm{pullback.snd} \fatsemi \mathtt{gmScalingP1\_chart}\,
      \bar k\,1\)) \emph{factors through the diagonal
      \(\Delta_{\overline{\mathbb P^1}}\)}. By Mathlib's
      \(\mathtt{IsClosedImmersion.lift\_iff\_range\_subset}\) (the universal property
      of a closed immersion: a morphism factors through a closed immersion iff its
      set-theoretic image is contained in the image of the immersion), the
      factorization is equivalent to the set-level containment of the cocycle's image
      in the image of \(\Delta\), which on points is exactly the chart equality on
      the \emph{closed points} of the intersection. The closed-point check is
      discrete (each closed point of the intersection lies over a single \(\bar k\)-
      point of \(D_+(X_0 \cdot X_1) \times \mathbb G_m\) and the two chart maps both
      evaluate to the same \(\bar k\)-point of \(\overline{\mathbb P^1}\)
      corresponding to that \((x, \lambda)\)), and bypasses the simp coverage gap
      that obstructs (III.a)/(III.b) entirely.
      \emph{Concrete prover recipe.} The path closes in four steps:
      \begin{enumerate}
        \item Fix the diagonal
          \(\Delta_{\overline{\mathbb P^1}} \colon \overline{\mathbb P^1}
            \longrightarrow \overline{\mathbb P^1} \times_{\bar k}
            \overline{\mathbb P^1}\)
          via Mathlib's \(\mathtt{CategoryTheory.Limits.prod.lift}\,(\mathbf
          1_{\overline{\mathbb P^1}})\,(\mathbf 1_{\overline{\mathbb P^1}})\), and
          obtain the closed-immersion witness from
          \(\mathtt{IsSeparated.diagonal\_isClosedImmersion}\) applied to the
          separated structure morphism
          \(\overline{\mathbb P^1}.\mathrm{hom} \colon \overline{\mathbb P^1} \to
          \Spec \bar k\).
        \item Build the pair-morphism
          \((\mathrm{chart}_0, \mathrm{chart}_1) \colon
            \mathrm{pullback}\,((\mathtt{gmScalingP1\_cover})_0.f)\,
              ((\mathtt{gmScalingP1\_cover})_1.f)
            \longrightarrow
            \overline{\mathbb P^1} \times_{\bar k} \overline{\mathbb P^1}\)
          via \(\mathtt{prod.lift}\) on the two named chart restrictions above.
        \item Show this pair factors through \(\Delta_{\overline{\mathbb P^1}}\)
          via the universal property of the intersection's pullback structure
          combined with \(\mathtt{IsClosedImmersion.lift}\) (kernel-inequality
          form; the \(\mathtt{\_iff\_range\_subset}\) variant is NOT in Mathlib
          at \texttt{b80f227} per iter-188 verification, so this step requires
          either upstream contribution or a project-side bridge plus
          \(\mathtt{IsReduced}\) on the intersection): produce the section
          \(\mathrm{(intersection)} \to \overline{\mathbb P^1}\) whose post-
          composition with \(\Delta\) equals \((\mathrm{chart}_0,
          \mathrm{chart}_1)\).
        \item The factorization reduces, by composing with the two projections
          \(\mathtt{prod.fst}\) and \(\mathtt{prod.snd}\) out of the codomain
          \(\overline{\mathbb P^1} \times_{\bar k} \overline{\mathbb P^1}\), to
          two scheme-equalities on the components --- each of which is a direct
          chart-bridge equality on \((\mathtt{gmScalingP1\_cover})_i\)'s structure
          (i.e.\ the \emph{single-chart} content, not the full cocycle), discharged
          using the same chart-ring iso \cref{def:proj_chart_ring_iso} and
          \(\mathtt{pullbackSpecIso}\) bridge that \cref{def:gmscaling_chart} uses.
          The chart-\(1\) instance of the per-component equality threads
          through the named \(\mathrm{Proj.appIso}\) evaluation
          \cref{lem:iotaGm_chart1_appIso_eval}, which packages the basic-open
          ring-map identification
          \(\mathtt{isLocElem} \mapsto [X_0/X_1] \mapsto 1\) at the global-sections
          level so the prover can target it directly without re-deriving the
          \((\mathrm{Proj.appIso}).\mathrm{inv}\) image of \(\mathtt{isLocElem}\)
          inline (the iter-188 through iter-191 budget-wall step). The chart-\(0\)
          instance is the symmetric statement with roles of \(X_0, X_1\)
          interchanged.
      \end{enumerate}
      \emph{Substrate hooks.} The categorical substrate is present at the
      pinned Mathlib commit \texttt{b80f227}:
      \(\mathtt{IsSeparated.diagonal\_isClosedImmersion}\),
      \(\mathtt{IsClosedImmersion.lift}\) (kernel-inequality form),
      \(\mathtt{CategoryTheory.Limits.prod.lift}\),
      \(\mathtt{pullback.lift}\), together with the standard
      \(\mathtt{IsSeparated}\) instance for \(\Proj\) of a graded ring (Stacks
      \texttt{01KU} / Hartshorne~II.4.6 diagonal-is-closed characterisation).
      % NOTE (iter-188 review): the \texttt{\_iff\_range\_subset} variant of
      % \texttt{IsClosedImmersion.lift} cited in the (III.c) recipe is NOT present
      % at b80f227 (verified iter-188 via lean\_leansearch on
      % \enquote{morphism factors through closed immersion iff range subset}
      % and \enquote{IsClosedImmersion lift range subset}; both return only the
      % kernel-inequality form). Closure requires either a Mathlib upstream
      % contribution of \texttt{IsClosedImmersion.lift\_iff\_range\_subset} or a
      % project-side build that also threads
      % \texttt{Algebra.TensorProduct.isDomain\_of\_isAlgClosed\_left}-style
      % reducedness of the intersection. See
      % \texttt{.archon/task\_results/AlgebraicJacobian\_Genus0BaseObjects\_GmScaling.lean.md}
      % \S\enquote{Mathlib lemmas confirmed NOT SHIPPED}.
      \emph{Cost estimate (post iter-188 falsification).} \(\sim 150\)--\(200\)
      LOC over 3--5 iters (Option B: project-side substrate including
      \(\mathtt{lift\_iff\_range\_subset}\) plus \(\mathtt{IsReduced}\) on the
      intersection); Option A (Mathlib upstream) is unbounded in time.
  \end{description}
  \medskip
  \noindent\textbf{(IV) Substantive ring-level residual.} All three pickup paths reduce,
  after their structural setup, to the same ring-level identity. On the basic-open
  intersection \(D_+(X_0 \cdot X_1) \subseteq \overline{\mathbb P^1}\) both chart
  coordinates are units (the chart-\(0\) coordinate \(t = X_1/X_0\) and the chart-\(1\)
  coordinate \(u = X_0/X_1\) are reciprocals: \(t \cdot u = 1\)), and the agreement
  reduces to the ring-level identity
  \[ \lambda \cdot u \ =\ \frac{1}{t} \cdot \lambda \quad
     \text{in } \mathrm{Localization.Away}\,t \otimes_{\bar k} \mathtt{GmRing}\,\bar k, \]
  i.e.\ to the unit-of-the-localisation identity \(t \cdot u = 1\)
  (\(\mathtt{IsLocalization.Away.mul\_invSelf}\)) together with commutativity of the
  tensor product. The proof routes through the chart-ring iso
  \cref{def:proj_chart_ring_iso} and the same \(\mathtt{pullbackSpecIso}\) bridge that
  \cref{def:gmscaling_chart} uses; consequently \cref{lem:gmscaling_chart_agreement}
  remains gated on \cref{def:gmscaling_chart} acquiring its concrete body \emph{and} on
  one of the (III.a)--(III.c) structural setups landing. See
  \(\mathtt{analogies/gmscaling{\text-}projection{\text-}idiom.md}\) (iter-184) for the
  full Recipe~1/2/3 decomposition and empirical simp-probe verification.
\end{lemma}
  
\begin{lemma}
\leanok
[Per-chart bridge \(\mathtt{gmScalingP1\_chart}\,i\,\fatsemi\,\mathbb P^1.\mathrm{hom}\)]
  \label{lem:gmscaling_chart_PLB_eq}

  \lean{AlgebraicGeometry.gmScalingP1_chart_PLB_eq}
  \uses{def:gmscaling_chart, def:gmscaling_cover, lem:chart_ring_iso_preserves_algebraMap}
  % Archon-original Lean encoding; the per-chart certificate that
  % \cref{lem:gmscaling_over_coherence} consumes after \(\mathtt{Scheme.Cover.hom\_ext}\) and
  % \(\mathtt{ι\_glueMorphisms\_assoc}\) reduce the over-side coherence to per-chart equations.
  % No external % SOURCE QUOTE: this is bespoke to the project's iter-171 chart-glue split of
  % \(\sigma_\times = \mathtt{gmScalingP1}\).
  For each chart index \(i \in \{0, 1\} = \mathrm{Fin}\,2\), the chart-\(i\) scheme morphism
  \(\mathtt{gmScalingP1\_chart}\,\bar k\,i\) (\cref{def:gmscaling_chart}) intertwines the
  structure morphism of the cover chart \((\mathtt{gmScalingP1\_cover}\,\bar k).f\,i\)
  (\cref{def:gmscaling_cover}) with the structure morphism of \(\mathbb P^1\):
  \[ \mathtt{gmScalingP1\_chart}\,\bar k\,i
     \,\fatsemi\,
     \mathtt{ProjectiveLineBar}\,\bar k.\mathrm{hom}
     \ =\
     (\mathtt{gmScalingP1\_cover}\,\bar k).f\,i
     \,\fatsemi\,
     \bigl((\mathtt{ProjectiveLineBar}\,\bar k) \otimes \mathtt{Gm}\,\bar k\bigr).\mathrm{hom}. \]
  This is the per-chart certificate that, after \(\mathtt{Scheme.Cover.hom\_ext}\) followed by
  \(\mathtt{ι\_glueMorphisms\_assoc}\), closes the over-side coherence
  \(\mathtt{gmScalingP1\_over\_coherence}\) (\cref{lem:gmscaling_over_coherence}) by which the
  glued total morphism \(\sigma_\times = \mathtt{gmScalingP1}\) (\cref{def:gaTranslationP1}) is
  promoted from a bare \(\mathbf{Sch}\)-level map to a morphism in \(\mathrm{Over}\,(\Spec \bar k)\).
  \medskip
  \noindent\textbf{Three-step reduction.} The proof has three steps. All three route through
  the project-side chart-bridge \(\mathtt{awayι\_comp\_PLB\_hom}\)
  (\(\mathtt{Genus0BaseObjects.lean}\), post-iter-175 split moves to
  \(\mathtt{Genus0BaseObjects/ChartBridge.lean}\)), which identifies \(\mathrm{Proj.awayι}
  \,\fatsemi\, \mathbb P^1.\mathrm{hom}\) with \(\Spec.\mathrm{map}\,(\mathrm{algebraMap}\,
  \bar k\,(\mathrm{Away}\,\mathcal A\,X_i))\) for the degree-\(1\) homogeneous generators
  \(X_i\).
  \begin{description}
    \item[(A) Collapse \(\mathrm{Proj.awayι} \fatsemi \mathbb P^1.\mathrm{hom}\).] Apply the
      chart-bridge \(\mathtt{awayι\_comp\_PLB\_hom}\) to rewrite the tail of the chart-\(i\)
      morphism, replacing \(\mathrm{Proj.awayι} \,\fatsemi\, \mathbb P^1.\mathrm{hom}\) with
      \(\Spec.\mathrm{map}\,(\mathrm{algebraMap}\,\bar k\,(\mathrm{Away}\,\mathcal A\,X_i))\).
    \item[(B) Merge the \(\Spec.\mathrm{map}\)'s.] Use \(\Spec.\mathrm{map\_comp}\) to merge the
      two consecutive \(\Spec.\mathrm{map}\)'s on the chart-\(i\) side, then apply the
      chart-ring iso's \(\bar k\)-algebra preservation
      (\cref{lem:chart_ring_iso_preserves_algebraMap}) together with
      \(\mathrm{MvPolynomial.eval}_2\mathrm{Hom}\_C\) to identify the merged ring map with
      \(\mathrm{algebraMap}\,\bar k\,(\mathrm{Away}\,\mathcal A\,X_i \otimes_{\bar k}
      \mathtt{GmRing}\,\bar k)\).
    \item[(C) Chase the source-side pullback isos.] Unfold the source-side iso
      \(\mathtt{gmScalingP1\_cover\_X\_iso}\,\bar k\,i\) (identifying the cover chart with
      \(\Spec\) of an \(\bar k\)-tensor product of chart-\(i\) and \(\mathtt{GmRing}\)). Then use
      Mathlib's \(\mathtt{pullbackSpecIso\_hom\_base}\),
      \(\mathtt{pullbackRightPullbackFstIso\_hom\_fst}\), and
      \(\mathtt{pullbackSymmetry\_hom\_comp\_fst}\) to chase the source side back to
      \((\mathtt{gmScalingP1\_cover}\,\bar k).f\,i\) composed with the structure morphism of
      the relative tensor product. The final identification with the RHS is via Mathlib's
      \(\mathtt{Over.tensorObj\_hom}\) and \(\mathtt{pullback.condition}\).
  \end{description}
  \emph{Status (axiom-clean since iter-180; updated iter-185 review).} Steps (A) and (B)
  are axiom-clean as of iter-174. Step (C) was closed kernel-clean
  \(\{\mathtt{propext}, \mathtt{Classical.choice}, \mathtt{Quot.sound}\}\) in iter-180 via
  the \(\mathtt{set\_option}\,\mathtt{backward.isDefEq.respectTransparency\,false}\) recipe
  (Decision-4 \texttt{ALIGN\_WITH\_MATHLIB} from the \texttt{pullbackspeciso-bypass} analogist
  consult). The earlier syntactic \(\mathrm{Fin}\)-literal mismatch and the
  \(\mathtt{Algebra.compHom}\)-driven instance synthesis chain were both retired by that
  recipe; both iter-177 \texttt{TEMP} project axioms were dropped.
  
\end{lemma}
\begin{lemma}
\leanok
[Over-structure coherence for the glued total morphism]
  \label{lem:gmscaling_over_coherence}

  \lean{AlgebraicGeometry.gmScalingP1_over_coherence}
  \uses{def:gmscaling_chart, def:gmscaling_cover, lem:gmscaling_chart_agreement, lem:gmscaling_chart_PLB_eq}
  % Archon-original Lean scaffold for the iter-171 chart-glue body of \(\sigma_\times\)
  % (\cref{def:gaTranslationP1}). No external % SOURCE QUOTE.
  The total morphism obtained by \(\mathtt{Scheme.Cover.glueMorphisms}\) from
  \(\mathtt{gmScalingP1\_chart}\,\bar k\) (\cref{def:gmscaling_chart}) and its cocycle
  \(\mathtt{gmScalingP1\_chart\_agreement}\,\bar k\) (\cref{lem:gmscaling_chart_agreement}) intertwines
  the structure morphisms to \(\Spec \bar k\):
  \[ (\mathtt{gmScalingP1\_cover}\,\bar k).\mathtt{glueMorphisms}\,
       (\mathtt{gmScalingP1\_chart}\,\bar k)\,
       (\mathtt{gmScalingP1\_chart\_agreement}\,\bar k)
     \fatsemi \mathtt{ProjectiveLineBar}\,\bar k.\mathrm{hom}
     \ =\
     ((\mathtt{ProjectiveLineBar}\,\bar k) \otimes \mathtt{Gm}\,\bar k).\mathrm{hom}. \]
  This is the \(\mathrm{Over}\,(\Spec \bar k)\)-side commutation that promotes the bare scheme map
  on \((\mathbb P^1 \otimes \mathbb G_m).\mathrm{left}\) to a genuine morphism in
  \(\mathrm{Over}\,(\Spec \bar k)\), supplying the \(\mathtt{Over.homMk}\) second argument that the
  body of \(\sigma_\times = \mathtt{gmScalingP1}\) (\cref{def:gaTranslationP1}) consumes.
  \medskip
  \noindent\textbf{Reduction.} By \(\mathtt{Scheme.Cover.hom\_ext}\) for the cover
  \(\mathtt{gmScalingP1\_cover}\) (\cref{def:gmscaling_cover}), the equation is equivalent to its
  chart-wise restrictions: for each \(i \in \mathrm{Fin}\,2\), composing first with
  \((\mathtt{gmScalingP1\_cover})_i.f \colon (\mathtt{gmScalingP1\_cover})_i \to (\mathbb P^1 \otimes
  \mathbb G_m).\mathrm{left}\) on the left, both sides land in \(\Spec \bar k\). Agreement there is
  automatic from the way \(\mathtt{gmScalingP1\_chart}\,\bar k\,i\) is built: it factors through
  \(\Spec.\mathrm{map}\) of the \(\bar k\)-algebra map
  \(\mathrm{algebraMap}\,\bar k\,(\mathrm{Away}\,\mathcal A\,X_i \otimes_{\bar k} \mathtt{GmRing})\),
  so its composition with the structure map down to \(\Spec \bar k\) matches the structure map of
  the cover chart. The argument is therefore mechanical given that \cref{def:gmscaling_chart}
  has a concrete body; like the cocycle, \cref{lem:gmscaling_over_coherence} is gated on
  
  \cref{def:gmscaling_chart}. See \(\mathtt{analogies/chart{\text-}bridge.md}\) (iter-173 in flight)
  
  for the chart-wise step.
\end{lemma}
\begin{lemma}
\leanok
[The \(\mathbb G_m\)-scaling action fixes the origin \(0 \in \mathbb P^1\)]
  \label{lem:gmScaling_fixes_zero}

  \lean{AlgebraicGeometry.gmScalingP1_collapse_at_zero}
  \uses{def:gaTranslationP1, def:p1bar_zero}
  % NOTE (iter-171 refresh of iter-170): the Lean body of `gmScalingP1_collapse_at_zero`
  % remains a gated `sorry` — but the gate is now NARROWER: iter-171 landed the
  % `gmScalingP1` body SKELETON (concrete `Over.homMk + Scheme.Cover.glueMorphisms`),
  % so the residual gate is on the three internal scaffold helpers `gmScalingP1_chart{,
  % _agreement, _over_coherence}`. Once `gmScalingP1_chart` (chart-`i` scheme morphism)
  % has a concrete body, this collapse lemma reduces via chart-1 to the ring-level
  % identity `X 0 ↦ 0 ⊗ λ = 0` in `Localization.Away t ⊗ GmRing` directly.
  % Archon-original; the fact is elementary (σ_×(0, λ) = λ · 0 = 0). No external % SOURCE QUOTE
  % needed for the Lean encoding.
  Let \(\sigma_\times \colon \mathbb P^1 \times \mathbb G_m \to \mathbb P^1\) be the scaling action
  of \cref{def:gaTranslationP1}. Then the composite
  \[ \mathbb G_m \xrightarrow{\;(\mathrm{toUnit}_{\mathbb G_m} \fatsemi 0_{\mathbb P^1},\ \mathbf 1_{\mathbb G_m})\;}
     \mathbb P^1 \times \mathbb G_m \xrightarrow{\;\sigma_\times\;} \mathbb P^1 \]
  equals the constant morphism \(\mathrm{toUnit}_{\mathbb G_m} \fatsemi 0_{\mathbb P^1}\). That is,
  \(\sigma_\times(0, \lambda) = 0\) for every \(\lambda \in \mathbb G_m\).
\end{lemma}
\begin{proof}
  Chart-level. On the affine chart \(\mathbb A^1 \times \mathbb G_m\) (coordinate \(x\)), the
  morphism \(\sigma_\times\) is the polynomial map \((x, \lambda) \mapsto \lambda \cdot x\). Restricting
  to the \(\mathbb G_m\)-section \(x = 0\) gives the constant ring map
  \(\bar k[t, t^{-1}] \to \bar k\), \(\lambda \mapsto 0\); equivalently the scheme morphism factors
  
  through the closed point \(0 \in \mathbb P^1\). The agreement with the other chart \(u = 1/x\) near
  
  \(\infty\) is automatic (the section \(x = 0\) does not meet that chart). The named equation is the
  Lean encoding of this fact, with \(\mathrm{toUnit}_{\mathbb G_m}\) supplying the
  \(\Spec \bar k\)-factor and \(\mathbf 1_{\mathbb G_m}\) the parameter \(\lambda\).
\end{proof}
% NOTE (iter-186 writer, gmscaling-chart-agreement-expansion): mini-blocks below pin the
% four product-stability instances on \(\mathbb P^1 \otimes \mathbb G_m\) that
% \cref{prop:morphism_P1_to_AV_constant}'s Lean basepoint helper consumes (one local-finiteness
% instance, one reducedness instance, and the two geometric-irreducibility instances). They
% are exported from `AlgebraicJacobian/Genus0BaseObjects/GmScaling.lean`. The two
% sorry-bearing instances (`gm_geomIrred`, `projGm_isReduced`) carry the Mathlib gap recorded
% in \cref{lem:gmscaling_chart_agreement}'s iter-186 NOTE.
\begin{lemma}

\leanok
[\((\mathbb P^1 \otimes \mathbb G_m).\mathrm{hom}\) is locally of finite type]
  \label{lem:projGm_locallyOfFiniteType}
  \lean{AlgebraicGeometry.projGm_locallyOfFiniteType}
  \uses{def:genus0_base_objects}
  % Archon-original Lean export to Lane B (\cref{prop:morphism_P1_to_AV_constant}'s
  % basepoint helper). No external % SOURCE QUOTE: the result is the standard
  % composition-and-base-change closure of \(\mathtt{LocallyOfFiniteType}\).
  The structure morphism of the product
  \((\overline{\mathbb P^1} \otimes \mathbb G_m).\mathrm{hom}\) to \(\Spec \bar k\) is
  locally of finite type. The proof is the standard decomposition
  \(\mathtt{pullback.fst} \fatsemi \overline{\mathbb P^1}.\mathrm{hom}\) followed by
  Mathlib's composition closure \(\mathtt{locallyOfFiniteType\_comp}\) and base-change
  stability \(\mathtt{locallyOfFiniteType\_isStableUnderBaseChange}\). Both factor
  morphisms \(\overline{\mathbb P^1}.\mathrm{hom}\) and \(\mathbb G_m.\mathrm{hom}\) are
  locally of finite type over \(\bar k\) (the former as \(\Proj\) of a finitely generated
  graded algebra, the latter as a localization of a polynomial ring).
\end{lemma}
\begin{lemma}

\leanok
[\(\mathbb G_m \to \Spec \bar k\) is geometrically irreducible \emph{(Mathlib gap)}]
  \label{lem:gm_geomIrred}
  \lean{AlgebraicGeometry.gm_geomIrred}
  \uses{def:gm, thm:gmRing_tensor_homogeneousAway_isDomain, lem:affineLine_geomIrred}
  % Archon-original Lean scaffold (sorry-bearing instance). No external % SOURCE QUOTE.
  % The mathematical content is standard but the Mathlib bridge (see below) is absent at
  % commit b80f227.
  Over the algebraically closed base field \(\bar k\) (the typeclass
  \(\mathtt{[IsAlgClosed\ \bar k]}\) is in scope), the structure morphism
  \(\mathbb G_m.\mathrm{hom} \colon \mathbb G_m \to \Spec \bar k\) is geometrically
  irreducible. The mathematical content is elementary --- \(\mathbb G_m\) is
  \(\Spec\,\bar k[t, t^{-1}]\) with \(\bar k[t, t^{-1}]\) a domain over the algebraically
  closed \(\bar k\), and any base change \(\mathrm{Spec}\,A \otimes_{\bar k} \bar k[t, t^{-1}]\)
  is the localisation of the polynomial ring \(A[t]\) at the powers of \(t\), still a domain
  when \(A\) is. The Mathlib bridge ``\(\mathtt{IsAlgClosed} + \mathtt{IrreducibleSpace}\) on
  the fibres \(\Rightarrow \mathtt{GeometricallyIrreducible}\) at scheme level'' is the
  recorded gap: the standard route via
  \(\mathtt{Algebra.TensorProduct.isDomain\_of\_isAlgClosed\_left}\) (the bridge
  ``algebraically closed factor + other factor a domain \(\Rightarrow\) tensor is a
  domain'') is absent at commit \(\mathtt{b80f227}\) (CONFIRMED gap, iter-185
  lean-vs-blueprint-checker). Stays as a typed sorry until upstream lands the bridge or a
  project-side \(\sim 30\)--\(50\) LOC shim via
  \(\mathtt{Function.Injective.isDomain}\) on the chain
  \(\mathtt{algebraMap}\,(\bar k \otimes A)\,(\bar k \otimes A)\) is contributed.
\end{lemma}
\begin{lemma}

\leanok
[\((\mathbb P^1 \otimes \mathbb G_m).\mathrm{hom}\) is geometrically irreducible]
  \label{lem:projGm_geomIrred}
  \lean{AlgebraicGeometry.projGm_geomIrred}
  \uses{lem:gm_geomIrred, lem:projectiveLineBar_geomIrred}
  % Archon-original Lean export to Lane B. No external % SOURCE QUOTE: the result follows
  % from the factor-wise scaffolds via Mathlib's `GeometricallyIrreducible.comp`.
  Over the algebraically closed base \(\bar k\) (\(\mathtt{[IsAlgClosed\ \bar k]}\) in
  scope), the structure morphism
  \((\overline{\mathbb P^1} \otimes \mathbb G_m).\mathrm{hom}\) is geometrically
  irreducible. The proof composes the per-factor scaffolds via Mathlib's
  \(\mathtt{GeometricallyIrreducible.comp}\) (with \(\mathtt{UniversallyOpen}\) handled
  by smoothness of each factor): the product structure morphism unfolds as
  \(\mathtt{pullback.fst}\,\overline{\mathbb P^1}.\mathrm{hom}\,\mathbb G_m.\mathrm{hom}
   \,\fatsemi\, \overline{\mathbb P^1}.\mathrm{hom}\), the
  \(\mathtt{pullback.fst}\) factor is geometrically irreducible by base-change stability
  from \cref{lem:gm_geomIrred}, and the second factor is the
  geometric-irreducibility of \(\overline{\mathbb P^1}\) (already in scope). Axiom-clean
  contingent on the two factor scaffolds (\cref{lem:gm_geomIrred} carries the Mathlib
  gap).
\end{lemma}
\begin{lemma}

\leanok
[\((\mathbb P^1 \otimes \mathbb G_m).\mathrm{left}\) is reduced \emph{(Mathlib gap)}]
  \label{lem:projGm_isReduced}
  \lean{AlgebraicGeometry.projGm_isReduced}
  \uses{lem:projlinebar_isReduced, def:gm, thm:gmRing_tensor_homogeneousAway_isDomain}
  % Archon-original Lean scaffold (sorry-bearing instance). No external % SOURCE QUOTE.
  % Carries the same Mathlib gap as \cref{lem:gm_geomIrred} on the
  % `tensor-of-domains-over-field-is-domain` bridge.
  The scheme underlying \((\overline{\mathbb P^1} \otimes \mathbb G_m).\mathrm{left}\)
  is reduced. The textbook argument is a chart-local one: cover
  \(\overline{\mathbb P^1} \otimes \mathbb G_m\) by
  \(\Spec\,\bar k[t, \lambda, \lambda^{-1}]\) charts (the product of the affine cover of
  \(\overline{\mathbb P^1}\) and the affine \(\mathbb G_m\)); each chart ring is a
  localisation of \(\bar k[X_0, X_1] \otimes_{\bar k} \bar k[\lambda]\) at the
  non-zero-divisor \(\lambda\), hence a domain, hence reduced. The bridge
  \(\mathtt{HomogeneousLocalization.Away}\)-is-domain \(+\)
  \(\mathtt{TensorProduct}\)-of-domains-over-field-is-domain is the recorded Mathlib gap
  (\(\mathtt{Algebra.TensorProduct.isDomain\_of\_isAlgClosed\_left}\) absent at
  \(\mathtt{b80f227}\); CONFIRMED gap, iter-185 lean-vs-blueprint-checker). The
  Smooth-\(\Rightarrow\)-GeometricallyReduced shortcut is also unavailable at scheme
  level. Stays as a typed sorry until either upstream bridge lands or a project shim is
  contributed.
\end{lemma}
\begin{lemma}
[\(\mathrm{Hom}(\mathbb G_a, A) = 0\): a homomorphism from the additive group into an abelian variety is trivial]
  \label{lem:hom_Ga_to_av_trivial}

  \lean{AlgebraicGeometry.hom_Ga_to_av_trivial}
  \uses{lem:eq_comp_of_isAffine_of_properIntegral, lem:hom_additivity_over_product, def:gaTranslationP1}
  % SOURCE: [Milne, Abelian Varieties], Proposition 3.9, \S I.3, doc p.~19--20 (read from references/abelian-varieties.pdf, PDF pages 25--26)
  % SOURCE QUOTE: "PROPOSITION 3.9. Every rational map A¹ ⇢ A or P¹ ⇢ A is constant."
  % NOTE (iter-164): the verbatim Milne Prop 3.9 statement+proof quotes are the prior-session-verified
  % copies also carried on \cref{lem:hom_from_Ga_trivial} below. This host lacks a PDF renderer
  % (pdftoppm) so the quotes were NOT re-rendered from references/abelian-varieties.pdf this session;
  % they are reproduced from the verified in-tree copy. See the writer report.
  \textit{Source: Milne, Abelian Varieties, Proposition~3.9.}
  \textbf{Off the genus-\(0\) critical path (iter-164).} This \(\mathrm{Hom}(\mathbb G_a, A) = 0\)
  lemma is \emph{superseded by the \(\mathbb G_m\)-scaling shortcut}
  (\cref{prop:morphism_P1_to_AV_constant}, \cref{rmk:base_case_fourth_route}), which closes the
  base case without it. It is retained only as the classical Milne Proposition~3.9 (\(\mathbb
  G_a/\mathbb G_m\) incompatibility) argument feeding the alternative route
  \cref{lem:hom_from_Ga_trivial}; no block on the genus-\(0\) critical path depends on it.
  Let \(A\) be an abelian variety over the algebraically closed field \(\bar k\), and let
  \(\varphi \colon \mathbb G_a \to A\) be a homomorphism of group objects (so \(\varphi(0) = 0_A\) and
  \(\varphi(x + y) = \varphi(x) + \varphi(y)\)). Then \(\varphi\) is the trivial homomorphism: it
  factors through the terminal object as the constant morphism at the identity,
  \(\varphi = \mathtt{toUnit}\,\mathbb G_a \fatsemi \eta_A\).
  \medskip
  \noindent\textbf{Equivalent completeness phrasing.} Equivalently: a homomorphism from an
  \emph{affine connected} group object into a \emph{complete} group object is trivial. That is the
  abstract content; the proof below realises it characteristic-free.
\end{lemma}
% SOURCE QUOTE PROOF: "PROOF. According to (3.2), α extends to a regular map on the whole of A¹.
% After composing α with a translation, we may suppose that α(0)=0. Then α is a homomorphism,
% α(x+y)=α(x)+α(y); all x,y ∈ A¹(k^al)=k^al. But A¹−{0} is also a group variety, and similarly,
% α(xy)=α(x)+α(y)+c; all x,y ∈ A¹(k^al)=k^al. This is absurd, unless α is constant." [Milne,
% Prop 3.9, \S I.3, doc p.~19--20.]
\begin{proof}
  \textit{Source: Milne, Abelian Varieties, Proposition~3.9.}
  Write \(A\) additively (an abelian variety is commutative). We give the conceptual
  \emph{completeness} proof as the abstract argument, then record Milne's elementary \(\mathbb
  G_a/\mathbb G_m\) incompatibility, which applies whenever \(\varphi\) extends to a morphism on
  \(\mathbb P^1\) (as it does in the genus-\(0\) application).
  \textbf{Completeness proof.} The image \(\varphi(\mathbb G_a)\) is a connected subgroup of \(A\)
  (image of the connected \(\mathbb G_a\) under a homomorphism). As the image of a homomorphism of
  algebraic groups it is a \emph{closed} subgroup of \(A\), hence \emph{complete} (closed in the
  complete \(A\)); as a homomorphic image of the \emph{affine} \(\mathbb G_a\) it is moreover
  \emph{affine}. A connected variety that is simultaneously complete and affine over \(\bar k\) is a
  single point: its ring of global regular functions is a field finite over \(\bar k\), hence equal
  to \(\bar k\) (the \(\Gamma = \bar k\) collapse already packaged for proper integral sources in
  \cref{lem:eq_comp_of_isAffine_of_properIntegral}). Therefore \(\varphi(\mathbb G_a)\) is one point;
  being a homomorphism, \(\varphi\) sends \(0\) there, so the point is \(0_A\) and
  \(\varphi = \mathtt{toUnit}\,\mathbb G_a \fatsemi \eta_A\).
  \medskip
  \noindent\textbf{Milne's elementary \(\mathbb G_a/\mathbb G_m\) incompatibility (when \(\varphi\)
  extends to \(\mathbb P^1\)).} Suppose \(\varphi = f|_{\mathbb G_a}\) for a morphism
  \(f \colon \mathbb P^1 \to A\) (the genus-\(0\) situation). Set \(g := f - f(1)\), so \(g(1) = 0_A\), and
  form \(h_\times := \sigma_\times \fatsemi g \colon \mathbb P^1 \times \mathbb G_m \to A\) with the
  scaling companion \(\sigma_\times\) of \cref{def:gaTranslationP1}. Applying the additive
  decomposition \cref{lem:hom_additivity_over_product} with \(V = \mathbb P^1\) (complete),
  \(W = \mathbb G_m\), base points \(1 \in \mathbb P^1\) and \(1 \in \mathbb G_m\) (so
  \(h_\times(1, 1) = g(1) = 0_A\)), the axis restrictions are \(x \mapsto g(x)\) and
  \(\lambda \mapsto g(\lambda)\), whence \(g(\lambda x) = g(x) + g(\lambda)\) on \(\bar k^\times\);
  re-expanding \(g\) gives the multiplicative relation
  \(\varphi(xy) = \varphi(x) + \varphi(y) + c\) with \(c = -\varphi(1)\). Now \(\varphi\) is at once
  additive, \(\varphi(x + y) = \varphi(x) + \varphi(y)\), and ``multiplicative-additive'',
  \(\varphi(xy) = \varphi(x) + \varphi(y) + c\). Computing \(\varphi((x + x')\,y)\) in two ways for
  \(x, x' \in \bar k\), \(y \in \bar k^\times\),
  \[ \varphi((x+x')y) = \varphi(xy) + \varphi(x'y) = \varphi(x) + \varphi(x') + 2\varphi(y) + 2c, \]
  while also \(\varphi((x+x')y) = \varphi(x+x') + \varphi(y) + c = \varphi(x) + \varphi(x') +
  \varphi(y) + c\); equating gives \(\varphi(y) = -c = \varphi(1)\) for every \(y \in \bar k^\times\).
  Hence \(\varphi\) is constant on the dense \(\mathbb G_m \subset \mathbb G_a\), so constant on
  \(\mathbb G_a\) (density, \(A\) separated), with value \(\varphi(0) = 0_A\). This is Milne's argument
  verbatim and uses only the proven \cref{lem:hom_additivity_over_product} plus the scaling action,
  avoiding the ``image is a closed subgroup'' input of the completeness proof. \qedhere
\end{proof}
\begin{remark}
  \label{rmk:base_case_fourth_route}
  \textbf{The primary route is the \(\mathbb G_m\)-scaling shortcut; why no Theorem~3.2 and no
  \(\mathrm{Hom}(\mathbb G_a, A) = 0\) are needed.} The genus-\(0\) base case ``\(f \colon \mathbb P^1
  \to A\) is constant'' is closed \emph{directly} by the proven additive decomposition
  \cref{lem:hom_additivity_over_product} composed with a total \(\mathbb G\)-action on \(\mathbb P^1\),
  with \emph{no} appeal to Milne's Theorem~3.2 / Lemma~3.3 (rational-map extension), the theorem of
  the cube, Auslander--Buchsbaum, or differentials. The reason the earlier draft thought
  Theorem~3.2 was needed --- that the additive defect \(\psi(x,y) = f(x+y) - f(x) - f(y)\) lives only
  on the affine \(\mathbb G_a \times \mathbb G_a\) and must be extended over \(\mathbb P^1 \times
  \mathbb P^1\) --- evaporates once one composes with a \emph{total} action morphism of
  \cref{def:gaTranslationP1}: the composite is already a total morphism out of \(\mathbb P^1 \times
  \mathbb G\), and the complete factor is \(\mathbb P^1\), exactly what
  \cref{lem:hom_additivity_over_product} (whose Lean signature needs only the first factor
  complete) consumes.
  \medskip
  \noindent\textbf{Primary route --- the scaling shortcut (strategy decision of iter-164).} The
  \emph{multiplicative} scaling action \(\sigma_\times \colon \mathbb P^1 \times \mathbb G_m \to
  \mathbb P^1\), \((x, \lambda) \mapsto \lambda x\) (\cref{def:gaTranslationP1}), closes the base case
  in one shot, \emph{without} the additive homomorphism step or \cref{lem:hom_Ga_to_av_trivial} at
  all. Normalise \(f(0) = 0_A\) (translate; here \(0\) is a scaling fixed point). Form \(h_\times :=
  \sigma_\times \fatsemi f\) and apply \cref{lem:hom_additivity_over_product} with \(V = \mathbb P^1\),
  \(W = \mathbb G_m\), base points \(0 \in \mathbb P^1\) and \(1 \in \mathbb G_m\) (so \(h_\times(0, 1) =
  f(0) = 0_A\)). Because \(0\) is a \emph{fixed point} of scaling, the \(W\)-axis restriction collapses:
  \(\lambda \mapsto h_\times(0, \lambda) = f(\lambda \cdot 0) = f(0) = 0_A\). The decomposition
  therefore reads \(h_\times = p \fatsemi f\) with the second summand the group identity, i.e.
  \[ \sigma_\times \fatsemi f = \mathrm{pr}_1 \fatsemi f \colon \mathbb P^1 \times \mathbb G_m \to A,
     \qquad\text{``}f(\lambda x) = f(x)\text{''}. \]
  Precomposing with \(\lambda \mapsto (1, \lambda)\) gives \(f|_{\mathbb G_m} = \text{const}\,f(1)\);
  since \(\mathbb G_m\) is dense in \(\mathbb P^1\) and \(A\) is separated, \(f\) is constant on \(\mathbb
  P^1\) (the dominant-source/separated-target handle \cref{thm:GrpObj_eq_of_eqOnOpen}, i.e.\ the
  \(\mathtt{ext\_of\_isDominant}\) family already used by \texttt{rigidity\_core}). This uses
  \emph{only} the proven \cref{lem:hom_additivity_over_product}, the total morphism \(\sigma_\times\),
  and density --- no \(\mathbb G_a\) additivity, no \cref{lem:hom_Ga_to_av_trivial}, and in
  particular \emph{no} ``image of a group homomorphism is a closed subgroup'' input (a deep Mathlib
  gap). It is the least-Mathlib-blocked realisation of the base case and is the route the live
  consumer \cref{prop:morphism_P1_to_AV_constant} formalises.
  \medskip
  \noindent\textbf{Retained alternative --- the \(\mathbb G_a\)-additive route.} The Milne-faithful
  argument routes through the translation action \(\sigma \colon \mathbb P^1 \times \mathbb G_a \to
  \mathbb P^1\): \(\sigma \fatsemi f\) is total, the additive decomposition makes \(f|_{\mathbb G_a}\) a
  
  homomorphism \(\mathbb G_a \to A\) (\cref{lem:hom_from_Ga_trivial}), and a homomorphism from the
  
  additive group into an abelian variety is trivial (\cref{lem:hom_Ga_to_av_trivial}, Milne's
  \(\mathbb G_a/\mathbb G_m\) incompatibility, Proposition~3.9). This reaches the same conclusion but
  needs the triviality lemma; it is retained for fidelity to Milne and is \emph{off} the genus-\(0\)
  critical path.
\end{remark}
\begin{lemma}
[A pointed \(\mathbb P^1 \to A\) is constant (Milne Prop 3.9, via the 4th route)]
  \label{lem:hom_from_Ga_trivial}

  \lean{AlgebraicGeometry.morphism_Ga_to_av_const}
  \uses{lem:hom_additivity_over_product, def:gaTranslationP1, lem:hom_Ga_to_av_trivial}
  % SOURCE: [Milne, Abelian Varieties], Proposition 3.9, \S I.3, doc p.~19--20 (read from references/abelian-varieties.pdf, PDF pages 25--26)
  % SOURCE QUOTE: "PROPOSITION 3.9. Every rational map A¹ ⇢ A or P¹ ⇢ A is constant."
  \textit{Source: Milne, Abelian Varieties, Proposition~3.9.}
  \textbf{Off the genus-\(0\) critical path (iter-164): retained alternative route.} This
  \(\mathbb G_a\)-additive close of the base case is \emph{superseded by the \(\mathbb G_m\)-scaling
  shortcut} now installed as the proof of \cref{prop:morphism_P1_to_AV_constant} (see
  \cref{rmk:base_case_fourth_route}). It is retained as the classical Milne Proposition~3.9
  argument (faithful to the \(\mathbb G_a \to A\) homomorphism / \(\mathrm{Hom}(\mathbb G_a, A) = 0\)
  presentation), but the live consumer \cref{prop:morphism_P1_to_AV_constant} no longer depends on
  it; it is provided here only for completeness.
  Let \(A\) be an abelian variety over \(\bar k\) and \(f \colon \mathbb P^1_{\bar k} \to A\) a
  morphism with \(f(0) = 0_A\) (the additive-identity \(\bar k\)-point of \(\mathbb G_a = \mathbb A^1
  \subset \mathbb P^1\); any base point reduces to this by a normalising translation in \(A\)).
  Then \(f\) is constant. Concretely: the translation action \(\sigma\) (\cref{def:gaTranslationP1})
  makes \(f|_{\mathbb G_a}\) an additive-group homomorphism \(\mathbb G_a \to A\), and a homomorphism
  \(\mathbb G_a \to A\) is trivial (\cref{lem:hom_Ga_to_av_trivial}), so \(f\) collapses to the
  constant \(0_A\).
\end{lemma}
% SOURCE QUOTE PROOF: "PROOF. According to (3.2), α extends to a regular map on the whole of A¹.
% After composing α with a translation, we may suppose that α(0)=0. Then α is a homomorphism,
% α(x+y)=α(x)+α(y); all x,y ∈ A¹(k^al)=k^al. But A¹−{0} is also a group variety, and similarly,
% α(xy)=α(x)+α(y)+c; all x,y ∈ A¹(k^al)=k^al. This is absurd, unless α is constant." [Milne,
% Prop 3.9, \S I.3, doc p.~19--20.]
\begin{proof}
  \textit{Source: Milne, Abelian Varieties, Proposition~3.9.}
  Normalise \(f(0) = 0_A\) by a translation in \(A\) (un-translate at the end). The proof is the
  \emph{4th route} (\cref{rmk:base_case_fourth_route}): the additive defect map is never built on
  the affine surface and never extended; instead we precompose \(f\) with the \emph{total}
  translation action \(\sigma\), and the complete factor is \(\mathbb P^1\) itself, so the proven
  additive decomposition applies directly. Milne's Theorem~3.2
  (\cref{thm:rational_map_to_av_extends}) is \emph{not} used.
  \textbf{\(f|_{\mathbb G_a}\) is an additive homomorphism, via the total action \(\sigma\).} Form
  \[ h := \sigma \fatsemi f \colon \mathbb P^1 \times \mathbb G_a \to A, \qquad
     \text{``}h(x, y) = f(x + y)\text{''}, \]
  using the translation action \(\sigma\) of \cref{def:gaTranslationP1}; this is a \emph{total}
  morphism (no rational-map extension needed). Apply the additive decomposition
  \cref{lem:hom_additivity_over_product} (Milne Cor~1.5) --- whose Lean signature requires only the
  \emph{first} factor complete --- with \(V = \mathbb P^1\) (proper), \(W = \mathbb G_a\), and base
  points \(0 \in \mathbb P^1\), \(0 \in \mathbb G_a\), noting \(h(0, 0) = f(0) = 0_A\). The axis
  restrictions are \(x \mapsto h(x, 0) = f(x)\) and \(y \mapsto h(0, y) = f(y)\), so the decomposition
  \(h = (p \fatsemi f) + (q \fatsemi (f|_{\mathbb G_a}))\) reads, on \(\bar k\)-points,
  \[ f(x + y) = f(x) + f(y), \qquad x \in \mathbb P^1(\bar k),\ y \in \mathbb G_a(\bar k) = \bar k. \]
  Restricting \(x\) to \(\mathbb G_a\) shows \(f|_{\mathbb G_a} \colon \mathbb G_a \to A\) is a
  homomorphism of group objects (it sends \(0 \mapsto 0_A\) and respects addition).
  \textbf{A homomorphism \(\mathbb G_a \to A\) is trivial.} By \cref{lem:hom_Ga_to_av_trivial}
  (the \(\mathrm{Hom}(\mathbb G_a, A) = 0\) lemma, the content of Milne's \(\mathbb G_a/\mathbb G_m\)
  incompatibility), \(f|_{\mathbb G_a}\) is the constant morphism at \(0_A\). Hence \(f\) is constant on
  the dense open \(\mathbb G_a \subset \mathbb P^1\); since \(\mathbb G_a\) is dense and \(A\) is
  separated, \(f\) is constant on all of \(\mathbb P^1\), with value \(f(0) = 0_A\) (i.e.\ the original
  base value before the normalising translation). \qedhere
  \medskip
  \noindent\emph{Note: this route is the demoted alternative.} The primary base-case proof is the
  \(\mathbb G_m\)-scaling shortcut now installed directly as the proof of
  \cref{prop:morphism_P1_to_AV_constant} (\cref{rmk:base_case_fourth_route}): the multiplicative
  scaling action \(\sigma_\times\) of \cref{def:gaTranslationP1} gives \(\sigma_\times \fatsemi f =
  \mathrm{pr}_1 \fatsemi f\) (the \(W\)-axis collapses because \(0\) is a scaling fixed point), so
  
  \(f|_{\mathbb G_m}\) is constant and \(f\) is constant by density --- bypassing both the \(\mathbb
  G_a\)-additivity step and \cref{lem:hom_Ga_to_av_trivial}, and avoiding the ``image is a closed
  subgroup'' Mathlib gap. The present \(\mathbb G_a\)-additive lemma is retained only for fidelity to
  \uses{lem:hom_additivity_over_product, def:gaTranslationP1, lem:gmScaling_fixes_zero}
  Milne; it is not on the genus-\(0\) critical path.
\end{proof}
\section{Morphisms from \(\mathbb P^1\) to an abelian variety are constant}
\begin{proposition}
\leanok
[A morphism \(\mathbb P^1 \to A\) is constant]
  \label{prop:morphism_P1_to_AV_constant}

  \lean{AlgebraicGeometry.morphism_P1_to_grpScheme_const}
  \uses{lem:hom_additivity_over_product, def:gaTranslationP1, lem:gmScaling_fixes_zero, thm:GrpObj_eq_of_eqOnOpen, lem:morphism_P1_to_grpScheme_const_aux}
  % SOURCE: [Milne, Abelian Varieties], Proposition 3.10 (specialised to a single \(\mathbb P^1\) factor), \S I.3, doc p.~20 (read from references/abelian-varieties.pdf, PDF page 26)
  % SOURCE QUOTE: "PROPOSITION 3.10. Every rational map \(\alpha \colon V \dashrightarrow A\)
  % from a unirational variety to an abelian variety is constant."
  % SOURCE QUOTE (unirational, same page): "A variety \(V\) over an algebraically closed field
  % is said to be unirational if there is a dominating rational map $\mathbb A^n
  % \dashrightarrow V\( with \)n = \dim V\(; equivalently, if \)k(V)$ can be embedded into
  % \(k(X_1, \ldots, X_n)\) (pure transcendental extension of \(k\)). A variety \(V\) over an
  % arbitrary field \(k\) is said to be unirational if \(V_{k^{al}}\) is unirational."
  \textit{Source: Milne, Abelian Varieties, Proposition~3.10 (case \(V = \mathbb P^1\)).}
  Let \(\bar k\) be algebraically closed and \(A\) an abelian variety over \(\bar k\). Then every
  morphism \(f \colon \mathbb P^1_{\bar k} \to A\) is constant: \(f(\mathbb P^1_{\bar k})\) is a
  single \(\bar k\)-point.
\end{proposition}
% SOURCE QUOTE PROOF: "PROOF. We may suppose that \(k\) is algebraically closed. By assumption
% there is a rational map \(\mathbb A^n \dashrightarrow V\) with dense image, and the composite
% of this with \(\alpha\) extends to a morphism $\beta \colon \mathbb P^1 \times \cdots \times
% \mathbb P^1 \to A$. An induction argument, starting from Corollary 1.5, shows that there are
% regular maps \(\beta_i \colon \mathbb P^1 \to A\) such that $\beta(x_1, \ldots, x_d) = \sum
% \beta_i(x_i)\(, and the lemma shows that each \)\beta_i$ is constant."
\begin{proof}
  This is the single-factor (\(d = 1\)) instance of Milne's Proposition~3.10, where the
  unirational variety is \(\mathbb P^1\) itself. We give the project's \textbf{primary} proof, the
  \emph{\(\mathbb G_m\)-scaling shortcut}: it proves constancy \emph{in one shot} from the proven
  additive decomposition (\cref{lem:hom_additivity_over_product}) and a scaling fixed point,
  without the intermediate ``\(f|_{\mathbb G_a}\) additive'' step and \emph{without} the
  \(\mathrm{Hom}(\mathbb G_a, A) = 0\) triviality lemma. The theorem of the cube is \emph{not} used
  (\cref{rmk:cube_not_needed}), and the map \(f\) is already a morphism out of the complete
  nonsingular \(\mathbb P^1\), so Milne's Theorem~3.2 (\cref{thm:rational_map_to_av_extends}) is
  \emph{not} invoked either.
  \textbf{Reduction to a base point.} Choose the scaling fixed point \(0 \in \mathbb G_a \subset
  \mathbb P^1\) (one of the two fixed points of \(\sigma_\times\)). Translating in the group \(A\) by
  \(-f(0)\), we may assume \(f(0) = 0_A\); un-translating at the end restores the original value.
  \textbf{The scaling shortcut.} Form the \emph{total} morphism
  \[ h := \sigma_\times \fatsemi f \colon \mathbb P^1 \times \mathbb G_m \to A, \qquad
     \text{``}h(x, \lambda) = f(\lambda x)\text{''}, \]
  using the multiplicative scaling action \(\sigma_\times\) of \cref{def:gaTranslationP1} (a genuine
  total morphism --- no rational-map extension needed). Apply the additive decomposition
  \cref{lem:hom_additivity_over_product} (Milne Cor~1.5) --- whose Lean signature requires only the
  \emph{first} factor complete --- with \(V = \mathbb P^1\) (proper), \(W = \mathbb G_m\), and base
  points \(0 \in \mathbb P^1\), \(1 \in \mathbb G_m\); here \(h(0, 1) = f(1 \cdot 0) = f(0) = 0_A\), as
  required. The decomposition writes \(h = (p \fatsemi f_0) + (q \fatsemi g)\) with \(f_0\) the
  \(V\)-axis restriction \(x \mapsto h(x, 1) = f(1 \cdot x) = f(x)\) (so \(f_0 = f\)) and \(g\) the
  \(W\)-axis restriction \(\lambda \mapsto h(0, \lambda) = f(\lambda \cdot 0) = f(0) = 0_A\).
  \textbf{The \(W\)-axis collapses (the load-bearing scaling fixed point).} Because \(0\) is a
  \emph{fixed point} of scaling, \(\sigma_\times(0, \lambda) = 0\) for all \(\lambda\)
  (\cref{lem:gmScaling_fixes_zero}), so the \(W\)-axis
  restriction \(g\) is the \emph{constant} morphism at \(0_A\) --- the identity of the hom-group
  \(\mathrm{Hom}(\mathbb G_m, A)\). The corresponding summand \(q \fatsemi g\) is therefore the group
  identity, and the decomposition collapses to
  \[ h = p \fatsemi f, \qquad\text{i.e.}\qquad
     \sigma_\times \fatsemi f = \mathrm{pr}_1 \fatsemi f, \qquad
     \text{``}f(\lambda x) = f(x)\text{''} \quad (x \in \mathbb P^1,\ \lambda \in \mathbb G_m). \]
  \textbf{Constancy via density.} Specialise \(x = 1\) (precompose with \(\lambda \mapsto (1,
  \lambda)\)): \(f(\lambda) = f(1)\) for all \(\lambda \in \mathbb G_m\), i.e.\ \(f|_{\mathbb G_m}\) is
  the constant morphism at \(f(1)\). Since \(\mathbb G_m\) is dense in the irreducible \(\mathbb P^1\)
  and \(A\) is separated, the two morphisms \(f\) and \(\mathrm{const}_{f(1)}\), agreeing on the
  non-empty open \(\mathbb G_m\), agree everywhere by the dominant-source/separated-target rigidity
  handle \cref{thm:GrpObj_eq_of_eqOnOpen} (the same \texttt{ext\_of\_isDominant} family
  \texttt{rigidity\_core} uses). Hence \(f\) is constant, with value \(f(1)\); in the normalised
  picture \(f(0) = 0_A\), so the constant is \(0_A\) (the original \(f(0)\) before translation).
  \medskip
  \noindent\emph{Multi-factor remark.} For the general unirational statement of Proposition~3.10,
  one writes \(\beta(x_1, \ldots, x_d) = \sum_i \beta_i(x_i)\) via the additive decomposition
  (\cref{lem:hom_additivity_over_product}) and applies this base case to each \(\beta_i \colon
  \mathbb P^1 \to A\); we need only the \(d = 1\) case here. The Milne-faithful \(\mathbb
  G_a\)-additive route (\cref{lem:hom_from_Ga_trivial}, via \(\mathrm{Hom}(\mathbb G_a, A) = 0\),
  \cref{lem:hom_Ga_to_av_trivial}) reaches the same conclusion and is retained as an alternative,
  but is \emph{not} on this critical path.
  % NOTE (iter-166 review): The Lean factors the proof into the outer translation reduction
  % (`AlgebraicGeometry.morphism_P1_to_grpScheme_const`, public) plus a private basepoint helper
  % (`AlgebraicGeometry.morphism_P1_to_grpScheme_const_aux`) carrying the body of the
  % \(\mathbb G_m\)-scaling shortcut. The chapter presents the proof at a single granularity; the
  % Lean split is structural-only and preserves the chapter's logical content. The helper carries
  % five honest scaffold sorries (three \(\mathbb P^1 \otimes \mathbb G_m\) product-instance
  % obligations [GeometricallyIrreducible / LocallyOfFiniteType / IsReduced], one
  % \(\mathtt{IsReduced}\ (\mathbb P^1)_{\text{left}}\) obligation, and one
  % \(\mathtt{IsDominant}\ \iota_{\mathbb G_m}.\text{left}\) obligation — the latter strictly
  % downstream of \(\sigma_\times\) landing a body); these are Mathlib bridge work and a Lane-2
  % dependency, not laundered hypotheses. Both `lean-vs-blueprint-checker-avr-iter166` and
  % `lean-auditor-iter166` confirmed: no iter-157-style laundering, basepoint hypothesis \(hf_0\)
  % is genuinely consumed (L975 of the Lean: `rw [← Category.assoc, hcollapse]; exact hf0`),
  % `lean_verify` axioms $\{\mathtt{propext}, \mathtt{sorryAx}, \mathtt{Classical.choice},
  % \mathtt{Quot.sound}\}$ propagate cleanly. Lifts to axiom-clean once the 5 helper sorries
  % discharge.
  % NOTE (iter-167 review): the 5 helper sorries above have been reduced to ONE. The four
  % product / \(\mathbb P^1\) instance obligations are now discharged via Lane A's named exports
  % (`projGm_geomIrred`, `projGm_locallyOfFiniteType`, `projGm_isReduced`,
  % `projectiveLineBar_isReduced` in `AlgebraicJacobian.Genus0BaseObjects`) and resolve by
  % `infer_instance` inside the helper. The remaining $\mathtt{IsDominant}\ \iota_{\mathbb G_m}
  % .\text{left}$ obligation has been promoted to a named file-local lemma
  % `AlgebraicGeometry.iotaGm_isDominant` (still `sorry`, strictly gated on Lane A shipping
  % the chartwise body of \(\sigma_\times = \mathtt{gmScalingP1}\)). So the helper now contains
  % zero inline `sorry`s and the AVR-side count for this proof is exactly one named bridge
  % (`iotaGm_isDominant`). `lean-vs-blueprint-checker-avr-iter167` confirms no laundering and
  % zero excuse-comments anywhere in AVR.lean.
  % NOTE (iter-168 review): of the 4 Lane A product/\(\mathbb P^1\) exports cited by the
  % iter-167 NOTE, `projectiveLineBar_isReduced` is now genuinely proven axiom-clean
  % (iter-168 substantive closure in `AlgebraicJacobian.Genus0BaseObjects`); the prior
  % "Mathlib gap" framing was incorrect — the bridge
  % `HomogeneousLocalization.Away` \(\to\) `IsDomain` is \(<10\) LOC via
  % `Function.Injective.isDomain` on the algebraMap into `Localization.Away`. The other 3
  % cited Lane A exports (`projGm_geomIrred`, `projGm_isReduced`, and the AVR-side
  % NOTE (iter-185 review): per `lean-vs-blueprint-checker-iter185-gmscaling`,
  % `gm_geomIrred` and `projGm_isReduced` are CONFIRMED Mathlib gaps at b80f227:
  % `Algebra.TensorProduct.isDomain_of_isAlgClosed_left` (or any "factor is alg-closed
  % + other is domain ⟹ tensor is domain" bridge) is absent. Both instances stay as
  % typed sorries until upstream lands the bridge, OR a project-side `~30-50` LOC shim
  % via `Function.Injective.isDomain` on the `algebraMap (kbar ⊗ A) (kbar ⊗ A)` chain
  % is contributed. The Lean docstring on each typed-sorry instance documents this gap
  % verbatim; no \(\mathtt{TEMP}\) axiom-by-axiom workaround is being attempted.
  
  % iter-168 also landed two upstream infrastructure pieces in `Genus0BaseObjects.lean` —
  % `projectiveLineBarAffineCover` (the 2-chart affine cover) and the chart-ring iso scaffold
  
  % `homogeneousLocalizationAwayIso` (forward direction axiom-clean; `aux_left` round-trip
  % still `sorry`) — that prepare the chartwise glue for the iter-169 attack on
  % `gmScalingP1`'s body via either the direct `Proj.fromOfGlobalSections` route or the iso
  % route. These new infrastructure pieces are not yet pinned by `\lean{...}` blocks in this
  % chapter; this is a blueprint-writer-iter169 dispatch.
\end{proof}
\begin{lemma}\leanok
[Pointed core of the \(\mathbb G_m\)-scaling shortcut]
  \label{lem:morphism_P1_to_grpScheme_const_aux}
  \lean{AlgebraicGeometry.morphism_P1_to_grpScheme_const_aux}
  \uses{lem:hom_additivity_over_product, lem:gmScaling_fixes_zero, lem:iotaGm_isDominant}
  Let \(\bar k\) be algebraically closed and \(A\) an abelian variety over \(\bar k\). If a morphism
  \(f \colon \mathbb P^1 \to A\) kills the base point \(0 \in \mathbb P^1\) (\(f(0) = 0_A\)), then
  \(f\) is the constant morphism at the identity \(0_A\). This is the basepoint-normalised core of
  \cref{prop:morphism_P1_to_AV_constant}: form \(h := \sigma_\times \fatsemi f\), feed it to the
  additive decomposition \cref{lem:hom_additivity_over_product} with \(V = \mathbb P^1\),
  \(W = \mathbb G_m\) and base points \(0\), \(1\); the \(W\)-axis collapses because \(0\) is a
  scaling fixed point (\cref{lem:gmScaling_fixes_zero}), leaving \(f(\lambda x) = f(x)\).
  Specialising along the dense inclusion \(\iota \colon \mathbb G_m \hookrightarrow \mathbb P^1\)
  (\cref{lem:iotaGm_isDominant}) and using separatedness of \(A\) forces \(f\) constant, and the
  basepoint hypothesis pins the constant value at \(0_A\).
\end{lemma}
\section{The headline: rigidity for a pointed genus-\(0\) curve}
\begin{theorem}
\leanok
[Rigidity: a pointed morphism from a genus-\(0\) curve to an abelian variety is constant]
  \label{prop:rigidity_genus0_curve_to_AV}

  \lean{AlgebraicGeometry.rigidity_genus0_curve_to_grpScheme}
  \uses{prop:morphism_P1_to_AV_constant, thm:genus_zero_curve_iso_p1}
  Let \(\bar k\) be an algebraically closed field (\emph{arbitrary characteristic}; no
  \(\mathtt{CharZero}\) hypothesis). Let \(C\) be a smooth proper geometrically irreducible curve
  over \(\bar k\) with \(\genus(C) = 0\), and let \(A\) be a smooth proper geometrically irreducible
  group scheme over \(\bar k\) (an abelian variety) with identity \(\eta_A\). Then for every
  morphism \(f \colon C \to A\) such that \(p \fatsemi f = \eta_A\) for some \(\bar k\)-point
  \(p \colon \mathbf 1 \to C\), the morphism \(f\) equals the constant morphism at the identity,
  \(f = \mathtt{toUnit}\,C \fatsemi \eta_A\).
  This is the project's committed characteristic-free statement that the genus-\(0\)
  Albanese-witness \(\mathtt{genusZeroWitness}\) consumes. Its signature mirrors
  \(\mathtt{AlgebraicGeometry.rigidity\_over\_kbar}\) of \cref{chap:RigidityKbar}
  \emph{verbatim except} that the \([\mathtt{CharZero}\ \bar k]\) instance is dropped: same
  curve typeclasses (\(\mathtt{SmoothOfRelativeDimension}\ 1\), \(\mathtt{IsProper}\),
  \(\mathtt{GeometricallyIrreducible}\) on \(C\)), same abelian-variety typeclasses
  (\(\mathtt{GrpObj}\), \(\mathtt{IsProper}\), \(\mathtt{Smooth}\), \(\mathtt{GeometricallyIrreducible}\)
  on \(A\)), and the same pointed-morphism conclusion shape \(f = \mathtt{toUnit}\,C \fatsemi
  \eta[A]\).
\end{theorem}
\begin{proof}
  By \cref{thm:genus_zero_curve_iso_p1} the genus-\(0\) curve \(C\) over the algebraically closed
  field \(\bar k\) is isomorphic to \(\mathbb P^1_{\bar k}\) via some isomorphism
  \(\varphi \colon C \xrightarrow{\ \sim\ } \mathbb P^1_{\bar k}\). The composite
  \(f \circ \varphi^{-1} \colon \mathbb P^1_{\bar k} \to A\) is then a morphism from
  \(\mathbb P^1_{\bar k}\) to the abelian variety \(A\), hence constant by
  \cref{prop:morphism_P1_to_AV_constant}: its image is a single \(\bar k\)-point \(a_0 \in A\).
  Therefore \(f\) itself is constant with value \(a_0\). The pointed hypothesis
  \(p \fatsemi f = \eta_A\) evaluates this constant at the \(\bar k\)-point \(p\) and pins it to the
  identity, \(a_0 = \eta_A\). Consequently \(f\) factors through the terminal object as the
  constant morphism at the identity, \(f = \mathtt{toUnit}\,C \fatsemi \eta_A\).
  No differential (\(df = 0\)) input, no \(H^0(C, \Omega_C) = 0\) / Serre-duality decomposition,
  and no Picard representability is used; the argument is valid in arbitrary characteristic.
  Once formalised it lets the \([\mathtt{CharZero}\ \bar k]\) hypothesis carried by
  \(\mathtt{rigidity\_over\_kbar}\) be dropped, and (after the \(\bar k \to k\) descent step hosted
  in \cref{chap:Jacobian}, not here) feeds \(\mathtt{genusZeroWitness}.\mathtt{key}\).
  % NOTE (iter-166 review): proof body landed iter-166 via clean iso-transport
  % (`obtain ⟨φ⟩ := genusZero_curve_iso_P1 _hgenus; set g := φ.inv ≫ f; obtain ⟨a₀, hga₀⟩ :=
  % morphism_P1_to_grpScheme_const g; pin a₀ = η[A] via `toUnit_unique` basepoint transport;
  % back-transport via `φ.hom_inv_id`'). `lean_verify` axioms $\{\mathtt{propext},
  % \mathtt{sorryAx}, \mathtt{Classical.choice}, \mathtt{Quot.sound}\}\( — \)\mathtt{sorryAx}$
  % propagates honestly through (a) `genusZero_curve_iso_P1` (RR-bridge sub-build, deferred to
  % iter-167+) and (b) `morphism_P1_to_grpScheme_const`'s 5 helper sorries. Lifts to axiom-clean
  % once both upstreams discharge; checker confirmed no laundering.
\end{proof}

\section{Genus-\(0\) base-object helper declarations (1-to-1 Lean coverage)}
\label{sec:genus0_helpers}
This section records, for 1-to-1 Lean\(\leftrightarrow\)blueprint coverage, the supporting
declarations of the four \texttt{Genus0BaseObjects} sub-files (\texttt{BareScheme},
\texttt{ChartIso}, \texttt{Points}, \texttt{GmScaling}). These are the construction
helpers, rewrite lemmas, and typeclass instances backing the headline objects \(\mathbb
P^1_{\bar k}\), \(\mathbb G_a\), \(\mathbb G_m\), the chart-ring iso, and the \(\mathbb
G_m\)-scaling action already pinned above. The substantive ones carry a short proof
sketch; the pure plumbing (instances, projection/round-trip components, \(\mathtt{eval}\)
rewrites) carries a one-line justification, its purpose being to carry the dependency edge.
All are Archon-original or directly Mathlib-backed; none quotes external source text.

\subsection*{\texttt{BareScheme.lean}: the bare scheme \(\mathbb P^1\) and its presentation supplement}

\begin{definition}[Standard \(\mathbb N\)-grading on \(\bar k[X_0, X_1]\)]
  \label{def:plb_grading}
  \lean{AlgebraicGeometry.projectiveLineBarGrading}
  The standard \(\mathbb N\)-grading on \(\bar k[X_0, X_1]\) by total degree: the
  homogeneous-component decomposition \(\mathtt{MvPolynomial.homogeneousSubmodule}\,(\mathrm{Fin}\,2)\,\bar k\),
  whose degree-\(n\) piece is the \(\bar k\)-submodule of degree-\(n\) homogeneous
  polynomials. This is the graded ring whose \(\Proj\) is \(\mathbb P^1_{\bar k}\).
\end{definition}

\begin{lemma}[Graded-ring structure on the standard grading]
  \label{lem:plb_grading_gradedRing}
  \lean{AlgebraicGeometry.projectiveLineBarGrading_gradedRing}
  \uses{def:plb_grading}
  The standard \(\mathbb N\)-grading of \cref{def:plb_grading} is a \(\mathtt{GradedRing}\).
\end{lemma}
\begin{proof}
  \uses{def:plb_grading}
  This is Mathlib's \(\mathtt{MvPolynomial.gradedAlgebra}\); proved directly in Lean.
\end{proof}

\begin{definition}[\(\bar k\)-algebra structure on the homogeneous away-localization]
  \label{def:algebra_kbar_away}
  \lean{AlgebraicGeometry.algebraKbarAway}
  \uses{def:plb_grading}
  For a polynomial \(f \in \bar k[X_0, X_1]\), the \(\bar k\)-algebra structure on the
  homogeneous localization \(\mathrm{Away}\,\mathcal A\,f\) obtained by composing the
  algebra map \(\bar k \to \mathcal A_0\) into the degree-\(0\) piece with the canonical map
  \(\mathcal A_0 \to \mathrm{Away}\,\mathcal A\,f\).
\end{definition}
\begin{proof}
  \uses{def:plb_grading}
  Mathlib ships only \(\mathtt{Algebra}\,\mathcal A_0\,(\mathrm{Away}\,\mathcal A\,f)\);
  this instance is the composite \(\bar k \to \mathcal A_0 \to \mathrm{Away}\,\mathcal A\,f\)
  via \(\mathtt{Algebra.compHom}\). Direct construction; proved directly in Lean.
\end{proof}

\begin{definition}[\(\mathbb P^1_{\bar k}\) as a bare scheme]
  \label{def:plb_scheme}
  \lean{AlgebraicGeometry.ProjectiveLineBarScheme}
  \uses{def:plb_grading}
  The projective line over \(\Spec \bar k\) as a bare scheme: \(\Proj \mathcal A\) of the
  standard \(\mathbb N\)-graded \(\bar k[X_0, X_1]\) (\cref{def:plb_grading}). The object
  \(\mathbb P^1_{\bar k}\) of \(\mathrm{Over}\,(\Spec \bar k)\) (\cref{def:genus0_base_objects})
  is this scheme equipped with its structure morphism.
\end{definition}

\begin{lemma}[The \(\mathrm{Over}\,(\Spec \bar k)\) structure on \(\mathbb P^1\)]
  \label{lem:plb_scheme_canOver}
  \lean{AlgebraicGeometry.projectiveLineBarScheme_canOver}
  \uses{def:plb_scheme, def:plb_grading}
  The natural \(\mathrm{Over}\,(\Spec \bar k)\) structure on \(\mathbb P^1_{\bar k}\), with
  structure morphism \(\Proj.\mathrm{toSpecZero}\,\mathcal A\) followed by the \(\Spec\) of
  the algebra map \(\bar k \to \mathcal A_0\).
\end{lemma}
\begin{proof}
  \uses{def:plb_scheme, def:plb_grading}
  Direct construction of the structure morphism; proved directly in Lean.
\end{proof}

\begin{definition}[Canonical generators of \(\mathrm{MvPolynomial}\,\alpha\,R\)]
  \label{def:mvpoly_generators}
  \lean{AlgebraicGeometry.mvPolyGenerators}
  The canonical \(\mathtt{Algebra.Generators}\) of \(R[\alpha] = \mathrm{MvPolynomial}\,\alpha\,R\)
  over \(R\), indexed by \(\alpha\), with generating family \(\alpha \ni s \mapsto X_s\) and
  section the identity (\(\mathtt{aeval}\,X\) is the identity on \(R[\alpha]\)).
\end{definition}
\begin{proof}
  This is \(\mathtt{Algebra.Generators.ofSurjective}\) of \(X\) with the section
  \(s \mapsto X_s\); proved directly in Lean.
\end{proof}

\begin{definition}[Relation-free presentation of \(\mathrm{MvPolynomial}\,\alpha\,R\)]
  \label{def:mvpoly_presentation}
  \lean{AlgebraicGeometry.mvPolyPresentation}
  \uses{def:mvpoly_generators}
  The canonical \(\mathtt{Algebra.Presentation}\) of \(R[\alpha]\) over \(R\) extending
  \cref{def:mvpoly_generators} with an empty relation index, the kernel of \(\mathtt{aeval}\,X\)
  being \(\bot\) since \(\mathtt{aeval}\,X\) is the identity.
\end{definition}
\begin{proof}
  \uses{def:mvpoly_generators}
  The relation index is \(\mathtt{PEmpty}\); \(\mathrm{span}\,\varnothing = \bot =
  \ker(\mathtt{aeval}\,X)\). Proved directly in Lean.
\end{proof}

\begin{definition}[Pre-submersive presentation of \(\mathrm{MvPolynomial}\,\alpha\,R\)]
  \label{def:mvpoly_presubmersive}
  \lean{AlgebraicGeometry.mvPolyPreSubmersivePresentation}
  \uses{def:mvpoly_presentation}
  The canonical \(\mathtt{Algebra.PreSubmersivePresentation}\) of \(R[\alpha]\) over \(R\)
  extending \cref{def:mvpoly_presentation}; the differential map is on the empty (zero)
  relation module.
\end{definition}
\begin{proof}
  \uses{def:mvpoly_presentation}
  The relation-to-variable map is \(\mathtt{PEmpty.elim}\), trivially injective. Proved
  directly in Lean.
\end{proof}

\begin{definition}[Submersive presentation of \(\mathrm{MvPolynomial}\,\alpha\,R\)]
  \label{def:mvpoly_submersive}
  \lean{AlgebraicGeometry.mvPolySubmersivePresentation}
  \uses{def:mvpoly_presubmersive}
  The canonical \(\mathtt{Algebra.SubmersivePresentation}\) of \(R[\alpha]\) over \(R\)
  extending \cref{def:mvpoly_presubmersive}; the Jacobian is the determinant of the
  \(0 \times 0\) matrix, namely \(1\), a unit. This is the witness from which
  \cref{lem:mvPolynomialFin_isStandardSmoothOfRelativeDimension} reads off standard
  smoothness of relative dimension \(n\).
\end{definition}
\begin{proof}
  \uses{def:mvpoly_presubmersive}
  The Jacobian is \(\det\) of the empty matrix \(= 1\), a unit
  (\(\mathtt{Matrix.det\_isEmpty}\)). Proved directly in Lean.
\end{proof}

\begin{lemma}[Per-chart degree-\(1\) homogeneity witness for the affine cover]
  \label{lem:plb_cover_fDeg}
  \lean{AlgebraicGeometry.projectiveLineBarAffineCover_fDeg}
  \uses{def:plb_grading}
  For each \(i \in \mathrm{Fin}\,2\), the cover element \((\![X_0, X_1]\!)_i\) lies in the
  degree-\((\![1,1]\!)_i\) piece of the standard grading; i.e. \(X_0, X_1\) are homogeneous of
  degree \(1\). This is the homogeneity input to the \(2\)-chart affine cover
  \cref{def:projlinebar_affine_cover}.
\end{lemma}
\begin{proof}
  \uses{def:plb_grading}
  Case analysis on \(i\) and \(\mathtt{MvPolynomial.isHomogeneous\_X}\). Proved directly in Lean.
\end{proof}

\begin{lemma}[Positive-degree witness for the affine cover indexing]
  \label{lem:plb_cover_hm}
  \lean{AlgebraicGeometry.projectiveLineBarAffineCover_hm}
  \uses{lem:plb_cover_fDeg}
  Each entry of the degree vector \((\![1,1]\!) : \mathrm{Fin}\,2 \to \mathbb N\) is positive,
  the positivity hypothesis for the basic-open charts \(D_+(X_0), D_+(X_1)\) of
  \cref{def:projlinebar_affine_cover}.
\end{lemma}
\begin{proof}
  \uses{lem:plb_cover_fDeg}
  Case analysis on the index and \(0 < 1\). Proved directly in Lean.
\end{proof}

\subsection*{\texttt{ChartIso.lean}: the chart-ring iso \(\mathrm{Away}\,\mathcal A\,(X_i) \cong \bar k[u]\)}

\begin{definition}\leanok
[The ``other'' \(\mathrm{Fin}\,2\) index]
  \label{def:otherFin}
  \lean{AlgebraicGeometry.otherFin}
  The involution \(\mathrm{Fin}\,2 \to \mathrm{Fin}\,2\) swapping \(0\) and \(1\); it names
  the complementary coordinate \(X_{\mathrm{other}\,i}/X_i\) in the chart-\(i\) affine
  coordinate.
\end{definition}

\begin{lemma}\leanok
[\(\mathrm{other}\,0 = 1\)]
  \label{lem:otherFin_zero}
  \lean{AlgebraicGeometry.otherFin_zero}
  \uses{def:otherFin}
  \(\mathrm{other}\,0 = 1\).
\end{lemma}
\begin{proof}\uses{def:otherFin} Definitional. Proved directly in Lean.\leanok
\end{proof}

\begin{lemma}\leanok
[\(\mathrm{other}\,1 = 0\)]
  \label{lem:otherFin_one}
  \lean{AlgebraicGeometry.otherFin_one}
  \uses{def:otherFin}
  \(\mathrm{other}\,1 = 0\).
\end{lemma}
\begin{proof}\uses{def:otherFin} Definitional. Proved directly in Lean.\leanok
\end{proof}

\begin{lemma}\leanok
[\(\mathrm{other}\,i \neq i\)]
  \label{lem:otherFin_ne}
  \lean{AlgebraicGeometry.otherFin_ne}
  \uses{def:otherFin}
  For each \(i \in \mathrm{Fin}\,2\), \(\mathrm{other}\,i \neq i\).
\end{lemma}
\begin{proof}\uses{def:otherFin} Decidable case analysis on \(i\). Proved directly in Lean.\leanok
\end{proof}

\begin{definition}\leanok
[The chart-\(i\) evaluation ring map]
  \label{def:chartEvalRingHom}
  \lean{AlgebraicGeometry.chartEvalRingHom}
  \uses{def:otherFin}
  The ring map \(\bar k[X_0, X_1] \to \bar k[u]\) (with \(\bar k[u] = \mathrm{MvPolynomial}\,\mathrm{Unit}\,\bar k\))
  sending \(X_i \mapsto 1\) and \(X_{\mathrm{other}\,i} \mapsto u\) (and constants to constants).
  This dehomogenisation underlies the forward chart-ring map (\cref{def:hlaway_to_mvpoly}).
\end{definition}

\begin{lemma}\leanok
[Chart evaluation sends \(X_i \mapsto 1\)]
  \label{lem:chartEvalRingHom_X_self}
  \lean{AlgebraicGeometry.chartEvalRingHom_X_self}
  \uses{def:chartEvalRingHom}
  The chart-\(i\) evaluation sends \(X_i\) to \(1\).
\end{lemma}
\begin{proof}\uses{def:chartEvalRingHom} Direct computation; proved directly in Lean.\leanok
\end{proof}

\begin{lemma}\leanok
[Chart evaluation sends \(X_{\mathrm{other}\,i} \mapsto u\)]
  \label{lem:chartEvalRingHom_X_other}
  \lean{AlgebraicGeometry.chartEvalRingHom_X_other}
  \uses{def:chartEvalRingHom, lem:otherFin_ne}
  The chart-\(i\) evaluation sends \(X_{\mathrm{other}\,i}\) to the generator \(u\).
\end{lemma}
\begin{proof}
  \uses{def:chartEvalRingHom, lem:otherFin_ne}
  \leanok
  By \(\mathtt{eval}_2\) on the variable, using \(\mathrm{other}\,i \neq i\)
  (\cref{lem:otherFin_ne}). Proved directly in Lean.
\end{proof}

\begin{lemma}\leanok
[Chart evaluation fixes constants]
  \label{lem:chartEvalRingHom_C}
  \lean{AlgebraicGeometry.chartEvalRingHom_C}
  \uses{def:chartEvalRingHom}
  The chart-\(i\) evaluation sends a constant \(C\,r\) to \(C\,r\).
\end{lemma}
\begin{proof}\uses{def:chartEvalRingHom} Direct computation; proved directly in Lean.\leanok
\end{proof}

\begin{definition}\leanok
[Forward chart-ring map \(\mathrm{Away}\,\mathcal A\,(X_i) \to \bar k[u]\)]
  \label{def:hlaway_to_mvpoly}
  \lean{AlgebraicGeometry.homogeneousLocalizationAwayToMvPoly}
  \uses{def:chartEvalRingHom, lem:chartEvalRingHom_X_self, def:plb_grading}
  The forward direction of the chart-ring iso (\cref{def:proj_chart_ring_iso}): the ring map
  \(\mathrm{Away}\,\mathcal A\,(X_i) \to \bar k[u]\) obtained from the chart evaluation
  \cref{def:chartEvalRingHom} by \(\mathtt{Localization.awayLift}\) (legitimate because the
  evaluation sends \(X_i\) to the unit \(1\), \cref{lem:chartEvalRingHom_X_self}), precomposed
  with the canonical map from the homogeneous localization to the ordinary localization.
\end{definition}

\begin{definition}\leanok
[Base ring map \(\bar k \to \mathrm{Away}\,\mathcal A\,(X_i)\)]
  \label{def:kbarToAwayRingHom}
  \lean{AlgebraicGeometry.kbarToAwayRingHom}
  \uses{def:plb_grading, def:algebra_kbar_away}
  The composite \(\bar k \to \mathcal A_0 \to \mathrm{Away}\,\mathcal A\,(X_i)\) of the algebra
  map into the degree-\(0\) piece with \(\mathtt{fromZeroRingHom}\); it carries the scalars of
  the inverse chart-ring map (\cref{def:mvpoly_to_hlaway}).
\end{definition}
\begin{proof}
  \uses{def:plb_grading, def:algebra_kbar_away}
  \leanok
  Composition of the two canonical ring maps; proved directly in Lean.
\end{proof}

\begin{definition}\leanok
[Inverse chart-ring map \(\bar k[u] \to \mathrm{Away}\,\mathcal A\,(X_i)\)]
  \label{def:mvpoly_to_hlaway}
  \lean{AlgebraicGeometry.mvPolyToHomogeneousLocalizationAway}
  \uses{def:kbarToAwayRingHom, def:otherFin, def:plb_grading}
  The inverse direction of the chart-ring iso: the ring map \(\bar k[u] \to \mathrm{Away}\,\mathcal A\,(X_i)\)
  given by the universal property of \(\bar k[u]\) over the base map \cref{def:kbarToAwayRingHom},
  sending the generator \(u\) to the degree-\(0\) element \(X_{\mathrm{other}\,i}/X_i\) (the
  canonical localization element \(\mathtt{isLocalizationElem}\)).
\end{definition}

\begin{lemma}\leanok
[Right round-trip of the chart-ring iso]
  \label{lem:hlaway_iso_aux_right}
  \lean{AlgebraicGeometry.homogeneousLocalizationAwayIso_aux_right}
  \uses{def:hlaway_to_mvpoly, def:mvpoly_to_hlaway}
  The composite \(\text{forward} \circ \text{inverse}\) is the identity on \(\bar k[u]\); i.e.
  \cref{def:hlaway_to_mvpoly} is a left inverse of \cref{def:mvpoly_to_hlaway}.
\end{lemma}
\begin{proof}
  \uses{def:hlaway_to_mvpoly, def:mvpoly_to_hlaway}
  \leanok
  By \(\mathtt{ringHom\_ext}\) it suffices to check on constants and on the generator \(u\); on
  constants both sides give \(C\,r\), and on \(u\) the away-lift computation returns \(u\)
  (using \(\mathtt{Localization.awayLift\_mk}\) and the chart-evaluation rewrites). Proved
  directly in Lean. Together with the left round-trip
  (\cref{lem:proj_chart_ring_iso_aux_left}) this assembles the iso
  \cref{def:proj_chart_ring_iso}.
\end{proof}

\begin{lemma}\leanok
[Per-chart smoothness of \(\Spec(\mathrm{Away}\,\mathcal A\,(X_i)) \to \Spec \bar k\)]
  \label{lem:plb_smooth_chart_X}
  \lean{AlgebraicGeometry.projectiveLineBar_smooth_chart_X}
  \uses{def:proj_chart_ring_iso, lem:chart_ring_iso_preserves_algebraMap, def:plb_grading, lem:plb_scheme_canOver}
  For each \(i\), the chart morphism \(\mathtt{awayι}\,(X_i) \fatsemi \mathbb P^1.\mathrm{hom}
  \colon \Spec(\mathrm{Away}\,\mathcal A\,(X_i)) \to \Spec \bar k\) is smooth of relative
  dimension \(1\).
\end{lemma}
\begin{proof}
  \uses{def:proj_chart_ring_iso, lem:chart_ring_iso_preserves_algebraMap, def:plb_grading}
  \leanok
  Rewriting \(\mathtt{awayι} \fatsemi \mathbb P^1.\mathrm{hom}\) through
  \(\mathtt{awayι\_toSpecZero}\) collapses the composite to \(\Spec\) of the base ring map
  \(\bar k \to \mathrm{Away}\,\mathcal A\,(X_i)\). By the affine-local criterion this reduces to
  standard smoothness of relative dimension \(1\) of that algebra map; transporting along the
  chart-ring iso \cref{def:proj_chart_ring_iso} (a \(\bar k\)-algebra isomorphism by
  \cref{lem:chart_ring_iso_preserves_algebraMap}) and the rename \(\mathrm{Fin}\,1 \cong
  \mathrm{Unit}\) identifies \(\mathrm{Away}\,\mathcal A\,(X_i)\) with \(\bar k[\mathrm{Fin}\,1]\),
  which is standard smooth of relative dimension \(1\)
  (\cref{lem:mvPolynomialFin_isStandardSmoothOfRelativeDimension}). Proved directly in Lean.
\end{proof}

\begin{lemma}\leanok
[Per-chart smoothness over the affine cover]
  \label{lem:plb_smooth_chart_aux}
  \lean{AlgebraicGeometry.projectiveLineBar_smooth_chart_aux}
  \uses{lem:plb_smooth_chart_X, def:projlinebar_affine_cover}
  For each \(i\), the cover map \((\mathtt{projectiveLineBarAffineCover})_i \fatsemi
  \mathbb P^1.\mathrm{hom}\) is smooth of relative dimension \(1\). This is the per-chart input
  to \cref{lem:projectiveLineBar_smoothOfRelDim}.
\end{lemma}
\begin{proof}
  \uses{lem:plb_smooth_chart_X, def:projlinebar_affine_cover}
  \leanok
  The cover's \(i\)-th map is \(\mathtt{awayι}\,(\![X_0, X_1]\!)_i\), which reduces
  definitionally to \(\mathtt{awayι}\,(X_i)\); apply \cref{lem:plb_smooth_chart_X}. Proved
  directly in Lean.
\end{proof}

\subsection*{\texttt{Points.lean}: standard \(\bar k\)-points and the groups \(\mathbb G_a\), \(\mathbb G_m\)}

\begin{definition}[Evaluation ring map of a \(\bar k\)-point of \(\mathbb P^1\)]
  \label{def:plb_evalIntoGlobal}
  \lean{AlgebraicGeometry.ProjectiveLineBar.evalIntoGlobal}
  For a vector \(v \colon \mathrm{Fin}\,2 \to \bar k\), the ring map \(\bar k[X_0, X_1] \to
  \Gamma(\Spec \bar k, \top)\) sending \(X_j \mapsto v_j\) followed by the inverse of the
  \(\Gamma\!\Spec\) iso. This is the underlying evaluation of the point \([v_0 : v_1]\).
\end{definition}

\begin{lemma}[Irrelevant ideal maps onto the unit ideal]
  \label{lem:plb_irrelevant_map_eq_top}
  \lean{AlgebraicGeometry.ProjectiveLineBar.irrelevant_map_eq_top}
  \uses{def:plb_evalIntoGlobal, def:plb_grading}
  If some coordinate \(v_i\) is a unit, then the image of the irrelevant ideal \((X_0, X_1)\)
  under \cref{def:plb_evalIntoGlobal} generates the unit ideal of \(\Gamma(\Spec \bar k, \top)\)
  --- the condition required to define a point of \(\Proj\).
\end{lemma}
\begin{proof}
  \uses{def:plb_evalIntoGlobal, def:plb_grading}
  \(X_i\) lies in the irrelevant ideal, and its image \(v_i\) is a unit; hence \(1\) lies in
  the image ideal. Proved directly in Lean.
\end{proof}

\begin{definition}[Constructing a \(\bar k\)-point of \(\mathbb P^1\) from a vector]
  \label{def:plb_pointOfVec}
  \lean{AlgebraicGeometry.ProjectiveLineBar.pointOfVec}
  \uses{def:plb_evalIntoGlobal, lem:plb_irrelevant_map_eq_top, def:plb_grading}
  Given \(v \colon \mathrm{Fin}\,2 \to \bar k\) with a unit coordinate, the \(\bar k\)-point
  \([v_0 : v_1]\) of \(\mathbb P^1_{\bar k}\), encoded as a section of \(\mathbb P^1.\mathrm{hom}\)
  via \(\Proj.\mathtt{fromOfGlobalSections}\) of \cref{def:plb_evalIntoGlobal} (the irrelevant
  condition supplied by \cref{lem:plb_irrelevant_map_eq_top}). The standard points
  \(0, 1, \infty\) (\cref{def:p1bar_zero,def:p1bar_one,def:p1bar_infty}) are instances.
\end{definition}
\begin{proof}
  \uses{def:plb_evalIntoGlobal, lem:plb_irrelevant_map_eq_top, def:plb_grading}
  The underlying map is \(\Proj.\mathtt{fromOfGlobalSections}\); the section condition is
  checked through \(\mathtt{fromOfGlobalSections\_toSpecZero}\) and the collapse of the inner
  composite onto the \(\Gamma\!\Spec\) iso. Proved directly in Lean.
\end{proof}

\begin{definition}[\(\mathbb G_a\) as a bare scheme]
  \label{def:gaScheme}
  \lean{AlgebraicGeometry.GaScheme}
  The underlying scheme of the additive group \(\mathbb G_a\) over \(\Spec \bar k\): the affine
  line \(\mathbb A^1 = \mathtt{AffineSpace}\,(\mathrm{Fin}\,1)\,(\Spec \bar k)\), whose global
  sections are \(\bar k[\mathrm{Fin}\,1]\). The object \(\mathbb G_a\)
  (\cref{def:ga}) is this scheme with its structure morphism.
\end{definition}

\begin{lemma}[\(\mathrm{Over}\,(\Spec \bar k)\) structure on \(\mathbb G_a\)]
  \label{lem:gaScheme_canOver}
  \lean{AlgebraicGeometry.gaScheme_canOver}
  \uses{def:gaScheme}
  The natural \(\mathrm{Over}\,(\Spec \bar k)\) structure on \(\mathbb G_a\), via the affine-space
  structure morphism.
\end{lemma}
\begin{proof}\uses{def:gaScheme} Inherited from \(\mathtt{AffineSpace.over}\); proved directly in Lean.\end{proof}

\begin{lemma}[\(\mathbb G_a\) is an affine morphism]
  \label{lem:ga_isAffineHom}
  \lean{AlgebraicGeometry.ga_isAffineHom}
  \uses{def:ga, def:gaScheme}
  The structure morphism \(\mathbb G_a \to \Spec \bar k\) is affine.
\end{lemma}
\begin{proof}\uses{def:ga, def:gaScheme} Free from the affine-space instance; proved directly in Lean.\end{proof}

\begin{lemma}[\(\mathbb G_a\) is locally of finite presentation]
  \label{lem:ga_lofp}
  \lean{AlgebraicGeometry.ga_locallyOfFinitePresentation}
  \uses{def:ga, def:gaScheme}
  The structure morphism \(\mathbb G_a \to \Spec \bar k\) is locally of finite presentation
  (the index \(\mathrm{Fin}\,1\) is finite).
\end{lemma}
\begin{proof}\uses{def:ga, def:gaScheme} Free from the affine-space instance; proved directly in Lean.\end{proof}

\begin{lemma}[\(\mathbb G_a\) is reduced]
  \label{lem:ga_isReduced}
  \lean{AlgebraicGeometry.ga_isReduced}
  \uses{def:ga, def:gaScheme}
  The underlying scheme of \(\mathbb G_a\) is reduced (its global sections \(\bar k[\mathrm{Fin}\,1]\)
  form a domain).
\end{lemma}
\begin{proof}
  \uses{def:ga, def:gaScheme}
  Transport \(\mathtt{IsReduced}\,(\Spec \bar k[\mathrm{Fin}\,1])\) across the affine-space iso.
  Proved directly in Lean.
\end{proof}

\begin{definition}[The ring \(\bar k[t, t^{-1}]\)]
  \label{def:gmRing}
  \lean{AlgebraicGeometry.GmRing}
  The ring \(\bar k[t, t^{-1}] = \mathtt{Localization.Away}\,t\), localization of \(\bar k[t] =
  \mathrm{MvPolynomial}\,\mathrm{Unit}\,\bar k\) at the powers of \(t = X_{()}\). This is the
  global-sections ring of \(\mathbb G_m\).
\end{definition}

\begin{definition}[\(\mathbb G_m\) as a bare scheme]
  \label{def:gmScheme}
  \lean{AlgebraicGeometry.GmScheme}
  The underlying scheme of the multiplicative group \(\mathbb G_m\) over \(\Spec \bar k\):
  \(\Spec \bar k[t, t^{-1}]\) (\cref{def:gmRing}). This affine encoding (rather than the
  basic-open of \(\mathbb A^1\)) keeps \(\mathbb G_m\) affine. The object \(\mathbb G_m\)
  (\cref{def:gm}) is this scheme with its structure morphism.
\end{definition}

\begin{lemma}[\(\mathrm{Over}\,(\Spec \bar k)\) structure on \(\mathbb G_m\)]
  \label{lem:gmScheme_canOver}
  \lean{AlgebraicGeometry.gmScheme_canOver}
  \uses{def:gmScheme, def:gmRing}
  The natural \(\mathrm{Over}\,(\Spec \bar k)\) structure on \(\mathbb G_m\), with structure
  morphism \(\Spec\) of the algebra map \(\bar k \to \bar k[t, t^{-1}]\).
\end{lemma}
\begin{proof}\uses{def:gmScheme, def:gmRing} Direct construction; proved directly in Lean.\end{proof}

\begin{lemma}[\(\mathbb G_m\) is affine]
  \label{lem:gm_isAffine}
  \lean{AlgebraicGeometry.gm_isAffine}
  \uses{def:gm, def:gmRing}
  The underlying scheme of \(\mathbb G_m\) is affine.
\end{lemma}
\begin{proof}\uses{def:gm, def:gmRing} It is \(\Spec\) of a ring; proved directly in Lean.\end{proof}

\begin{lemma}[\(\mathbb G_m\) is locally of finite presentation]
  \label{lem:gm_lofp}
  \lean{AlgebraicGeometry.gm_locallyOfFinitePresentation}
  \uses{def:gm, def:gmRing}
  The structure morphism \(\mathbb G_m \to \Spec \bar k\) is locally of finite presentation
  (a localization of a polynomial ring at a single element is finitely presented).
\end{lemma}
\begin{proof}
  \uses{def:gm, def:gmRing}
  From \(\mathtt{Algebra.FinitePresentation}\,\bar k\,\bar k[t, t^{-1}]\) via the
  \(\Spec\)-criterion. Proved directly in Lean.
\end{proof}

\begin{lemma}[\(\mathbb G_m\) is reduced]
  \label{lem:gm_isReduced}
  \lean{AlgebraicGeometry.gm_isReduced}
  \uses{def:gm, def:gmRing}
  The underlying scheme of \(\mathbb G_m\) is reduced (\(\bar k[t, t^{-1}]\) is a domain).
\end{lemma}
\begin{proof}\uses{def:gm, def:gmRing} \(\Spec\) of a reduced ring is reduced; proved directly in Lean.\end{proof}

\begin{lemma}[\(\bar k[t, t^{-1}]\) is an integral domain]
  \label{lem:gmRing_isDomain}
  \lean{AlgebraicGeometry.gmRing_isDomain}
  \uses{def:gmRing}
  The ring \(\bar k[t, t^{-1}]\) is an integral domain.
\end{lemma}
\begin{proof}
  \uses{def:gmRing}
  Localization of the domain \(\bar k[t]\) at the powers of the nonzero element \(t\) (a set of
  non-zero-divisors) preserves the domain property. Proved directly in Lean. Load-bearing for
  \cref{lem:gm_irreducibleSpace} and \cref{lem:gm_geomIrred}.
\end{proof}

\begin{lemma}[\(\mathbb G_m\) is irreducible]
  \label{lem:gm_irreducibleSpace}
  \lean{AlgebraicGeometry.gm_irreducibleSpace}
  \uses{def:gm, lem:gmRing_isDomain}
  The underlying topological space of \(\mathbb G_m\) is irreducible.
\end{lemma}
\begin{proof}
  \uses{def:gm, lem:gmRing_isDomain}
  \(\Spec\) of the domain \(\bar k[t, t^{-1}]\) (\cref{lem:gmRing_isDomain}) is irreducible.
  Proved directly in Lean.
\end{proof}

\begin{definition}[Units-of-global-sections functor]
  \label{def:gmHomFunctor}
  \lean{AlgebraicGeometry.gmHomFunctor}
  \uses{def:gm}
  The functor \((\mathrm{Over}\,(\Spec \bar k))^{\mathrm{op}} \to \mathbf{Grp}\) sending
  \(T \mapsto \Gamma(T_{\mathrm{left}}, \top)^\times\), with functorial action by \(\mathtt{Units.map}\)
  of the pullback on global sections. This is the presheaf of groups that \(\mathbb G_m\)
  represents (\cref{def:gmHomFunctor_representableBy}).
\end{definition}
\begin{proof}\uses{def:gm} Functoriality is the identity/composition laws of \(\mathtt{Units.map}\); proved directly in Lean.\end{proof}

\begin{definition}[Forward of the representability bijection]
  \label{def:gmHomEquiv_toFun}
  \lean{AlgebraicGeometry.gmHomEquiv_toFun}
  \uses{def:gm, def:gmRing}
  Given \(f \colon T \to \mathbb G_m\), the unit of \(\Gamma(T_{\mathrm{left}}, \top)^\times\)
  obtained as the image of the standard generator \(t \in \bar k[t, t^{-1}]\) (a unit) under the
  induced ring map \(\bar k[t, t^{-1}] \to \Gamma(T_{\mathrm{left}}, \top)\).
\end{definition}

\begin{lemma}[The inverse morphism is over \(\Spec \bar k\)]
  \label{lem:gmHomEquiv_invFun_isOver}
  \lean{AlgebraicGeometry.gmHomEquiv_invFun_isOver}
  \uses{def:gm, def:gmRing}
  For a unit \(u \in \Gamma(T_{\mathrm{left}}, \top)^\times\), the scheme morphism
  \(T_{\mathrm{left}} \to \mathbb G_m\) built from \(t \mapsto u\) commutes with the structure
  morphisms to \(\Spec \bar k\).
\end{lemma}
\begin{proof}
  \uses{def:gm, def:gmRing}
  The lifted ring map restricted along \(\bar k \to \bar k[t, t^{-1}]\) recovers the structure
  ring map of \(T\); apply \(\mathtt{IsLocalization.Away.lift\_comp}\) and the \(\Gamma\dashv\Spec\)
  unit triangle. Proved directly in Lean.
\end{proof}

\begin{definition}[Backward of the representability bijection]
  \label{def:gmHomEquiv_invFun}
  \lean{AlgebraicGeometry.gmHomEquiv_invFun}
  \uses{lem:gmHomEquiv_invFun_isOver, def:gm}
  Given a unit \(u \in \Gamma(T_{\mathrm{left}}, \top)^\times\), the morphism \(T \to \mathbb G_m\)
  determined by the ring map \(\bar k[t, t^{-1}] \to \Gamma(T_{\mathrm{left}}, \top)\) sending
  \(t \mapsto u\) (lifted through the localization since \(u\) is a unit), bundled with the
  over-\(\Spec \bar k\) commutativity \cref{lem:gmHomEquiv_invFun_isOver}.
\end{definition}

\begin{lemma}[Left round-trip of the representability bijection]
  \label{lem:gmHomEquiv_left_inv}
  \lean{AlgebraicGeometry.gmHomEquiv_left_inv}
  \uses{def:gmHomEquiv_invFun, def:gmHomEquiv_toFun}
  \(\text{invFun} \circ \text{toFun} = \mathrm{id}\) on \(\mathrm{Hom}(T, \mathbb G_m)\).
\end{lemma}
\begin{proof}
  \uses{def:gmHomEquiv_invFun, def:gmHomEquiv_toFun}
  Reduce through the \(\Gamma\dashv\Spec\) injectivity (\(\mathtt{ext\_to\_Spec}\)) and the
  uniqueness of the away-lift; compare the resulting ring maps on constants and on \(t\). Proved
  directly in Lean.
\end{proof}

\begin{lemma}[Right round-trip of the representability bijection]
  \label{lem:gmHomEquiv_right_inv}
  \lean{AlgebraicGeometry.gmHomEquiv_right_inv}
  \uses{def:gmHomEquiv_invFun, def:gmHomEquiv_toFun}
  \(\text{toFun} \circ \text{invFun} = \mathrm{id}\) on \(\Gamma(T_{\mathrm{left}}, \top)^\times\).
\end{lemma}
\begin{proof}
  \uses{def:gmHomEquiv_invFun, def:gmHomEquiv_toFun}
  The image of the generator \(t\) under the lifted map is exactly \(u\)
  (\(\mathtt{IsLocalization.Away.lift\_eq}\)), after cancelling the \(\Gamma\dashv\Spec\) unit
  triangle. Proved directly in Lean.
\end{proof}

\begin{lemma}[Naturality of the representability bijection]
  \label{lem:gmHomEquiv_homEquiv_comp}
  \lean{AlgebraicGeometry.gmHomEquiv_homEquiv_comp}
  \uses{def:gmHomEquiv_toFun, def:gmHomFunctor}
  Precomposition by \(\varphi \colon T' \to T\) corresponds to \(\mathtt{Units.map}\) of
  \(\varphi\) on the units side; i.e. the forward map \cref{def:gmHomEquiv_toFun} is natural in
  \(T\) for the functor \cref{def:gmHomFunctor}.
\end{lemma}
\begin{proof}
  \uses{def:gmHomEquiv_toFun, def:gmHomFunctor}
  Unfold both sides via \(\mathtt{comp\_appTop}\) and transfer the generator's image through the
  away-lift. Proved directly in Lean.
\end{proof}

\begin{definition}[\(\mathbb G_m\) represents the units functor]
  \label{def:gmHomFunctor_representableBy}
  \lean{AlgebraicGeometry.gmHomFunctor_representableBy}
  \uses{def:gmHomFunctor, def:gmHomEquiv_toFun, def:gmHomEquiv_invFun, lem:gmHomEquiv_left_inv, lem:gmHomEquiv_right_inv, lem:gmHomEquiv_homEquiv_comp, def:gm}
  The \(\mathtt{RepresentableBy}\) witness exhibiting \(\mathbb G_m\) as a representing object
  for the units-of-global-sections functor \cref{def:gmHomFunctor}: the natural bijection
  \(\mathrm{Hom}(-, \mathbb G_m) \cong \Gamma(-, \top)^\times\) with components
  \cref{def:gmHomEquiv_toFun}/\cref{def:gmHomEquiv_invFun}, round-trips
  \cref{lem:gmHomEquiv_left_inv}/\cref{lem:gmHomEquiv_right_inv}, and naturality
  \cref{lem:gmHomEquiv_homEquiv_comp}. This is the witness feeding the \(\mathtt{GrpObj}\,\mathbb G_m\)
  structure \cref{def:gm_grpObj}.
\end{definition}
\begin{proof}
  \uses{def:gmHomEquiv_toFun, def:gmHomEquiv_invFun, lem:gmHomEquiv_left_inv, lem:gmHomEquiv_right_inv, lem:gmHomEquiv_homEquiv_comp}
  Assemble the per-\(T\) equivalence from its four components and the naturality square. Proved
  directly in Lean.
\end{proof}

\begin{lemma}[\(\mathbb G_m\) is smooth over \(\Spec \bar k\)]
  \label{lem:gm_smooth}
  \lean{AlgebraicGeometry.gm_smooth}
  \uses{def:gm_grpObj, lem:gm_lofp, lem:gm_isReduced}
  Over an algebraically closed \(\bar k\), the structure morphism \(\mathbb G_m \to \Spec \bar k\)
  is smooth.
\end{lemma}
\begin{proof}
  \uses{def:gm_grpObj, lem:gm_lofp, lem:gm_isReduced}
  A reduced, locally-of-finite-presentation group scheme over an algebraically closed field is
  smooth (\(\mathtt{smooth\_of\_grpObj\_of\_isAlgClosed}\)), using the group structure
  \cref{def:gm_grpObj}. Proved directly in Lean.
\end{proof}

\subsection*{\texttt{GmScaling.lean}: chart bridge, cover isos, and the \(\sigma_\times\) cocycle}

\begin{lemma}\leanok
[Chart bridge: \(\mathtt{awayι} \fatsemi \mathbb P^1.\mathrm{hom} = \Spec\) of the base map]
  \label{lem:awayι_comp_PLB_hom}
  \lean{AlgebraicGeometry.awayι_comp_PLB_hom}
  \uses{def:genus0_base_objects, def:plb_grading, def:algebra_kbar_away}
  For any homogeneous \(f\) of positive degree, \(\mathtt{awayι}\,\mathcal A\,f \fatsemi
  \mathbb P^1.\mathrm{hom} = \Spec(\bar k \to \mathrm{Away}\,\mathcal A\,f)\). Generic in the
  element \(f\) and its degree \(m\) so it applies to both the cover generators \((\![X_0, X_1]\!)_i\)
  and the merged generator \(X_0 X_1\).
\end{lemma}
\begin{proof}
  \uses{def:plb_grading, def:algebra_kbar_away}
  A rewrite chasing \(\mathtt{awayι\_toSpecZero}\), \(\Spec.\mathtt{map\_comp}\), and the
  \cref{def:algebra_kbar_away} defeq. Proved directly in Lean.
\end{proof}

\begin{definition}\leanok
[Chart-1 ring map of \(\sigma_\times\)]
  \label{def:gmscaling_chart1_ringMap}
  \lean{AlgebraicGeometry.gmScalingP1_chart1_ringMap}
  \uses{def:gmRing}
  At the \(\bar k[u]\)-level, the ring map sending the affine coordinate \(u = X_0/X_1\) to
  \(u \otimes \lambda\), where \(\lambda = t \in \bar k[t, t^{-1}]\) --- the local form of the
  scaling \((x, \lambda) \mapsto \lambda x\) near \(0\).
\end{definition}

\begin{definition}\leanok
[Chart-0 ring map of \(\sigma_\times\)]
  \label{def:gmscaling_chart0_ringMap}
  \lean{AlgebraicGeometry.gmScalingP1_chart0_ringMap}
  \uses{def:gmRing}
  At the \(\bar k[u]\)-level, the ring map sending the affine coordinate \(t = X_1/X_0\) to
  \(t \otimes \lambda^{-1}\), where \(\lambda^{-1} = \mathtt{invSelf}\,t \in \bar k[t, t^{-1}]\)
  --- the local form of the scaling near \(\infty\), regular because \(\lambda\) is invertible.
\end{definition}

\begin{definition}\leanok
[Chart-\(i\) source of the scaling cover]
  \label{def:gmscaling_cover_X_iso}
  \lean{AlgebraicGeometry.gmScalingP1_cover_X_iso}
  \uses{lem:awayι_comp_PLB_hom, def:gmscaling_cover, lem:plb_cover_hm, lem:plb_cover_fDeg, def:gmRing}
  The identification of the \(i\)-th chart of the scaling cover \cref{def:gmscaling_cover} with
  \(\Spec\big((\mathrm{Away}\,\mathcal A\,(\![X_0, X_1]\!)_i) \otimes_{\bar k} \bar k[t, t^{-1}]\big)\),
  built by composing \(\mathtt{pullbackSymmetry}\), \(\mathtt{pullbackRightPullbackFstIso}\), the
  \cref{lem:awayι_comp_PLB_hom} rewrite, and \(\mathtt{pullbackSpecIso}\). Used to define the
  chart map \cref{def:gmscaling_chart}.
\end{definition}
\begin{proof}
  \uses{lem:awayι_comp_PLB_hom, lem:plb_cover_hm, lem:plb_cover_fDeg}
  The hoisted degree/positivity witnesses \cref{lem:plb_cover_fDeg}/\cref{lem:plb_cover_hm} keep
  the bridge chain syntactic so \(\mathtt{pullbackSpecIso}\) applies uniformly in \(i\). Proved
  directly in Lean.
\end{proof}

\begin{definition}\leanok
[Cross-chart intersection source of the scaling cover]
  \label{def:gmscaling_cover_intersection_X_iso}
  \lean{AlgebraicGeometry.gmScalingP1_cover_intersection_X_iso}
  \uses{lem:awayι_comp_PLB_hom, def:gmscaling_cover, lem:plb_cover_fDeg, lem:plb_cover_hm, def:gmRing}
  The identification of the pullback of the two charts \((\mathtt{cover}.f\,0)\),
  \((\mathtt{cover}.f\,1)\) of \cref{def:gmscaling_cover} with
  \(\Spec\big((\mathrm{Away}\,\mathcal A\,(X_0 X_1)) \otimes_{\bar k} \bar k[t, t^{-1}]\big)\),
  the cross chart over the merged degree-\(2\) generator \(X_0 X_1\). Built by the
  \(\mathtt{Proj.pullbackAwayιIso}\) route, mirroring \cref{def:gmscaling_cover_X_iso}.
\end{definition}
\begin{proof}
  \uses{lem:awayι_comp_PLB_hom, lem:plb_cover_fDeg, lem:plb_cover_hm}
  A chain of pullback-pasting isos merging the two away-charts into the away-chart at \(X_0 X_1\),
  then the standard \(\mathtt{pullbackSpecIso}\) recipe at degree \(2\). Proved directly in Lean.
\end{proof}

\begin{lemma}\leanok
[Helper: \(\bar k[u] \otimes_{\bar k} K\) is a domain]
  \label{lem:isDomain_mvPolyUnit_tensor}
  \lean{AlgebraicGeometry.isDomain_mvPolyUnit_tensor}
  For any field \(K\) over \(\bar k\), the ring \(\bar k[u] \otimes_{\bar k} K\) is an integral
  domain (the ``geometrically irreducible affine line'' tensor identity).
\end{lemma}
\begin{proof}
  Through the algebra isomorphisms \(\bar k[u] \otimes_{\bar k} K \cong K \otimes_{\bar k} \bar k[u]
  \cong K[u]\) (\(\mathtt{Algebra.TensorProduct.comm}\), \(\mathtt{MvPolynomial.algebraTensorAlgEquiv}\)),
  and \(K[u]\) is a domain. Proved directly in Lean.
\end{proof}

\begin{lemma}\leanok
[\(\mathbb A^1\) is geometrically irreducible over \(\Spec \bar k\)]
  \label{lem:affineLine_geomIrred}
  \lean{AlgebraicGeometry.affineLine_geomIrred}
  \uses{lem:isDomain_mvPolyUnit_tensor}
  The structure morphism \(\Spec \bar k[u] \to \Spec \bar k\) is geometrically irreducible: for
  every field extension \(K/\bar k\) the base change \(\Spec(\bar k[u] \otimes_{\bar k} K) \cong
  \Spec K[u]\) is irreducible.
\end{lemma}
\begin{proof}
  \uses{lem:isDomain_mvPolyUnit_tensor}
  By the closed-under-isomorphism criterion it suffices that each base change is irreducible;
  the base-change ring is a domain by \cref{lem:isDomain_mvPolyUnit_tensor}, so its \(\Spec\) is
  irreducible, transported across \(\mathtt{pullbackSpecIso}\). Proved directly in Lean. This is
  the geometric-irreducibility input to \cref{lem:gm_geomIrred}.
\end{proof}

\begin{lemma}\leanok
[Cross-chart agreement of \(\sigma_\times\) (cocycle)]
  \label{lem:gmscaling_chart_agreement_cross01}
  \lean{AlgebraicGeometry.gmScalingP1_chart_agreement_cross01}
  \uses{def:gmscaling_cover_intersection_X_iso, def:gmscaling_chart, lem:gmscaling_chart_PLB_eq, lem:projlinebar_isReduced}
  The substantive \((0,1)\) cross case of the \(\sigma_\times\) gluing condition: on the overlap
  \(D_+(X_0 X_1) \times \mathbb G_m\), the chart-\(0\) map (\(t \mapsto t \otimes \lambda^{-1}\))
  and the chart-\(1\) map (\(u \mapsto u \otimes \lambda\)) agree, i.e.
  \[ \mathtt{pullback.fst} \fatsemi (\mathtt{chart}\,0) = \mathtt{pullback.snd} \fatsemi (\mathtt{chart}\,1). \]
  This is the cocycle identity feeding \cref{lem:gmscaling_chart_agreement}.
\end{lemma}
\begin{proof}
  \uses{def:gmscaling_cover_intersection_X_iso, lem:gmscaling_chart_PLB_eq, lem:projlinebar_isReduced}
  On the intersection both chart coordinates are units with \(t \cdot u = 1\); substituting
  \(u = 1/t\) into the chart-1 image \(\lambda u\) gives \(\lambda/t\), which matches the
  chart-0 image \(t/\lambda\) read in the complementary coordinate, so the underlying ring-level
  identity reduces to \(\lambda \cdot u \cdot t = \lambda\), true since \(u t = 1\).
  Structurally one first checks agreement after post-composition with \(\mathbb P^1.\mathrm{hom}\)
  (via the per-chart bridge \cref{lem:gmscaling_chart_PLB_eq} and \(\mathtt{pullback.condition}\)),
  then upgrades to agreement of the two morphisms themselves through the closed-immersion
  diagonal of the separated \(\mathbb P^1.\mathrm{hom}\) (using reducedness of \(\mathbb P^1\),
  \cref{lem:projlinebar_isReduced}) and the cross-chart identification
  \cref{def:gmscaling_cover_intersection_X_iso}. The final factorization through the diagonal is
  not yet closed in Lean (the requisite range-containment substrate is absent from Mathlib at the
  current pin), so this lemma carries a residual sorry; the ring-level identity above is the
  remaining content.
\end{proof}

\subsection{The canonical \(\mathbb G_m \hookrightarrow \mathbb P^1\) inclusion: dominance via the chart-\(1\) open immersion}

\begin{lemma}\leanok
[Chart-\(1\) range containment of the point \([1:1]\)]
  \label{lem:iotaGm_r_1_range_subset}
  \lean{AlgebraicGeometry.iotaGm_r_1_range_subset}
  \uses{def:p1bar_one}
  The topological image of the \(\bar k\)-point \([1:1] \in \mathbb P^1\) (\cref{def:p1bar_one})
  lies inside the open image of the chart-\(1\) affine open
  \(\mathrm{Proj.away}\iota\,X_1 \colon \Spec(\mathrm{Away}\,\mathcal A\,X_1) \to \Proj\,\mathcal A\),
  i.e.\ inside the basic open \(D_+(X_1)\). Indeed \([1:1]\) is the global-sections point with
  \(X_0, X_1 \mapsto 1\), so the value of \(X_1\) is the unit \(1\), and the preimage of
  \(D_+(X_1)\) under this point is all of \(\Spec \bar k\).
\end{lemma}

\begin{definition}\leanok
[Chart-\(1\) lift of the point \([1:1]\)]
  \label{def:iotaGm_r_1}
  \lean{AlgebraicGeometry.iotaGm_r_1}
  \uses{lem:iotaGm_r_1_range_subset, def:p1bar_one}
  The factor \(r_1 \colon \Spec \bar k \to \Spec(\mathrm{Away}\,\mathcal A\,X_1)\) of the
  \(\bar k\)-point \([1:1] \in \mathbb P^1\) (\cref{def:p1bar_one}) through the chart-\(1\) open
  immersion \(\mathrm{Proj.away}\iota\,X_1\), produced by the universal property of the open
  immersion together with the range-containment witness \cref{lem:iotaGm_r_1_range_subset}.
\end{definition}

\begin{lemma}\leanok
[Factorisation of \([1:1]\) through chart \(1\)]
  \label{lem:iotaGm_r_1_fac}
  \lean{AlgebraicGeometry.iotaGm_r_1_fac}
  \uses{def:iotaGm_r_1}
  The chart-\(1\) lift \(r_1\) (\cref{def:iotaGm_r_1}) factors the point \([1:1]\) through the
  chart-\(1\) open immersion: \(r_1 \fatsemi \mathrm{Proj.away}\iota\,X_1 = [1:1]\). Immediate from
  the defining factorisation property of the open-immersion lift.
\end{lemma}

\begin{definition}\leanok
[Chart-\(1\) evaluation ring map \(\mathrm{Away}\,\mathcal A\,X_1 \to \bar k\)]
  \label{def:kbarChart1Ring}
  \lean{AlgebraicGeometry.kbarChart1Ring}
  \uses{def:hlaway_to_mvpoly}
  The \(\bar k\)-algebra evaluation map
  \(\mathrm{Away}\,\mathcal A\,X_1 \to \mathrm{MvPolynomial}\,\{\ast\}\,\bar k \to \bar k\), the
  composite of the chart-\(1\) ring iso \(\mathtt{homogeneousLocalizationAwayToMvPoly}\)
  (\cref{def:hlaway_to_mvpoly}, sending the canonical generator \(\mathtt{isLocElem}\) to the
  affine coordinate \(u\)) with evaluation of \(u\) at \(\lambda = 1\). This is the ring-map shadow
  of the point \([1:1]\) read in chart \(1\).
\end{definition}

\begin{lemma}\leanok
[The point \([1:1]\) is the Spec of the chart-\(1\) evaluation map]
  \label{lem:kbarChart1Ring_specMap_fac}
  \lean{AlgebraicGeometry.kbarChart1Ring_specMap_fac}
  \uses{def:kbarChart1Ring, lem:iotaGm_r_1_fac,
    lem:awayi_appIso_top_inv_apply_isLocElem}
  The Spec of the chart-\(1\) evaluation ring map (\cref{def:kbarChart1Ring}), post-composed with
  the chart-\(1\) open immersion \(\mathrm{Proj.away}\iota\,X_1\), recovers the point \([1:1]\):
  \(\Spec(\mathtt{kbarChart1Ring}) \fatsemi \mathrm{Proj.away}\iota\,X_1 = [1:1]\).
\end{lemma}
\begin{proof}
  Reduce through the lift factorisation \cref{lem:iotaGm_r_1_fac} to a focused ring-map identity
  on the affine chart \(1\): both the Spec of the evaluation map and the point \([1:1]\) factor
  as the same open-immersion lift, so it suffices to check the underlying ring maps agree on the
  affine chart's global sections, which they do by construction of \(\mathtt{kbarChart1Ring}\).
\end{proof}

\begin{lemma}\leanok
[The chart-\(1\) lift is the Spec of the evaluation map]
  \label{lem:iotaGm_r_1_eq_specMap}
  \lean{AlgebraicGeometry.iotaGm_r_1_eq_specMap}
  \uses{def:iotaGm_r_1, def:kbarChart1Ring, lem:kbarChart1Ring_specMap_fac}
  The chart-\(1\) lift \(r_1\) (\cref{def:iotaGm_r_1}) coincides with the Spec of the chart-\(1\)
  evaluation ring map (\cref{def:kbarChart1Ring}): \(r_1 = \Spec(\mathtt{kbarChart1Ring})\). Both
  satisfy the same lift factorisation through \(\mathrm{Proj.away}\iota\,X_1\)
  (\cref{lem:iotaGm_r_1_fac}, \cref{lem:kbarChart1Ring_specMap_fac}), so they agree by uniqueness
  of the open-immersion lift.
\end{lemma}

\begin{lemma}\leanok
[Inner-lift compatibility for the chart-\(1\) section]
  \label{lem:iotaGm_inner_lift_compat}
  \lean{AlgebraicGeometry.iotaGm_inner_lift_compat}
  \uses{def:p1bar_one}
  The compatibility identity needed to build the chart-\(1\) section as a pullback lift: the two
  legs \((\iota_{\mathbb G_m} \fatsemi [1:1]) \fatsemi \mathbb P^1.\mathrm{hom}\) and
  \(\mathbf 1 \fatsemi \mathbb G_m.\mathrm{hom}\) agree, which reduces to the structure-map
  compatibility \(\mathtt{Over.w}\) of the point \([1:1]\) (\cref{def:p1bar_one}) over the base.
\end{lemma}

\begin{definition}\leanok
[Chart-\(1\) section of \(\mathbb G_m\) into the scaling cover]
  \label{def:iotaGm_chart1_section}
  \lean{AlgebraicGeometry.iotaGm_chart1_section}
  \uses{lem:iotaGm_inner_lift_compat, def:iotaGm_r_1, lem:iotaGm_r_1_fac}
  The section \(s \colon \mathbb G_m \to (\sigma_\times\text{-cover}).X_1\) into the chart-\(1\)
  member of the scaling cover (\cref{def:gmscaling_cover}), built as a pullback lift with
  \(\mathbb G_m.\mathrm{hom} \fatsemi r_1\) (\cref{def:iotaGm_r_1}) as its chart component and the
  inner lift assembled from \cref{lem:iotaGm_inner_lift_compat}; its defining factorisation uses
  \cref{lem:iotaGm_r_1_fac}.
\end{definition}

\begin{lemma}\leanok
[Chart-\(1\) composition is an open immersion]
  \label{lem:iotaGm_chart1_composition_isOpenImmersion}
  \lean{AlgebraicGeometry.iotaGm_chart1_composition_isOpenImmersion}
  \uses{def:iotaGm_chart1_section, lem:iotaGm_chart1_appIso_eval}
  The composite of the chart-\(1\) section (\cref{def:iotaGm_chart1_section}) with the chart-\(1\)
  scaling map (\cref{def:gmscaling_chart}) is an open immersion \(\mathbb G_m \to \mathbb P^1\). It
  realises the canonical inclusion \(\lambda \mapsto [\lambda : 1]\) as the three-step composite
  \(\Spec \bar k[t, t^{-1}] \hookrightarrow \Spec \bar k[u] \cong \Spec(\mathrm{Away}\,\mathcal A\,X_1)
  \hookrightarrow \mathbb P^1\): the first map is a localisation open immersion, the middle is the
  chart-ring iso, and the last is the chart-\(1\) open immersion, with the composition identity
  supplied by \cref{lem:iotaGm_chart1_appIso_eval}.
\end{lemma}

\begin{lemma}\leanok
[The canonical inclusion \(\mathbb G_m \hookrightarrow \mathbb P^1\) is an open immersion]
  \label{lem:iotaGm_isOpenImmersion}
  \lean{AlgebraicGeometry.iotaGm_isOpenImmersion}
  \uses{def:iotaGm_chart1_section, lem:iotaGm_inner_lift_compat, lem:iotaGm_chart1_composition_isOpenImmersion}
  The canonical morphism \(\iota \colon \mathbb G_m \to \mathbb P^1\) (the ``specialise the scaling
  action at \(x = 1\)'' map) is an open immersion. After unfolding it as a pullback lift glued
  through the scaling cover, it agrees with the chart-\(1\) composite of
  \cref{lem:iotaGm_chart1_composition_isOpenImmersion} (via \cref{def:iotaGm_chart1_section} and
  \cref{lem:iotaGm_inner_lift_compat}), which is an open immersion.
\end{lemma}

\begin{lemma}\leanok
[The image of \(\mathbb G_m \hookrightarrow \mathbb P^1\) is open]
  \label{lem:iotaGm_range_isOpen}
  \lean{AlgebraicGeometry.iotaGm_range_isOpen}
  \uses{lem:iotaGm_isOpenImmersion}
  The topological image of the canonical inclusion \(\iota \colon \mathbb G_m \to \mathbb P^1\) is
  open in \(\mathbb P^1\). Immediate from \(\iota\) being an open immersion
  (\cref{lem:iotaGm_isOpenImmersion}).
\end{lemma}

\begin{lemma}\leanok
[The canonical inclusion \(\mathbb G_m \hookrightarrow \mathbb P^1\) is dominant]
  \label{lem:iotaGm_isDominant}
  \lean{AlgebraicGeometry.iotaGm_isDominant}
  \uses{lem:iotaGm_range_isOpen, lem:projectiveLineBar_geomIrred}
  The canonical inclusion \(\iota \colon \mathbb G_m \to \mathbb P^1\) is dominant: its image is
  dense. Its image is a non-empty open (\cref{lem:iotaGm_range_isOpen}) of the irreducible
  \(\mathbb P^1\) (geometric irreducibility, \cref{lem:projectiveLineBar_geomIrred}, over the
  one-point base \(\Spec \bar k\)), and a non-empty open of an irreducible space is dense.
\end{lemma}
